\documentclass[11pt,a4paper]{extbook}

\usepackage{fontspec}
\usepackage{polyglossia}
\setdefaultlanguage{russian}
\setotherlanguage{english}
\usepackage{geometry}
\usepackage{graphicx}
\usepackage{caption}
\usepackage{hyperref}
\usepackage{tocloft}
\usepackage{fancyhdr}
\usepackage{xcolor}
\usepackage{listings}
\usepackage{array}
\usepackage{longtable}
\usepackage{booktabs}
\usepackage{microtype}
\usepackage{amsmath}
\usepackage{amssymb}
\usepackage{amsthm}

\defaultfontfeatures{Ligatures=TeX}
\setmainfont{DejaVu Serif}
\setsansfont{DejaVu Sans}
\setmonofont{DejaVu Sans Mono}[Scale=0.85]

\tolerance=2000
\emergencystretch=3em

\geometry{left=2.4cm, right=2.4cm, top=2.5cm, bottom=2.5cm}

\hypersetup{
    colorlinks=true,
    linkcolor=black,
    urlcolor=blue,
    citecolor=black,
    linktoc=all,
}

\usepackage{titlesec}
\titleformat{\chapter}[hang]{\Large\bfseries}{\thechapter.}{18pt}{\Large\bfseries}
\titleformat{\section}[hang]{\large\bfseries}{\thesection.}{12pt}{\large\bfseries}
\titleformat{\subsection}[hang]{\normalsize\bfseries}{\thesubsection.}{10pt}{\normalsize\bfseries}

\setlength{\headheight}{24pt}
\pagestyle{fancy}
\fancyhf{}
\lhead{\small\nouppercase{\leftmark}} 
\cfoot{\thepage}

\lstdefinestyle{code}{
    basicstyle=\small\ttfamily,
    breaklines=true,
    frame=single,
    numbers=left,
    numberstyle=\tiny\color{gray},
    keywordstyle=\color{blue},
    commentstyle=\color{gray},
    stringstyle=\color{red!70!black},
    showstringspaces=false,
    extendedchars=true,
    columns=fullflexible,
    keepspaces=true,
}
\lstdefinestyle{pseudo}{
    basicstyle=\small\ttfamily,
    breaklines=true,
    frame=leftline,
    framerule=1pt,
    rulecolor=\color{gray!50},
    xleftmargin=1em,
    mathescape=true,
    showstringspaces=false,
    extendedchars=true,
    columns=fullflexible,
    keepspaces=true,
}
\lstset{style=code}

\usepackage{tikz}
\usetikzlibrary{arrows.meta,positioning,fit,backgrounds,calc,shapes.geometric}

\usepackage[most]{tcolorbox}
\newtcolorbox{intuition}[1][Интуиция]{colback=gray!5,colframe=gray!45,title={#1},fonttitle=\bfseries,breakable}
\newtcolorbox{pitfall}[1][Ловушка]{colback=red!4,colframe=red!35,title={#1},fonttitle=\bfseries,breakable}
\newtcolorbox{workedbox}[1][Разбор примера]{colback=blue!3,colframe=blue!30,title={#1},fonttitle=\bfseries,breakable}
\newtcolorbox{openq}[1][Открытый вопрос]{colback=orange!6,colframe=orange!60,title={#1},fonttitle=\bfseries,breakable}

\theoremstyle{definition}
\newtheorem{definition}{Определение}[chapter]
\newtheorem{example}{Пример}[chapter]
\newtheorem{principle}{Принцип}[chapter]

\newcommand{\Neigh}{\mathcal{N}}
\newcommand{\conf}{\mathrm{conf}}
\newcommand{\mood}{\mathrm{mood}}
\newcommand{\Reff}{R_{\mathrm{eff}}}
\newcommand{\KL}{D_{\mathrm{KL}}}
\newcommand{\Var}{\mathrm{Var}}
\newcommand{\E}{\mathbb{E}}

\title{Конструирование графовой памяти Рин\\[0.4em]
\large Расширенный учебный курс по архитектуре памяти, динамике, эпистемике и инженерии}
\author{Расширенная учебная версия проектного мануала RIN}
\date{\today}

\begin{document}

\frontmatter
\maketitle

\vspace{1.5em}
\begin{center}
\begin{minipage}{0.84\textwidth}\itshape\small
Эта версия намеренно шире базового учебного дубля. Её задача — не просто перевести
архитектурный замысел в более мягкую педагогическую форму, а развернуть его в курс,
где рядом стоят интуиция, определения, формулы, инженерные ограничения, разборы
сценариев и аппарат проверки. Базовый roadmap сохранён; расширения, которых не было
в исходном тексте в явном виде, вводятся открыто и используются как развитие, а не
как тихая замена исходных решений.
\end{minipage}
\end{center}
\vspace{1.5em}

\tableofcontents

\mainmatter

\chapter{Как читать этот курс}

\section{Не FAQ, а курс}

Эта версия текста сознательно написана не как памятка для уже посвящённых и не как
монография, где читатель обязан заранее жить внутри терминологии. Жанр здесь другой:
\emph{курс по архитектуре памяти}. Это значит, что изложение должно делать не только
две вещи — сообщать решения и хранить корректность, — но и третью: показывать,
\emph{почему} решение вообще появилось, на какой вопрос отвечает и чем рискует.

Отсюда три обязательства текста.
\begin{enumerate}
\item Определение появляется раньше, чем начинает незаметно работать.
\item Формула появляется там, где без неё теряется механизм, а не для украшения.
\item Инженерный компромисс называется компромиссом, а не тайком выдаётся за закон природы.
\end{enumerate}

\section{Три уровня утверждений}

В курсе важно не смешивать три рода высказываний.

\begin{definition}[Архитектурное утверждение]
Утверждение архитектурно, если от него зависит сам класс строящейся системы:
например, что память хранится как явный типизированный граф; что retrieval,
внутренняя жизнь и сон реализуются как политики одного обходчика; что факт и
гипотеза различаются эпистемически, а не только текстово.
\end{definition}

\begin{definition}[Инженерное упрощение]
Утверждение инженерно, если оно выбрано как первое рабочее приближение, но в
принципе может быть заменено без смены класса системы: например, жёсткий порог
пересечения вместо полноценной кластерной консолидации, LSH/MinHash вместо более
тонкой структуры сходства, разреженный Inspector вместо тяжёлого explainability-слоя.
\end{definition}

\begin{definition}[Исследовательская ветвь]
Утверждение относится к исследовательской ветви, если оно не нужно для Stage~1,
но обещает новый класс эффектов: клеточные комплексы, динамическая память, резервуарные
вычисления, квантовые ассоциативные схемы.
\end{definition}

\begin{pitfall}[Самая частая методологическая ошибка]
Если инженерное упрощение обсуждается как будто это и есть архитектура, система
становится хрупкой не из-за механики, а из-за языка описания. Разработчик начинает
бояться менять то, что на самом деле было временным. Обратная ошибка не менее вредна:
если архитектурное решение выдаётся за временную деталь, его начинают размывать
«локальными оптимизациями», и система теряет идентичность.
\end{pitfall}

\section{Что именно расширено по сравнению с базовым дублем}

В этом курсе есть места, где добавлена не только педагогика, но и техника.
Это сделано открыто.

\begin{enumerate}
\item Retrieval расписан не только качественно, но и как оператор на графе: с
матрицей переходов, рестартом, температурой, hitting time и спектральными
ограничениями.
\item Эпистемика развернута в виде каналов свидетельств, а не только одного
скалярного поля уверенности.
\item Экспериментальный раздел усилен явной диагностикой failure modes.
\item Добавлены сквозные разборы, где память прослеживается от записи события до
ревизии производных узлов.
\end{enumerate}

\begin{intuition}[Почему это не «раздуть текст»]
Когда архитектура сложна, краткость бывает двух разных сортов. Бывает краткость
хорошая: выкинули повтор и оставили идею. А бывает краткость опасная: скрыли
промежуточные мосты, и тогда текст выглядит компактным, но держится на том, что
читатель сам за автора достроит половину механизма. Расширенный курс борется именно
со вторым типом краткости.
\end{intuition}

\section{Три прохода чтения}

Для такого текста полезно иметь не один, а три режима чтения.

\paragraph{Первый проход — карта.} Нужно понять, что система вообще считает памятью,
какие в ней есть процессы и где граница между графом, логом, рантаймом и языковой
моделью.

\paragraph{Второй проход — механика.} Здесь уже важны формулы, типы узлов, поля,
операторы консолидации, режимы обходчика и пропагация уверенности.

\paragraph{Третий проход — споры и открытые ветви.} На этом уровне читаются не
столько ответы, сколько развилки: где исходное решение твёрдое, где оно временное,
что можно заменить, а что нельзя, не разрушив замысел.

\section{Главный критерий хорошего текста о памяти}

Хороший текст об архитектуре памяти должен выдерживать двоякую проверку.

\begin{principle}[Два экзамена]
Первый экзамен — \emph{технический}: после чтения должно быть возможно собрать Stage~1,
придумать тесты, найти диагностические переменные и не потерять roadmap. Второй
экзамен — \emph{педагогический}: должно быть понятно, почему эта архитектура вообще
выглядит именно так, а не как список случайно соседствующих трюков.
\end{principle}

Этим двум экзаменам и подчинён весь дальнейший порядок курса.

\chapter{Введение: зачем памяти нужен субстрат}

\section{Постановка задачи}

У локального AI-компаньона есть требование, которое почти невозможно
обойти архитектурной хитростью: он должен вести себя не как одноразовый
ответчик на очередной запрос, а как сущность с накопленной биографией. Это
означает по меньшей мере четыре вещи.

\begin{enumerate}
\item Он должен связывать события между собой, а не хранить их как изолированные записи.
\item Он должен ошибаться не хаотически, а осмысленно: делать предварительные выводы, сомневаться в них и позднее их пересматривать.
\item Он должен меняться во времени: недавнее должно звучать иначе, чем давнее; важное должно переживать неважное; повторяющееся должно сворачиваться в более устойчивые структуры.
\item Он должен вспоминать контекстно: не одинаково по любой команде, а в зависимости от текущего вопроса, эмоционального состояния и уже активированного окружения.
\end{enumerate}

Плоская история сообщений решает лишь малую часть этой задачи. Векторная база
воспоминаний продвигается дальше: она позволяет искать смыслово похожие записи.
Но ни история сообщений, ни набор независимых векторных элементов сами по себе
не дают ответа на более глубокий вопрос: где в памяти живут связи между фактами,
уверенность в них, их происхождение и их эволюция во времени?

\section{Почему здесь появляется граф}

Слово ``граф'' легко превратить в пустую модную вывеску. В этой книге под
графом понимается не престижная структура данных сама по себе, а минимально
естественный ответ на перечисленные выше требования.

Если память должна хранить не только ``что произошло'', но и ``с чем это
связано'', ``насколько этому доверяют'', ``из чего это было выведено'', ``как это
влияет на соседние убеждения'' и ``когда это было важно'', то у архитектуры
почти неизбежно возникают:

\begin{itemize}
\item узлы как носители содержания;
\item рёбра как носители отношений;
\item метаданные на узлах и рёбрах;
\item процессы над графом, а не просто чтение из хранилища.
\end{itemize}

\begin{intuition}[Граф здесь не украшение]
Если убрать из памяти отношения и оставить только список фрагментов текста,
система ещё сможет что-то ``вытаскивать по похожести'', но уже не сможет внятно
ответить, почему одно суждение усиливает другое, откуда возникла гипотеза, что
следует пересмотреть после опровержения факта и по какому маршруту вообще
течёт активация. Иначе говоря, в такой системе не останется самой механики
памяти; останется только склад.
\end{intuition}

\section{Центральный тезис книги}

\begin{principle}[Главная формула всего курса]
Память Рин рассматривается как один типизированный темпоральный граф с
метаданными на узлах и рёбрах. Retrieval, обход, вывод, консолидация,
забывание, ночная переработка, настроение, личность и языковое выражение
трактуются как операторы над этим субстратом или как сцепленные с ним
процессы, а не как набор несогласованных подсистем с разными моделями мира.
\end{principle}

Именно этот тезис задаёт порядок изложения. Сначала вводится архитектурная
карта системы (глава 3). Затем даётся словарь основных объектов (глава 4).
После этого разбирается субстрат памяти (глава 5) и его многомасштабная
динамика (глава 11). Далее рассматриваются основные процессы: настроение
(глава 7), обход графа (глава 6), fusion и консолидация (глава 10), эпистемика
и ревизия убеждений (главы 13 и 14), влияние топологии (глава 8), ночная
работа (глава 12), слой личности и голоса (главы 16 и 17), инженерный рантайм
(глава 18) и, наконец, экспериментальная проверка (глава 22).

\section{Четыре процесса, которые нельзя настраивать по отдельности}

При первом знакомстве с задачей возникает соблазн мыслить модулями: сначала
``сделать retrieval'', потом ``добавить забывание'', потом ``приделать сон''.
Такой подход удобен управленчески, но почти всегда ошибочен содержательно.
В памяти Рин особенно тесно сцеплены четыре процесса.

\begin{enumerate}
\item Извлечение релевантного. Что именно всплывает в ответ на запрос, внутренний импульс или текущую сцену.
\item Пропагация уверенности. Как новый факт усиливает, ослабляет или ставит под сомнение соседние убеждения.
\item Забывание и консолидация. Что тускнеет, что переживает время, а что сворачивается в обобщение.
\item Ночная переработка. Как память в фоне инспектирует сама себя, чинит противоречия и перестраивает структуру.
\end{enumerate}

Эти процессы читают и пишут одни и те же поля одного и того же графа.
Изменение уверенности меняет то, что будет извлечено. Изменение топологии
меняет маршруты обхода. Консолидация меняет последующую пропагацию. Ночной
скраб меняет судьбу гипотез, созданных дневными операторами.

\begin{pitfall}[Ошибка независимой настройки]
Если каждый оператор проектируется как будто остальные существуют где-то
рядом, но не вмешиваются в его поведение, то система может хорошо работать в
локальном тесте и при этом распадаться в общем цикле. Например, ``удачный
retrieval'' на фиксированном графе может оказаться неудачным после первой же
волны консолидации; ``аккуратная пропагация'' может превратиться в лавину при
иной топологии; ``красивое забывание'' может уничтожить материал, без которого
ревизия убеждений станет невозможной.
\end{pitfall}

Именно поэтому завершающая часть книги посвящена не красивым примерам, а
экспериментальному стенду с контролируемой истиной (глава 22). Для такой
архитектуры особенно важно мерить не изолированные фрагменты поведения, а их
совместную физику.

\section{Чего эта книга не утверждает}

Текст следует читать как рабочую архитектурную теорию памяти для RIN, а не
как финальную общую теорию разума. В нём есть три разных слоя утверждений, и
их полезно различать с самого начала.

\begin{enumerate}
\item Твёрдые инженерные решения. Например, требование единого субстрата или различение активации и уверенности.
\item Рабочие гипотезы. Например, конкретные формы обхода, консолидации или эпистемических предохранителей.
\item Дальние линии развития. Например, переход к гиперграфам, клеточным комплексам или иным более выразительным субстратам.
\end{enumerate}

\begin{intuition}[Граница между моделью и реальностью]
Ни одна архитектура памяти не обязана быть истинной в сильном смысле слова;
она обязана быть содержательно плодотворной, инженерно управляемой и
проверяемой на стенде. Поэтому в этой книге важнее не вопрос ``похоже ли это
на окончательную истину о памяти'', а вопрос ``даёт ли такая схема хорошие
механизмы, понятные сбои и измеримые компромиссы''.
\end{intuition}

\section{Как читать этот курс}

Книга задумана так, чтобы не требовать от читателя постоянного угадывания,
``что именно автор имел в виду''. Ради этого соблюдаются два правила.

\begin{enumerate}
\item Новый объект по возможности сначала мотивируется, затем называется, затем формализуется.
\item Формула появляется там, где без неё уже становится неточно, а не там, где удобно сразу обозначить символы.
\end{enumerate}

Если нужен быстрый вход в систему, достаточно прочитать главы 3, 4, 5 и 20.
Если нужен полный контур, текст лучше читать по порядку: в этой версии книги
порядок глав соответствует логике постепенного усложнения модели.

\chapter{Память как система: первый обзор}
\label{ch:architecture}

\section{Зачем нужна отдельная карта системы}

Прежде чем обсуждать отдельные механизмы, полезно увидеть всю систему как один
замкнутый цикл. Иначе почти неизбежно возникает ложное ощущение, что память
состоит из набора независимых блоков: где-то отдельно лежит граф, где-то
отдельно работает retrieval, где-то отдельно находится настроение, а языковая
модель ``просто забирает контекст''. Для архитектуры памяти это опасная
иллюзия. Почти каждый блок влияет на остальные двояко: он не только читает
состояние системы, но и перестраивает его.

\section{Что именно строится}

Рассматривается локальный AI-компаньон, память которого должна обладать
следующими свойствами:

\begin{itemize}
\item быть долговременной, а не ограниченной текущим окном контекста;
\item быть ассоциативной, а не только хронологической;
\item различать наблюдение, гипотезу, предпочтение, реактивный паттерн и состояние отношений;
\item допускать пересмотр убеждений;
\item учитывать эмоциональную и личностную окраску того, что всплывает;
\item поддерживать фоновую переработку памяти, а не только обслуживание очередного запроса.
\end{itemize}

Чтобы эти свойства не рассыпались по разным хранилищам и разным схемам
данных, вводится единый субстрат.

\begin{principle}[Единый субстрат]
Память Рин моделируется как один типизированный темпоральный граф с
метаданными на узлах и рёбрах. Все основные процессы памяти либо
непосредственно работают на этом графе, либо сцеплены с ним так, что их
состояние можно выразить через чтение и запись к его полям.
\end{principle}

\section{Четыре слоя системы}

Если смотреть на систему сверху, в ней удобно различать четыре слоя.

\begin{enumerate}
\item Субстрат. Сам граф памяти: узлы, рёбра, уверенность, провенанс, валентность, временной слой, важность и другие поля (глава 5).
\item Операторы. Обход, retrieval, fusion, консолидация, пропагация, забывание, сон и скраб. Это динамика памяти как таковой (главы 6, 10, 11, 14, 12).
\item Состояния. Настроение, эпистемические статусы, вспомогательная метакогнитивная динамика. Эти сущности не следует путать с графом, но они двусторонне с ним связаны (главы 7, 13).
\item Выражение. Слой, который превращает внутреннюю структуру в реплику, тон, фокус внимания и устойчивый характер (главы 16, 17).
\end{enumerate}

\begin{center}
\begin{tikzpicture}[font=\small,>=Stealth,
  layer/.style={draw,rounded corners,minimum width=82mm,align=center,inner sep=6pt,fill=gray!5},
  lbl/.style={font=\bfseries}]
  \node[layer] (expr) {{\bfseries Выражение}\\[1pt]
     \footnotesize личность \quad$\cdot$\quad голос \quad$\cdot$\quad языковая композиция};
  \node[layer,below=7mm of expr] (state) {{\bfseries Состояния}\\[1pt]
     \footnotesize настроение \quad$\cdot$\quad эпистемические статусы \quad$\cdot$\quad метакогнитивные сигналы};
  \node[layer,below=7mm of state] (op) {{\bfseries Операторы}\\[1pt]
     \footnotesize retrieval $\cdot$ обход $\cdot$ fusion $\cdot$ консолидация\\
     \footnotesize пропагация $\cdot$ забывание $\cdot$ сон};
  \node[layer,below=7mm of op] (sub) {{\bfseries Субстрат}\\[1pt]
     \footnotesize граф памяти: узлы, рёбра, типы, метаданные};
  \draw[<->,thick] (expr)--(state);
  \draw[<->,thick] (state)--(op);
  \draw[<->,thick] (op)--(sub);
\end{tikzpicture}
\captionof{figure}{Общая карта системы. Стрелки двусторонние: слой не просто потребляет нижележащий, но и изменяет его состояние.}
\label{fig:layers}
\end{center}

\begin{pitfall}[Чего нельзя делать с этой схемой]
Эти четыре слоя нельзя понимать как четыре независимые базы или четыре
микросервиса с отдельными картинами мира. Слой выражения не должен
придумывать собственную память поверх графа. Слой состояний не должен
дублировать содержательное знание в виде параллельного скрытого графа.
Операторы не должны тайно хранить свою ``локальную истину'', не
согласованную с общим субстратом.
\end{pitfall}

\begin{pitfall}[``Один субстрат'' --- дисциплина, а не догма ``всё есть граф'']
Тезис о едином субстрате легко перетянуть слишком далеко и начать складывать
в граф вообще всё: runtime-очереди, streaming-флаги, транзиентные тики
настроения, состояние планировщика. Это превращает граф не в память, а в
свалку состояний.

\[
\begin{gathered}
\text{граф} = \text{память}, \qquad
\text{event log} = \text{история},\\
\text{телеметрия} = \text{самонаблюдение}, \qquad
\text{scheduler state} = \text{управление},\\
\text{prompt context} = \text{временная сборка для LLM}.
\end{gathered}
\]

Эти сущности связаны, но это разные хранилища с разными законами жизни.
Настроение --- сцепленный сбоку модуль; самотелеметрия метакогниции в
инженерном варианте естественно оформляется как кольцевой буфер; планировщик
--- как конечный автомат режимов работы. Правило простое: если данное
транзиентно или относится к управлению, истории или упаковке контекста, то
оно не должно механически жить в графе памяти.
\end{pitfall}

\section{Что течёт по системе}

Через эту архитектуру проходят несколько разных ``величин'', и их полезно не
смешивать.

\begin{definition}[Активация]
Активация --- это краткоживущая подсветка узла: мера того, насколько он важен
сейчас, в данном локальном ходе памяти.
\end{definition}

\begin{definition}[Уверенность]
Уверенность --- это медленно меняющаяся степень убеждённости в содержании
узла. Она не должна путаться с активацией: можно быть уверенным в факте и при
этом не держать его в фокусе; можно, наоборот, временно активировать
сомнительную гипотезу.
\end{definition}

\begin{definition}[Валентность]
Валентность --- аффективная окраска содержания. Она отвечает не за истинность
узла, а за то, как этот узел переживается и какую роль играет в настроенческой
динамике.
\end{definition}

\begin{definition}[Структурное изменение]
Структурное изменение --- это рождение, укрепление, ослабление, слияние или
удаление узлов и рёбер. В отличие от активации, оно меняет саму геометрию
памяти.
\end{definition}

\section{Сквозной пример}

Полезно проследить полный цикл на одном эпизоде. Пусть во время стрима
зритель пишет: ``Опять слил ранг в доте''.

\begin{workedbox}[От реплики до перестройки памяти]
\begin{enumerate}
\item Кодирование события. В систему поступает наблюдение о конкретном эпизоде. Оно получает тип, базовую уверенность, временной слой и привязку к участнику взаимодействия.
\item Поиск якорей. По содержанию и контексту сцены находятся ближайшие узлы: прежние упоминания доты, прошлые реплики этого зрителя, возможные эмоциональные реакции на проигрыш.
\item Расширение окрестности. От якорей запускается ограниченный обход. Он поднимает не только прямые совпадения, но и соседние ассоциативные структуры.
\item Локальный вывод. Совместная активация нескольких узлов может породить новую гипотезу: например, что в данной сцене уместен мягкий подкол, а не сухое сочувствие.
\item Выражение. На сцену выходит слой личности и голоса: из всплывшего материала строится конкретная реплика в характерном стиле.
\item Обратная запись. Реакция зрителя становится частью новых данных. Если подкол сработал, укрепляется соответствующий паттерн; если нет, появляется основание для пересмотра.
\item Фоновая переработка. Позднее эта сцена может войти в более устойчивое знание: ``данный зритель часто шутит о рейтинге'', ``такая интонация обычно безопасна'', ``эти две темы часто связаны''.
\end{enumerate}
\end{workedbox}

Этот пример показывает главный архитектурный факт: реплика является не
конечной целью отдельного генератора, а одним моментом внутри цикла памяти.
Система не просто говорит; она одновременно перестраивает себя.

\section{Что не входит в граф}

Тезис о едином субстрате не означает, что ``всё вообще'' надо хранить в
памяти-графе.

\begin{itemize}
\item Граф памяти хранит долговременное содержательное знание и его обоснования.
\item Event log хранит историю коммитов и удалений.
\item Телеметрия хранит самонаблюдение и служебные сигналы.
\item Планировщик хранит режимы работы и бюджет вычислений.
\item Промпт-контекст хранит временную сборку того, что подаётся в LLM на конкретном шаге.
\end{itemize}

\begin{intuition}[Почему это различение принципиально]
Если event log, граф, телеметрия и prompt context свалить в одну кучу,
получится не богатая память, а неразличимая смесь истории, управления и
содержания. Такой объект неудобен не только технически. Он делает почти
невозможным ответ на вопрос, что именно является памятью, а что лишь следом
работы памяти.
\end{intuition}

\section{Роль архитектурной главы в остальной книге}

Эта глава не заменяет дальнейшие определения и механизмы. Её задача скромнее:
задать правильную оптику. После неё должно быть ясно не ``как реализуется
каждая функция'', а:

\begin{itemize}
\item что в системе считается субстратом;
\item какие процессы над ним текут;
\item почему они не независимы;
\item где проходит граница между памятью, состоянием и выражением.
\end{itemize}

Следующий шаг поэтому естественен: после общей карты нужен словарь основных
объектов. Этой задаче посвящена глава 4.

\chapter{Основные объекты памяти}
\label{ch:definitions}

\section{Зачем вообще нужна отдельная глава понятий}

В технических текстах о памяти часто возникает неприятная ситуация: читатель
уже встретил ``уверенность'', ``активацию'', ``провенанс'', ``консолидацию'' и
``walker'', но ещё не получил компактной картины того, как эти слова
соотносятся друг с другом. Формально каждое из них потом где-то определяется,
однако понимание успевает распасться на локальные фрагменты.

Задача этой главы обратная: собрать в одном месте основные объекты, не
пытаясь здесь же полностью исчерпать все механизмы. Поэтому формулировки
ниже намеренно коротки. Их роль --- дать рабочий словарь и развести сущности,
которые особенно легко спутать.

\section{Объекты субстрата}

\begin{definition}[Граф памяти]
Граф памяти $G = (V, E)$ --- направленный динамический граф, в котором узлы
$V$ несут содержательные единицы памяти, а рёбра $E$ несут отношения между
ними. Граф называется типизированным, потому что и узлы, и рёбра имеют виды;
темпоральным, потому что память различает временные слои и возраст
содержания; динамическим, потому что операторы могут менять и узлы, и рёбра.
\end{definition}

\begin{definition}[Узел памяти]
Узел памяти $v \in V$ --- это единица содержания, с которой можно работать как
с объектом: наблюдение, факт, эпизод, гипотеза, предпочтение, паттерн
реакции, модель отношения, элемент самопредставления и так далее. Узел хранит
не только текст или embedding, но и поля состояния: тип узла, уверенность,
провенанс, временной слой, активацию, валентность, важность и служебные
флаги.
\end{definition}

\begin{intuition}[Почему узел --- это больше, чем ``запись'']
Если узел понимать только как текстовый фрагмент, память превращается в
коллекцию документов. Узел в этой архитектуре ближе к ``объекту мышления'':
он не просто что-то содержит, но ещё может быть активирован, усилен, ослаблен,
проверен, связан с посылками и включён в дальнейший вывод.
\end{intuition}

\begin{definition}[Тип узла]
\texttt{node\_kind} --- вид узла, определяющий, как именно его следует читать и
обрабатывать. В частности, именно тип узла задаёт смысл поля confidence. Для
наблюдаемого факта это степень доверия к истинности; для предпочтения ---
устойчивость оценки; для реактивного паттерна --- надёжность модели ответа;
для mood-summary это может быть вообще не числовая истина, а служебная
свёртка состояния.
\end{definition}

\begin{pitfall}[Почему тип узла не декоративен]
В памяти проекта сосуществуют по меньшей мере observation, claim, hypothesis,
episode, preference, reaction\_pattern, relationship\_state, boundary,
skill/procedure, self\_model и mood-summary. Эти сущности нельзя честно
измерять одной и той же шкалой: ``уверенность 0.8'', относящаяся к факту, и
``уверенность 0.8'', относящаяся к предпочтению, означают разные инженерные
вещи.
\end{pitfall}

\begin{definition}[Ребро]
Ребро $(u \to v)$ --- направленная связь между двумя узлами, несущая тип
отношения. Примеры: таксономическое включение, причинная поддержка,
противоречие, ассоциация, участие в одном эпизоде, отношение к одному
участнику взаимодействия.
\end{definition}

\begin{definition}[Два канала ребра]
На ребре живут два разных числовых канала:
\begin{itemize}
\item вес перехода $w_{uv} \ge 0$, который нужен обходу и говорит, насколько охотно маршрут пойдёт по этой связи;
\item знак $s_{uv} \in \{+1, -1\}$, который нужен эпистемике и говорит, усиливает ли эта связь убеждение или ослабляет его.
\end{itemize}
Эти каналы не следует смешивать: маршрут может часто ходить по ребру, которое
при этом вносит отрицательный вклад в уверенность, и наоборот.
\end{definition}

\begin{definition}[Провенанс (обоснование)]
Провенанс --- это информация о происхождении узла: из какого наблюдения,
события, источника или набора посылок он возник. Для выводных узлов
провенанс играет роль обоснования, то есть justification set. Важно, что
provenance --- это не полный трек обхода: маршрут walker'а шумен, может
храниться отдельно и сам по себе не является обоснованием. Именно provenance
делает возможной ретракцию: если опровергнута посылка, можно найти и её
следствия.
\end{definition}

\begin{definition}[Время: системное и содержательное]
Время в памяти удобно разводить на два канала. Системное время --- это
минимальный набор меток, нужный самим операторам: \texttt{layer} (сцена,
сессия, день, долгий паттерн) и, как правило, одна служебная метка вроде
времени кодирования или последнего доступа. Содержательное время --- это ответ
на вопрос ``когда это произошло'' внутри самого мира памяти. Оно может
моделироваться отдельными узлами и быть реконструктивным, а значит, и
ошибочным. Интервалы валидности (\texttt{start}/\texttt{end}) имеет смысл
вешать лишь на те факты, для которых это содержательно нужно, а не
механически на каждый узел.
\end{definition}

\section{Поля состояния}

\begin{definition}[Уверенность]
Уверенность $\conf(v) \in [-1, 1]$ --- знаковая степень убеждённости в
содержании узла. Положительный знак означает ``скорее да'', отрицательный ---
``скорее нет'', нуль соответствует незнанию или эпистемической
неопределённости. Такая запись удобна тем, что различает два принципиально
разных состояния: ``неверие'' и ``отсутствие знания''.
\end{definition}

\begin{equation}
\conf = 2P - 1,
\end{equation}

где $P$ --- эквивалентная вероятность истинности. Это соответствие полезно, но
в книге предпочтителен именно знаковый скаляр: он лучше согласуется с
противоречием и ретракцией. Практически удобно считать, что зона незнания
задаётся условием $|\conf(v)| < h$ с грубым ориентиром $h \approx 0{,}3$.
Именно такие узлы являются естественными кандидатами на
\texttt{need\_to\_discover}, пересмотр или удаление. Уверенность меняется
медленно: через пропагацию и ревизию убеждений, а не при каждом локальном
подсвечивании. Принципиальное ограничение: выводный узел не получает
$|\conf| = 1$.

\begin{pitfall}[Чего нельзя требовать от confidence]
Одно число не обязано кодировать всю эпистемическую жизнь узла. Оно не
заменяет тип узла, не заменяет провенанс и не заменяет внешние механизмы
неопределённости, такие как Beta-модель реакций или флаги \texttt{need\_to\_*}.
Поэтому вопрос ``почему бы не хранить всё в одном идеальном скаляре'' почти
всегда поставлен неверно.
\end{pitfall}

\begin{definition}[Активация]
Активация $a_v(t) \ge 0$ --- краткоживущая мера того, насколько узел
``подсвечен сейчас''. Активация относится к локальной динамике памяти, а не к
долговременной вере. Она может возрастать из-за запроса, обхода, совместной
активности соседей или внутренней заинтересованности системы.
\end{definition}

\begin{equation}
a(t + 1) = \lambda \tilde{P}^{\top} a(t) + s(t), \qquad \lambda \in (0, 1),
\end{equation}

где $\tilde{P}$ --- матрица переходов, а $s(t)$ --- внешний или внутренний
вброс активации. Эквивалентно об этом можно думать как об ожидаемой частоте
посещений узла роем в коротком окне. Активацию важно не путать ни с
уверенностью, ни со стационарным распределением $\pi$: первое отвечает за
убеждение, второе --- за долгосрочную геометрию потока.

\begin{definition}[Валентность и настроение]
Валентность --- аффективная окраска узла. Она отвечает не за то, насколько узел
истинный, а за то, как он переживается и как он взаимодействует с настроением.
В книге обычно предполагается двумерное описание ``валентность--возбуждение'',
но точная параметризация вторична по отношению к самой идее аффективной
несущей. Настроение \texttt{mood} --- отдельный низкоразмерный вектор
состояния, сцепленный с графом двусторонне, но не живущий внутри каждого узла
по отдельности.
\end{definition}

\begin{definition}[Значимость]
Значимость (\texttt{importance}, \texttt{salience}) --- мера того, насколько
объект памяти важен для дальнейшей судьбы системы. Значимость влияет на
консолидацию, забывание и приоритеты переработки. Она может частично зависеть
от эмоционального сдвига, от роли объекта в долгих паттернах и от структурной
центральности. В качестве дешёвого прокси значимости удобно использовать
величину изменения настроения $|\Delta \mood|$; более глубокий, но и более
дорогой сигнал --- ошибка предсказания.
\end{definition}

\section{Операторы над памятью}

\begin{definition}[Random Walker]
Random Walker --- блуждающий агент, который движется по графу по
вероятностному правилу, зависящему от структуры рёбер, релевантности,
текущего настроения, неопределённости и других факторов. В типовой форме:
\[
P(A \to B) \propto w_{AB} \cdot rel(B) \cdot mc(B) \cdot unc(B)^{\beta_u},
\]
где $w_{AB}$ --- вес перехода, $rel(B)$ --- релевантность узла текущему
запросу или сцене, $mc(B)$ --- mood-congruent bias, а $unc(B)^{\beta_u}$ ---
смещение к неопределённому содержанию в тех режимах, где это нужно. Один
walker даёт локальную траекторию; множество walker'ов даёт статистическую
картину потока по памяти.
\end{definition}

Функционал одного walker'а в минимуме таков: (i) перейти по вероятностному
правилу; (ii) с вероятностью $\alpha$ телепортироваться, то есть
перезапуститься из другого узла; (iii) повысить активацию посещённого узла;
(iv) по условию запустить операцию вроде retrieval, fusion или инспекции
\texttt{need\_to\_*}; (v) записать след как последовательность посещённых узлов.

\begin{definition}[Рой обходчиков]
Рой (\texttt{swarm}) --- набор walker'ов, работающих параллельно под общей
политикой. Именно рой, а не одиночный маршрут, лучше соответствует идее
распределённой активации памяти. На достаточной статистике следов рой даёт и
эмпирическую картину маршрутов по графу, что особенно важно для
экспериментов.
\end{definition}

\begin{definition}[Три политики обхода]
Один и тот же движок обхода должен уметь работать как минимум в трёх режимах:
retrieval под запрос, внутренняя жизнь / любопытство с повышенным интересом к
неопределённому и сон, где включается более широкое стохастическое покрытие с
телепортом. Эти режимы различаются не типом сущности, а параметрами той же
машины.
\end{definition}

\begin{definition}[Retrieval]
Retrieval --- извлечение содержательно релевантной части памяти. Здесь это не
один шаг, а двухходовая процедура: сначала находятся якоря быстрым индексом,
затем от якорей запускается ограниченное расширение окрестности посредством
обхода. Этот порядок принципиален: walker хорош для graph expansion, но плох
как механизм поиска входа. В результате retrieval зависит не только от
семантической близости запроса, но и от настроения, то есть допускает
mood-congruent recall.
\end{definition}

\begin{definition}[Fusion]
Fusion --- механизм рождения нового узла из согласованной конфигурации
совместно активированных посылок. Обычно результат fusion --- гипотеза или
обобщение с ограниченной уверенностью и явным провенансом. Для выводных
узлов принципиально: результат fusion не стартует как ``окончательный факт'',
а сохраняет провизорный статус до дополнительной проверки.
\end{definition}

\begin{definition}[Консолидация]
Консолидация --- процесс перехода от множества частных следов к более
устойчивым структурам памяти. Она может включать объединение дублей,
summary-overlay, рождение обобщений и более позднее огрубление графа. Важно,
что консолидация --- не один жёсткий оператор, а семейство подопераций:
fusion/вывод, summary-overlay, dedup/merge и только на более позднем этапе
разрушающее огрубление.
\end{definition}

\begin{definition}[Забывание]
Забывание --- ослабление или удаление тех элементов памяти, которые не
поддерживаются ни недавним использованием, ни значимостью, ни более глубокой
структурной ролью. В этой архитектуре забывание не противоположно
консолидации, а образует с ней единую динамику перераспределения памяти во
времени. На базовом уровне его удобно связывать с частотой и свежестью
доступа, то есть с медленно меняющейся базовой активацией.
\end{definition}

\begin{definition}[Сон]
Сон --- фоновый режим памяти, в котором обход становится более стохастическим
и более широким, чем в дневном retrieval. В этом режиме происходят скраб,
пересмотр противоречий, консолидация, случайные инспекции и иногда
неожиданные дальние fusion-связки. Тот же режим удобно использовать и для
подтягивания внешнего материала по флагам \texttt{need\_to\_discover}.
\end{definition}

\section{Сцепленные с памятью процессы}

\begin{definition}[Настроение]
Настроение \texttt{mood} --- отдельный вектор состояния, который не хранится
внутри каждого узла, но двусторонне сцеплен с графом. Настроение влияет на
маршруты активации и на форму выражения; граф, в свою очередь, поставляет
настроению материал и обратную связь.
\end{definition}

\begin{definition}[Личность]
Личность --- не таблица жёстко прошитых весов, а способ, которым мотивация,
настроение, эпистемическая осторожность и поведенческие приоритеты смещают
работу памяти и её последующее выражение. Внутри архитектуры личность
проявляется не как отдельный \texttt{store}, а как \texttt{bias} во время обхода
и как стиль на выходе.
\end{definition}

\begin{definition}[Голос]
Голос --- языковой слой, который превращает внутренний контекст памяти в
конкретную реплику. Голос не хранит сами факты вместо памяти, а композирует
высказывание из того, что память уже подняла и подготовила. Часть характера
живёт именно здесь; факты, \texttt{provenance} и эпистемические отношения
живут в графе.
\end{definition}

\section{Стартовое состояние и инженерная перспектива}

\begin{definition}[Эмбриональный граф]
Эмбриональный граф --- стартовое заполнение памяти из внешних баз знаний и
здравого смысла, прежде всего в духе ConceptNet + ATOMIC. Это не готовая
личность и не законченная память, а начальная, нейтральная по характеру среда,
в которой могут запускаться собственные механизмы retrieval, fusion,
консолидации и ревизии.
\end{definition}

\begin{definition}[Stage 1]
Stage 1 --- первый рабочий вертикальный срез системы, в котором теоретическая
модель памяти доводится до минимально наблюдаемой инженерной машины: с
\texttt{event log}, \texttt{write-path}, \texttt{typed node store}, retrieval,
\texttt{prompt packing}, инспектором памяти и синтетическими тестами
(глава 20).
\end{definition}

\begin{pitfall}[Почему не вводится более богатая эпистемическая тройка]
В исходной архитектуре сознательно удерживается знаковый скаляр
$\conf \in [-1, 1]$, а не полный кортеж belief/disbelief/uncertainty. Это
решение не объявляется окончательно истинным: более богатая субъективная
логика остаётся допустимой альтернативой, если разнесённые сигналы
неопределённости окажутся неудобными. Но на текущем этапе roadmap делает
ставку именно на скаляр confidence, provenance, node\_kind и дополнительные
специальные поля неопределённости там, где они реально нужны.
\end{pitfall}

\section{Что должно остаться после этой главы}

После этой главы полезно удерживать три различения.

\begin{enumerate}
\item Узел не равен строке текста. Это содержательный объект с историей и состоянием.
\item Активация не равна уверенности. Первая отвечает за фокус сейчас, вторая --- за долговременное убеждение.
\item Оператор не равен полю хранения. Retrieval, fusion, консолидация и сон --- это процессы над графом, а не разные виды записей в одном и том же столбце.
\end{enumerate}

На этом словарь готов. Следующая глава уже может говорить не об именах
объектов, а о самом субстрате памяти как инженерной конструкции: какие поля
действительно нужны графу и почему именно они.

\chapter{Субстрат: типизированный темпоральный граф}
\label{ch:substrate}

\section{Что такое узел и что такое ребро}

Память Рин --- направленный граф $G = (V, E)$, в котором узлы несут содержание
(воспоминания, факты, сущности, концепты), а рёбра --- типизированные и
знаковые отношения между ними. Ключевая мысль главы: и узлы, и рёбра несут
метаданные, и именно метаданные, а не сам факт связи, делают возможными все
дальнейшие механизмы.

\begin{definition}[Узел памяти]
Узел $v \in V$ --- это запись с содержанием $x_v$ и набором полей состояния,
над которыми работают операторы памяти. Узел одновременно является и
``данными'' (что помнится), и ``убеждением'' (насколько это считается
истинным), и ``следом'' (насколько активно и свежо).
\end{definition}

\subsection{Схема узла}

\begin{longtable}{p{0.24\textwidth} p{0.10\textwidth} p{0.56\textwidth}}
\toprule
\textbf{Поле} & \textbf{Тип} & \textbf{Смысл и кто с ним работает} \\
\midrule
\endhead
\texttt{id} & int & идентификатор \\
\texttt{content} & text & содержание в естественном виде \\
\texttt{embedding} $x_v$ & vec$[d]$ & векторное представление; retrieval-релевантность, мост баз \\
\texttt{node\_kind} & enum & вид узла (observation\slash claim\slash hypothesis\slash preference\slash \dots); задаёт смысл \texttt{confidence} (см.\ раздел~\ref{sec:nodekind}) \\
\texttt{confidence} $c_v$ & $[-1, 1]$ & для truth-apt узлов --- знаковое убеждение ($+$да/$-$нет/$0$хз); для остальных --- надёжность оценки (раздел~\ref{sec:nodekind}) \\
\texttt{valence} & vec$[2]$ & валентность--возбуждение; сцепка с настроением \\
\texttt{layer}\slash\texttt{ts} & enum\slash time & масштаб (сессия\dots жизнь) и/или метка времени; консолидация \\
\texttt{provenance} & list & источники и/или обоснования; TMS, ревизия \\
\texttt{need\_to\_*} & flags & ``исследовать\slash пересмотреть''; любопытство, скраб \\
\texttt{importance} & $[0, 1]$ & значимость: ставится salience при кодировании (гл.~\ref{ch:personality}), подкрепляется центральностью; гейтит прунинг (гл.~\ref{ch:timescales}) \\
\texttt{domains} & tags & членство в доменах знаний; несёт importance по членству, а не диффузией (гл.~\ref{ch:timescales}) \\
\texttt{activation} $a_v$ & float & транзиентная активация; обход \\
\texttt{degree}\slash\texttt{centrality} & float & производные структурные; топология \\
\bottomrule
\end{longtable}

\subsection{Схема ребра}

Ребро $(u \to v)$ несёт направленность (\texttt{src}, \texttt{dst}), тип
отношения \texttt{type} (\textsc{isA}, \textsc{partOf}, \textsc{supports},
\textsc{contradicts}, \textsc{associates}, \dots) и два независимых числовых
канала: вес перехода $w_{uv} \ge 0$ (нормируется по узлу, для обхода) и знак
$s_{uv} \in \{+1, -1\}$ (полярность отношения, для пропагации уверенности).
Знак --- хранимый лейбл, внутреннее свойство связи: ``эта связь за или
против?'', а не вычисляемая величина. Его нельзя ``посчитать по пути'', он
известен при создании ребра. Изменение уверенности $\Delta \conf$ --- уже
производное, оно вычисляется в пропагации с использованием знака
(глава~\ref{ch:propagation}).

\begin{intuition}[Веса --- про свидетельство, а не про характер]
Забегая вперёд (глава~\ref{ch:personality}): вес ребра выражает
\emph{эвиденциальную} силу связи и уверенность, а не ``характер''. Соблазн
задать личность статическими множителями на рёбра-реакции
(``грусть$\to$сарказм $\times 2{,}0$'') мы разберём и отвергнем: характер
должен входить как смещение во время обхода от текущего состояния, а не как
вшитый в граф множитель.
\end{intuition}

\section{Почему рёбра направленные, типизированные и знаковые}

Плоская модель ``узлы и связи одинакового сорта'' ломается сразу на трёх
вещах: направленность (асимметрия причинности и времени), тип (нельзя валить
covalent-связь, близость, ``подтверждает'' и ``является частью'' в один
линейный канал) и знак (без него не отличить ``всё усилилось'' от ``всё
перевернулось'').

Правильный оператор поэтому не общий лапласиан сглаживания, а
\emph{relation-specific} передача сообщения. Сообщение по ребру:
\begin{equation}
m_{u \to v} = \phi_{r_{uv}}\!\big(h_u,\, h_v,\, e_{uv},\, \Delta t,\, c_{uv},\, s_{uv}\big),
\label{eq:message}
\end{equation}
где $r_{uv}$ --- тип отношения, $\Delta t$ --- возраст или временная
дистанция, $c_{uv}$ --- уверенность, $s_{uv}$ --- источник или эпистемический
статус. Затем узел агрегирует входящие сообщения --- и здесь важно не одну
статистику, а набор (сумма, среднее, максимум, дисперсия по типам отношений),
потому что для памяти существенны и кратность подтверждений, и пропорции, и
наличие экстремального свидетельства.

\begin{workedbox}[Почему одна статистика теряет смысл]
Ключевое свойство \emph{среднего} --- инвариантность к дублированию. Пусть про
факт есть свидетельства $\{+, -\}$ (один за, один против) и в другом случае
$\{+, +, -, -\}$ (двое за, двое против). Считая $+ = 1$, $- = -1$, получаем
\[
\mathrm{mean}\{+, -\} = 0 = \mathrm{mean}\{+, +, -, -\}.
\]
Среднее одинаково, хотя во втором случае свидетелей вдвое больше. Для памяти
это принципиально: ``четверо разделились'' --- не то же, что ``один на один''.
\emph{Сумма} же сохраняет кратность, а \emph{максимум} по знаку --- наличие
сильного контрсвидетельства. Отсюда правило: агрегатор хранит \emph{набор}
статистик (сумма, среднее, максимум, дисперсия), а не одну, --- иначе
теряется либо кратность, либо пропорция, либо экстремум.
\end{workedbox}

\begin{intuition}[Аналогия с электрической цепью]
Полезно думать о графе как о цепи с проводниками разного сорта.
Направленность --- это диоды; тип ребра --- материал проводника со своей
проводимостью; знак --- это полярность. Плоская модель --- цепь, где все
элементы заменены одинаковыми резисторами: она проведёт ток, но потеряет всю
логику схемы.
\end{intuition}

\begin{center}
\begin{tikzpicture}[>=Stealth,font=\footnotesize,
  mem/.style={draw, rounded corners, minimum height=7mm, inner sep=3pt},
  inf/.style={draw, rounded corners, fill=blue!5, minimum height=7mm, inner sep=3pt}]
  \node[mem] (a) {разбил кружку};
  \node[mem, below=9mm of a] (d) {юзер хотел новую};
  \node[inf, right=26mm of a] (r) {реакция: грусть};
  \node[inf, right=22mm of d] (x) {вывод: уместно облегчение};
  \draw[->,thick] (a) -- node[above,font=\scriptsize]{$+$ xReact} (r);
  \draw[->,dashed] (a) -- (x);
  \draw[->,dashed] (d) -- (x);
  \draw[->,thick,red!60!black] (x) -- node[right,font=\scriptsize]{$-$ ослабляет} (r);
\end{tikzpicture}
\captionof{figure}{Противоречие не самостоятельное, а \emph{контекстное}. Базовое $+$xReact ``разбил $\to$ грусть'' ослабляется не узлом ``хотел новую'' самим по себе, а совместным выводом ``разбил $+$ хотел новую $\to$ уместно облегчение'' (пунктир --- провенанс; это $n$-арная конфигурация, глава~\ref{ch:fusion}). Именно она делает возможным уместное ``поздравляю''.}
\label{fig:cup}
\end{center}

\section{Виды узлов и что значит уверенность}
\label{sec:nodekind}

Раньше мы говорили: ``узел --- это одновременно данные и убеждение''. Красиво,
но для реализации грубо: не всякий узел --- утверждение об истинности.
``Португалия имеет такой-то режим магазинов'', ``этот вывод поддержан
источником'', ``гипотеза House / Cast Away ложна'' --- да, это truth-apt
утверждения, и для них $\conf \in [-1, 1]$ означает истинность. Но
``Ярославу заходит такой подкол'', ``Рин не любит такой тон'', ``этот человек
устал'', ``у Анны есть кот'' --- это предпочтения, реакции, склонности,
состояния. ``Уверенность 0.8, что Ярослав любит $X$'' --- не то же самое, что
``вероятность истинности факта 0.9''. Предпочтения меняются, зависят от
контекста, имеют исключения.

Поэтому у узла есть поле \texttt{node\_kind}, и оно задаёт смысл
\texttt{confidence}:

\begin{center}
\begin{tabular}{p{0.22\textwidth} p{0.28\textwidth} p{0.28\textwidth}}
\toprule
\textbf{node\_kind} & \textbf{что это} & \textbf{смысл confidence} \\
\midrule
observation & прямое наблюдение & истинность (высокая) \\
claim & утверждение о мире & истинность \\
hypothesis & вывод или гипотеза & истинность (низкая, провизорно) \\
episode & конкретное событие & было / не было \\
preference & вкус или склонность & надёжность оценки, не истинность \\
reaction\_pattern & ``на $X$ реагирую $Y$'' & надёжность, Beta-постериор \\
relationship\_state & состояние отношения / кто-что-знает & надёжность модели \\
boundary & граница или нежелание & устойчивость \\
skill/procedure & навык / как-делать & освоенность \\
self\_model & убеждение о себе & надёжность самооценки \\
mood\_summary & свёртка настроения & это метка, не убеждение \\
\bottomrule
\end{tabular}
\end{center}

Для не-truth-apt видов confidence читается как надёжность или устойчивость
оценки, а не ``вероятность истинности''. Предпочтения и реакции лучше хранить
как модель с условиями, а не как голый скаляр истинности.

\begin{principle}[Confidence priors = source\_trust $\times$ kind, а не плоская константа]
Импортированным узлам, то есть эмбриональному графу (ConceptNet, ATOMIC),
нельзя давать плоскую уверенность на всё. Baseline должен зависеть от двух
осей:
\[
\mathrm{confidence}(\mathrm{node}) =
\mathrm{source\_trust}(\mathrm{source}) \times
\mathrm{baseline}(\mathrm{node\_kind}).
\]
По \texttt{kind}: ConceptNet-концепт --- truth-apt claim, стартует выше;
ATOMIC-эпизод --- социально-каузальная оценка реакции, стартует ниже, потому
что это надёжность мягкого прогноза, а не истинность. По источнику: разные
базы заслуживают разного доверия.
\end{principle}

Два принципиальных решения здесь особенно важны.

\begin{itemize}
\item Источник --- это провенанс, а не метка. \texttt{source} уже несётся в
провенансе (\texttt{/c/en/...} vs \texttt{atomic}); не надо плодить
эпистемические метки \texttt{prior\_conceptnet} / \texttt{prior\_atomic}. Метка
\texttt{prior} --- это эпистемический класс для packer'а, а источник ---
ортогональные метаданные для калибровки.
\item Не выводить node-confidence из веса ребра. Вес ConceptNet --- это
корроборация отношения, канал ребра, а не истинность узла. Вес гонит
диффузию, confidence узла --- про доверие к содержанию.
\end{itemize}

\begin{intuition}[Почему знаковый скаляр, а не тройка мнения]
Напрашивается более богатое представление --- тройка субъективной логики
\emph{belief / disbelief / uncertainty}, где неопределённость несёт собственную
координату. Мы сознательно берём один знаковый скаляр $c \in [-1, 1]$
(эквивалент $2P - 1$: $+1$ ``да'', $-1$ ``нет'', $0$ ``не знаю''). Причины
две. Во-первых, скаляр уже разводит неверие ($c < 0$) и незнание
($c \approx 0$). Во-вторых, третью ось --- ``сколько за этим свидетельств /
насколько оценка устойчива'' --- мы и так несём отдельно и разнотипно: через
\texttt{node\_kind}, через \texttt{need\_to\_discover}, через
Beta-неопределённость реакций (гл.~\ref{ch:personality}) и провенанс.
Неопределённость у нас не потеряна, она просто не запихнута в один кортеж.
\end{intuition}

\begin{intuition}[Кто-что-знает --- тоже вид убеждения]
Особо отметим \texttt{relationship\_state} в смысле моделирования чужих
знаний. ``Рин знает $X$'', ``Ярослав знает, что Анна сказала $X$'', ``Анна не в
курсе, что Рин это запомнила'' --- это разные узлы: убеждения Рин о знаниях
других. Это не про приватность-как-политику, а про то, что theory-of-mind ---
законный и важный вид памяти: без него Рин не отличит ``я это знаю'' от ``он
думает, что я этого не знаю''.
\end{intuition}

\subsection{От словаря к поведению: операционная матрица}
\label{sec:kindmatrix}

Список видов --- это словарь; реализация спросит \emph{поведение}. Каждый
\texttt{node\_kind} обязан ответить на семь вопросов: какие операторы его
читают, какие могут в него писать, как считается confidence, как работает
распад, какие типы рёбер допустимы, можно ли его fusion-ить и можно ли
показывать LLM в context-packer'е (гл.~\ref{ch:runtime}).

Полную матрицу $11 \times 7$ заполнять сейчас --- значит выдумывать точность,
которой нет. Поэтому ниже --- каркас с наиболее решающими столбцами и
уверенными клетками; прочее помечено как TBD.

\begin{center}
\begingroup
\small
\setlength{\tabcolsep}{4pt}
\begin{tabular}{@{}p{0.18\textwidth} p{0.21\textwidth} p{0.16\textwidth} p{0.13\textwidth} p{0.16\textwidth}@{}}
\toprule
\textbf{kind} & \textbf{смысл confidence} & \textbf{распад} & \textbf{fusion?} & \textbf{в packer?} \\
\midrule
observation & истинность & медленный & да (посылка) & да \\
claim & истинность & медленный & да & да \\
hypothesis & истинность (низкая) & быстрый, если не подтверждена & да & да, с меткой ``гипотеза'' \\
preference & надёжность оценки & по использованию & осторож\-но & да \\
\shortstack[l]{reaction\_\\pattern} & Beta-надёжность & по успеху / провалу & нет, это модель & как контекст \\
\shortstack[l]{relationship\_\\state} & надёжность модели & медленный & осторож\-но & \shortstack[l]{по ситуации\\(TBD)} \\
boundary & устойчивость & очень медленный & нет & да \\
\shortstack[l]{self\_\\model} & надёжность самооценки & медленный & да (метакогниция) & да \\
\shortstack[l]{mood\_\\summary} & это метка & транзиент\-но & нет & редко \\
\bottomrule
\end{tabular}
\endgroup
\end{center}

Столбцы ``кто читает / кто пишет / допустимые рёбра'' определяются там, где
живёт соответствующий оператор, а не дублируются здесь. Матрица --- это
\emph{контракт} поведения на тип, и она сознательно неполна: заполняется по
мере постройки.

\section{Провенанс и обоснование}
\label{sec:prov}

Поле \texttt{provenance} шире, чем ``источник''. Для \emph{исходных} фактов
провенанс --- откуда факт узнан. Для узлов, \emph{порождённых выводом}
(fusion), провенанс должен хранить \emph{обоснование} --- множество
родительских узлов, из которых вывод получен. Эта деталь оказывается
критической в главе~\ref{ch:epistemics}: именно наличие обоснований позволяет,
опровергнув посылку, автоматически пометить на пересмотр её следствия. Без
записанных обоснований опровержение посылки не находит потомков --- и
рождается ``ползучая вера''.

\section{Что значит «коммит» и кто его гейтит}

Важнейшая граница схемы --- между активным рабочим убеждением и записью. Не
всё, что всплыло или сгенерировалось, попадает в граф как факт.

\begin{principle}[Гипотеза, а не факт]
Всё, что порождено выводом или генерацией (fusion, реакция, ассоциация),
пишется провизорно: с низкой уверенностью, с провенансом-обоснованием и, при
необходимости, с флагом \texttt{need\_to\_discover}. Повышение до ``записи'',
то есть переход в более медленный слой, происходит только через накопление
подтверждающих свидетельств. Управление --- на границе коммита.
\end{principle}

\begin{lstlisting}[style=pseudo]
def commit(G, item, kind):
    if kind == "observed":          # прямое наблюдение
        c0 = source_trust(item.source)
    if kind == "inferred":          # порождено fusion / выводом
        c0 = LOW                    # гипотеза, не факт
        item.provenance = item.parents      # обоснование для TMS
    node = G.add_node(item, confidence=c0, layer="session")
    propagate(G, node)              # см. главу про лавину
    return node
\end{lstlisting}

\section{Кража схемы у Wikidata}

Есть готовая, обкатанная модель данных, уже решившая ровно те задачи, что
встают перед метаданными Рин, --- модель утверждений \emph{Wikidata}. Брать
её \emph{содержание} для существа, мимикрирующего под человека, не нужно:
человек помнит мало фактов, его сила в ассоциативности. Но её \emph{схему}
стоит перенять целиком:

\begin{itemize}
\item квалификаторы (точка во времени, начало, конец) --- темпоральная
валидность: факт истинен в интервале, а не вечно;
\item references --- провенанс и обоснование;
\item ранг (\emph{preferred} / \emph{normal} / \emph{deprecated}) ---
недеструктивная ревизия: устаревшее не удаляют, а депрецируют, помечая новое
предпочтительным; история сохраняется для машинерии поддержки истинности.
\end{itemize}

Иными словами, факт стоит моделировать как
\[
\text{(утверждение)} + \text{(интервал валидности)} +
\text{(провенанс)} + \text{(статус / ранг)},
\]
и это буквально схема Wikidata. Не изобретать метаданные заново, а взять
проверенный шаблон.

\section{Врождённое и приобретённое: не граница, а континуум}

Соблазнительно разделить узлы на ``врождённые константы'' (эмбриональный
граф) и ``приобретённые'' (нажитые, с временными метками). Как первое
приближение полезно, но как жёсткая граница неверно. Консолидация
\emph{двигает} записи через границу (эпизод $\to$ обобщённое знание), а
ревизия должна уметь \emph{понижать} и ``врождённое'' --- иначе унаследованную
ошибку здравого смысла Рин не исправит никогда. Правильная картина ---
\emph{континуум глубины консолидации}, где порог опрокидывания убеждения
растёт с глубиной, но нигде не становится бесконечным. К оси
``пластичность -- упрямство'' мы ещё не раз вернёмся
(главы~\ref{ch:propagation}, \ref{ch:epistemics}).

\chapter{Обход графа: один движок, три политики}
\label{ch:walker}

\section{Случайное блуждание как универсальный механизм}

Почти всё, что Рин делает со своей памятью, --- это обход графа роем случайных
блуждающих агентов, распространяющих активацию. Retrieval, внутренняя жизнь
и сон --- не три разных движка, а три политики одного движка, различающиеся
смещением перехода, вероятностью телепорта и набором операций.

Общее правило перехода --- нормированное произведение факторов:
\begin{equation}
P(A \to B) \propto
w_{AB}
\cdot \mathrm{rel}(B \mid q)
\cdot \mathrm{mc}(B \mid \mood)
\cdot \mathrm{unc}(B)^{\beta_u},
\label{eq:trans}
\end{equation}
где $\mathrm{mc}$ --- из \eqref{eq:moodcong}, а показатель $\beta_u$
включает/выключает любопытство.

\begin{intuition}[Что такое «температура»]
«Температура» --- одна ручка: насколько выбор следующего шага случаен против
жадного. Формально переход берётся softmax-ом от логарифма правой части
\eqref{eq:trans}, делённого на температуру $T$: при $T \to 0$ --- жадный шаг по
сильнейшему ребру (эксплуатация), при большом $T$ --- почти равновероятный
выбор соседа (исследование). Термин из имитации отжига; мистики нет.
\end{intuition}

\section{Три политики как параметризация}

\begin{center}
\begin{tabular}{>{\raggedright\arraybackslash}p{0.30\textwidth} c c c >{\raggedright\arraybackslash}p{0.25\textwidth}}
\toprule
Политика & $\alpha$ (телепорт) & $\beta_u$ & $T$ & Что запускает \\
\midrule
retrieval (цель) & $0$ & $0$ & низкая & извлечение под запрос \\
внутренняя жизнь (любопытство) & мал. & $>0$ & средняя & inspect, fusion у неопределённых \\
сон (покрытие) & большой & $0$ & высокая & fusion/inspect у случайных \\
\bottomrule
\end{tabular}
\end{center}

\smallskip\noindent\textbf{Retrieval --- целенаправленная.}
Стартует от заякоренных запросом узлов, собирает окрестность активированного
содержания. Подчеркнём (глава~\ref{ch:personality}): обход собирает ингредиенты,
а не готовую реакцию.

\smallskip\noindent\textbf{Внутренняя жизнь --- по неопределённости.}
Смещена к узлам с высоким \texttt{need\_to\_discover}, низкой уверенностью,
новизной ($\beta_u > 0$). «Любопытство» --- целенаправленный поиск
неопределённости, не случайность.

\smallskip\noindent\textbf{Сон --- телепортирующий обходчик.}
Политика стохастического покрытия: с вероятностью $\alpha$ телепорт в
случайный узел, иначе --- шаг по \eqref{eq:trans}. Это случайное блуждание с
рестартом (персонализированный PageRank). Его стационарное распределение
удовлетворяет
\begin{equation}
\pi = (1 - \alpha) P^{\top}\pi + \alpha r,
\label{eq:ppr}
\end{equation}
где $r$ --- распределение рестарта. Телепорт делает блуждание эргодическим ---
это ключевое свойство для эпистемической гигиены (глава~\ref{ch:topology}).

\begin{intuition}[Три способа гулять по городу]
Retrieval --- идти по адресу (есть цель, выбираешь осмысленный путь).
Внутренняя жизнь --- идти туда, где интересно и непонятно (тянет к незнакомым
переулкам). Сон --- время от времени садиться в случайное такси и ехать в
случайный район: не потому, что там дело, а чтобы не застревать в одном
квартале. Один пешеход, три правила выбора улицы.
\end{intuition}

\section{Что walker делает с графом: читает, предлагает, но не коммитит}
\label{sec:proposecommit}

Естественный вопрос: обходчик сам правит граф --- сливает узлы, меняет
уверенности, ретрактит? Ответ, определяющий всю конкуренцию в системе: нет.
Walker --- это сенсор-предлагатель, а не писатель. Мы разводим три роли,
которые легко слить в одну и потом мучиться с гонками:
\begin{itemize}
\item чтение (обход) --- walker ходит по графу и собирает активированное
содержание; это то, чем занимаются все три политики;
\item предложение --- итог обхода не «правка», а типизированное предложение
операции: не TODO-текст, а структура вида
\texttt{\{op: consolidate, members:[A, B, C, D], rationale, support\_provenance, walker\_id, t\}},
которая кладётся в буфер предложений;
\item коммит --- отдельная, сериализованная стадия (тот же
schema-writer / оператор консолидации, гл.~\ref{ch:fusion}) вычитывает буфер и
применяет или отклоняет операции, дописывая событие в лог
(гл.~\ref{ch:runtime}, раздел~\ref{sec:eventlog}).
\end{itemize}

Это ровно наш принцип «выход fusion --- гипотеза, а не факт»,
распространённый на все структурные операции: обход рождает кандидата,
коммиттер его судит. Зачем такая дисциплина, а не «walker сам чинит граф на
месте»: во-первых, конкурентная безопасность --- структурный писатель ровно
один, значит двух одновременных структурных правок не бывает по построению
(гонки, гл.~\ref{ch:runtime}); во-вторых, аудируемость --- каждая правка проходит
через event-log и объяснима реплеем; в-третьих, эпистемика --- предложение это
гипотеза, а коммит это отдельный гейтируемый шаг, где работает независимость
свидетельств и карантин (гл.~\ref{ch:epistemics}). Walker предлагает ---
коммиттер располагает.

\begin{intuition}[Единственное, что walker пишет сам --- но это тоже запись]
Полный запрет на запись был бы неудобен: обход должен оставлять след «здесь
были, это подсветили». Поэтому walker напрямую трогает только поле активации /
яркости $a_v$ (и связанную телеметрию внимания). Точная формулировка важна во
избежание тихих багов: это не «не запись», а неструктурная коммутативная
запись --- не узлы, не рёбра, не уверенности, а состояние внимания, которому
всё равно нужны TTL, затухание и \texttt{access\_reason} (гл.~\ref{ch:timescales}).
«Коммутативная» значит, что итог не зависит от порядка параллельных записей.
Но конкретно для яркости не годится аддитивный счётчик: на хабах, через которые
проходят сотни walker'ов, $a_v \mathrel{+}= 1$ улетает в космос. Берём max-CRDT:
$a_v \leftarrow \max(a_v, e_{rw})$, где $e_{rw}$ --- энергия зашедшего обходчика.
Тогда яркость ограничена сверху энергией агента, спад предсказуем, а сотня
одновременных подсветок даёт корректный итог без блокировок. Всё, что
структурно, --- только через предложение и коммит.
\end{intuition}

\section{Фоновая разведка: один процесс, а не «день/ночь»}
\label{sec:background}

Разделение на «дневное любопытство» и «ночной обходчик» --- удобная, но
неточная рамка. Точнее: есть один всегда включённый фоновый разведчик, а
«день» и «ночь» --- два конца континуума по свободным ресурсам, а не
переключатель.

Формализуем «мощности, внимание, фокус». Пусть $\rho(t) \in [0, 1]$ ---
foreground-нагрузка: насколько текущий момент требует активного ответа
($\rho \approx 1$ в напряжённом диалоге, $\rho \approx 0$ в простое). Общий
бюджет $B_{\text{total}}$ делится с приоритетом переднего плана; фоновой
разведке достаётся
\[
B_{\text{bg}} = B_{\text{total}} (1 - \rho).
\]

Фокус --- это активное множество $A_{\text{focus}}$ (узлы высокой активации
$a_v$). Параметры фонового роя --- функции $B_{\text{bg}}$: число walker'ов
$\propto B_{\text{bg}}$; доля телепортов $\alpha$ и «широта» растут при падении
$\rho$; стартовое распределение интерполирует
\[
r = (1 - \beta)\cdot \mathrm{anchor}(A_{\text{focus}}) + \beta\cdot \mathrm{uniform},
\qquad \beta \uparrow \text{ при } \rho \downarrow.
\]

Тогда прежние «любопытство» и «сон» --- это два режима одного процесса:
высокий $\rho$ (мало бюджета) $\Rightarrow$ мало walker'ов, $\alpha \approx 0$,
старт от фокуса --- локальная разведка вокруг текущей темы; низкий $\rho$
(много бюджета) $\Rightarrow$ много walker'ов, большой $\alpha$, старт почти
равномерный --- широкий телепорт-скраб (глава~\ref{ch:sleep}).

\begin{intuition}[Мысль на перекуре]
Именно поэтому яркая ассоциация может прийти днём, «на ровном месте»:
фоновый разведчик работает всегда, и даже на скудном дневном бюджете иногда
набредает на далёкую связь. Развёрнутый пример такой дневной находки ---
в главе~\ref{ch:fusion} (House / Cast Away).
\end{intuition}

\subsection{Поля активации: быстрое внимание и медленный след --- не путать}
\label{sec:brightness}

Здесь легко совершить ошибку --- свалить в одно поле «яркость», «свежесть»,
«долговременную силу» и «важность». Соблазнительно, но потом это поле начнёт
делать слишком много, и, что хуже, диагностика начнёт тренировать память:
скраб, дебаг и тесты будут поднимать ту же переменную, что и настоящее
вспоминание, и случайно укреплять узлы. Поэтому разведём разные по темпу
состояния:
\begin{itemize}
\item $a_{\text{work}}(v,t) \in [0,1]$ --- рабочая яркость внимания. Живёт
секунды-минуты, формирует зону фокуса $A_{\text{focus}}$ и вход дневных
walker'ов, быстро гаснет. Это неструктурная коммутативная запись обходчика
(max-CRDT выше).
\item $B_{\text{base}}(v)$ --- медленная сила следа (это и есть ACT-R base-level,
гл.~\ref{ch:timescales}). Живёт дни-месяцы, обновляется только настоящими
значимыми обращениями (по \texttt{access\_reason}), влияет на decay / забывание.
\item $\mathrm{importance}(v)$ --- значимость от appraisal / домена / структуры
(гл.~\ref{ch:personality}), не путать с частотой обращений.
\item $\mathrm{coverage\_deficit}(v)$ --- «давно не инспектировали», рычаг
ночного скраба.
\end{itemize}

Ключевое разделение: настоящее вспоминание поднимает и $a_{\text{work}}$
(мгновенно), и $B_{\text{base}}$ (чуть-чуть, медленно); а скраб / дебаг / тест
трогают в лучшем случае $\mathrm{coverage\_deficit}$ и никогда не укрепляют
след. Правило быстрой яркости:
\[
a_{\text{work}}(v,t+\Delta t) = a_{\text{work}}(v,t)e^{-\Delta t/\tau(v)}
\quad \text{(между визитами),}
\]
\[
a_{\text{work}}(v) \leftarrow \max\!\bigl(a_{\text{work}}(v), e_{rw}\bigr)
\quad \text{(genuine-визит),}
\]
где $\tau(v)$ не универсальна (важные / консолидированные остывают медленнее),
а $e_{rw}$ --- энергия обходчика (max, не сумма, --- иначе хабы взрываются).

\begin{openq}[max теряет «навязчивость» --- возможно, нужен второй канал]
Честный компромисс max: он ограничивает яркость (хорошо против хабов), но
теряет повторный фокус в коротком окне. Сто обходчиков через хаб --- max
правильно не взрывает; но если один узел многократно всплывает в текущем
рассуждении, иногда важно знать не только пик, но и интегральную экспозицию
(«это навязчиво лезет прямо сейчас»). Возможное решение --- два канала:
$a^{\text{peak}}_{\text{work}} = \max(\ldots)$ для яркости и отдельный
\texttt{visit\_mass\_window} для краткосрочной статистики (он не укрепляет
$B_{\text{base}}$ автоматически). На Stage 1 не обязателен; помечаем как
возможное расширение.
\end{openq}

Теперь смещение входа обходчиков. Распределение рестарта $r$ из
\eqref{eq:ppr} берётся не равномерным, и день с ночью различаются формой
смещения, а не разными движками:
\begin{itemize}
\item День (высокий $\rho$, радиус $r_0$): $r \propto a_{\text{work}} \times
\mathrm{relevance} \times \mathrm{uncertainty}$. Обходчик эксплуатирует текущий
фокус --- близкие ассоциации «на сейчас», далеко не отходит.
\item Ночь (низкий $\rho$, радиус $r_1 > r_0$): $r \propto \mathrm{uniform} +
\mathrm{coverage\_deficit} + \mathrm{need\_to\_revise} + \text{low-}\pi$ регионы.
Именно так, а не просто «шире от дневного фокуса» --- иначе сон останется
заложником дневного внимания. Сон нужен, чтобы уйти туда, куда внимание
никогда не ходит: anti-entropy scrub, телепорт-эргодичность, проверка
противоречий (гл.~\ref{ch:sleep}).
\end{itemize}

Лаконично: день --- эксплойт локально у горячего, ночь --- исследуй заброшенное.
Поля разные, но связанные: $a_{\text{work}}$ ведёт дневной фокус,
$\mathrm{coverage\_deficit}$ и низкая $\pi$ ведут ночь, $B_{\text{base}}$ решает,
что переживёт забывание.

\begin{intuition}[Мотыльки на остывающий свет: откуда дневные догадки]
Если дневной ретрив --- «слепой матричный ток» (раздел~\ref{sec:semgate}),
откуда берутся внезапные дневные озарения? Не из matvec. После дневного ответа
подсвеченные узлы $a_{\text{work}}$ гаснут экспоненциально минутами-часами ---
остаётся тёплое послесвечение. Одиночные фоновые RW (тот самый остаточный
бюджет $B_{\text{bg}}$) через формулу рестарта
$r \propto a_{\text{work}} \times \mathrm{relevance} \times \mathrm{uncertainty}$
получают повышенную вероятность телепортироваться именно в эти недавно-горячие
зоны --- летят на остывающий свет, как мотыльки. Кружа в тёплой области, они
сталкивают концепты на своих траекториях и роняют кандидатов в thought-inbox
(гл.~\ref{ch:runtime}). Важно: роняют именно кандидатов (в основном шум),
которых обязан просеять salience-гейт ($|\Delta \mood|$, ошибка предсказания) ---
иначе Рин завалит себя мусорным fusion. Так разделение труда честное: matvec
подсвечивает, мотыльки ползают по послесвечению, salience решает, что вообще
достойно всплыть догадкой.
\end{intuition}

\begin{intuition}[Где «живут» обходчики]
Walker'ы не объекты, которые где-то постоянно «живут». Это фоновый процесс
планировщика с бюджетом $B_{\text{bg}}$: каждый обход рождается из
распределения рестарта $r$, проходит, оставляет предложения и подсветку --- и
гаснет. Persist между ними не сами walker'ы, а поле яркости $a_{\text{work}}$ и
буфер предложений. Спросить «где живут обходчики» --- всё равно что спросить,
где живут волны: волн как вещей нет, есть среда (поле $a_{\text{work}}$) и то,
что её возмущает (планировщик).
\end{intuition}

\begin{openq}[Как измерять foreground-нагрузку $\rho$]
$\rho$ управляет тем, насколько агрессивно разведчик забирает ресурсы, но как
её мерить в рантайме --- не зафиксировано. Кандидаты: наличие активного
пользовательского хода (есть незакрытая реплика $\Rightarrow \rho$ высока);
длина очереди foreground-задач; «важность» текущего события (по appraisal,
глава~\ref{ch:personality}); загрузка вычислителя. Вероятно, комбинация, но
веса и динамика (как быстро $\rho$ спадает в простое) пока открыты.
\end{openq}

\section{Точки входа: как query превращается в якоря}
\label{sec:entrypoints}

Ретрив у нас двухшаговый: шаг 1 --- найти точки входа (якоря) в граф, шаг 2 ---
растечься из них диффузией (раздел~\ref{sec:semgate}). Про шаг 2 сказано много,
а про шаг 1 --- только «быстрым индексом». Распишем его, потому что это первая
строчка кода ретрива.

\smallskip\noindent\textbf{Сразу уточнение терминологии.}
Query не «разбивают на эмбеддинги». Его разбивают на семантические единицы,
и уже каждую единицу отдельно эмбеддят. Это разные вещи.

\smallskip\noindent\textbf{Кто разбивает --- не LLM.}
Ретрив --- рефлекс (System 1, ниже LLM-слота, гл.~\ref{ch:runtime}), поэтому
превращение query в якоря должно быть дешёвым и LLM-free на горячем пути.
Делает это тот же экстрактор, что и на write-path
(\texttt{raw\_event -> parse/extract candidates}, раздел~\ref{sec:writepath}),
только без коммита. Красивая симметрия: понимание запроса = извлечение
кандидатов из события минус запись. На Stage 1 экстрактор простой: NER +
ключевые фразы + лемматизация (для русского морфология обязательна ---
«плавки / плавок / плавками» сводятся к лемме). LLM подключается к
переформулировке мутного запроса только на эскалации --- редкий System 2, а не
норма.

\smallskip\noindent\textbf{Два сорта якорей.}
Экстрактор выдаёт две разновидности, и ищутся они по-разному:
\begin{itemize}
\item Сущностные / лексические --- именованные сущности и конкретные термины
(«Азоры», «плавки», «9 сентября»). Им нужен точный матч, а не «похожее по
вектору».
\item Холистический --- эмбеддинг всего запроса (или его клауз) для
ассоциативной диффузии.
\end{itemize}

Если единиц несколько --- это мульти-векторный ретрив, а не одно усреднённое.
Mean-pool «плавки на Азоры 9 сентября» размазал бы три резкие сущности в кашу;
поэтому каждую единицу сажаем отдельным seed'ом, а узлы, подсвеченные
несколькими seed'ами, получают буст (пересечение = релевантнее).

\smallskip\noindent\textbf{Чем эмбеддим.}
Тем же энкодером, что и \texttt{node.embedding} при записи (гл.~\ref{ch:substrate}).
Это не опция: косинус осмыслен, только если query и узлы в одном пространстве.
Вектор узла уже кэширован (считался один раз коммиттером), так что на шаге
ретрива через энкодер гоним только query-единицы, узлы не трогаем.

\begin{flushleft}\ttfamily
1 units      = extract(query + active\_focus)      \# NER + keyphrases + леммы (+ кореф)\\
2 anchors    = \{\}\\
3 for u in units:\\
4 \quad if u.is\_named\_entity:\\
5 \qquad anchors |= exact\_index.lookup(lemma(u)) \# точный/alias-индекс (forensic-подобный)\\
6 \quad anchors |= ann\_index.knn(embed(u), k)     \# вектор-ANN (HNSW), top-k по косинусу\\
7 anchors    = resolve\_entities(anchors)          \# упоминания -> канонические узлы\\
8 r          = weight\_by\_match\_confidence(anchors) \# точное > сильный вектор > слабый\\
9 \# r -> распределение рестарта для локальной разреженной диффузии (шаг 2)
\end{flushleft}

То есть точки входа приходят из двух индексов: точного / alias (для именованных
сущностей --- «Азоры» должны попасть в тот самый узел) и вектор-ANN (для
семантической близости). Их объединение, взвешенное по уверенности матча, и
есть seed-набор. Плюс entity resolution: разные упоминания сводятся к
каноническому узлу, чтобы «плавки» били в один узел, а не в пять эпизодов
покупки.

\smallskip\noindent\textbf{Дальше.}
Seed-набор становится начальным впрыском $a(0)=s(0)$ для локальной затухающей
диффузии (fast-path вокруг якорей, раздел~\ref{sec:semgate}), а расширение
гейтируется условным throttle. Важно: retrieval --- это политика с телепортом
$\alpha = 0$ (таблица политик выше), то есть источник впрыскивается один раз и
импульс затухает; рестарт (член $s(t>0)$, он же $\alpha>0$ в \eqref{eq:ppr}) ---
это сон, а не retrieval. Цепочка целиком: экстрактор дал якоря $\to$
ANN+точный индекс нашли точки входа $\to$ импульс затух $\to$ packer срезал по
порогу.

\begin{openq}[Что на Stage 1 сознательно упрощаем]
Открыты: гранулярность разбиения (по сущностям / клаузам / предложениям ---
тюнить на бенче); кореференс --- «поищи их» требует обогатить query текущим
фокусом $a_{\text{work}}$ до эмбеддинга, иначе якорей нет; порог эскалации к
LLM-переформулировке; качество alias-резолвинга под русскую морфологию.
Механика (два сорта якорей, два индекса, мульти-seed, resolution) ---
зафиксирована.
\end{openq}

\section{Семантический вентиль: дёшево и без протечек по омонимии}
\label{sec:semgate}

В \eqref{eq:trans} множитель $\mathrm{rel}(B \mid q)$ --- это семантический
вентиль: он не даёт фронту расползтись из зоны запроса. Возникает законный
страх за производительность: если на каждом шаге RW считать сходство контента
узла с $q$ заново, это катастрофа. Ответ: на лету ничего не эмбеддим.
Эмбеддинг узла считается один раз при записи (коммиттером) и лежит в поле
\texttt{embedding} (гл.~\ref{ch:substrate}). Тогда вентиль на шаге --- это
скалярное произведение фиксированного вектора $q$ на кэшированный вектор узла,
$O(d)$; а раз ретрив у нас --- усечённая локальная затухающая диффузия из
seed'ов (политика retrieval: телепорт $\alpha = 0$, источник впрыскивается один
раз --- это не PPR; рестарт-форма \eqref{eq:ppr} зарезервирована за сном), весь
фронт считается одним matvec, где вектор активации на каждом шаге умножается на
маску сходства с $q$. Как только значения падают ниже порога $\varepsilon$ ---
граница замкнулась. Никаких рекурсивно плодящихся агентов: «группа вспышек,
слабеющих к окраинам» --- это результат диффузии, а не объектная механика.

Инженерная честность про этот дневной fast-path: matvec идёт не по всему
графу. Умножать $\tilde{P}^{\top}$ на $|V|$-мерный вектор на графе в сотни
тысяч узлов --- лечь. Это локальная разреженная диффузия на ограниченном
носителе: 2--3 хопа вокруг seed-якорей, найденных ANN-индексом (HNSW), ---
субматрица смежности окрестности, а не всё $V$-пространство. И тайминги честные
--- миллисекунды, не микросекунды (основное время --- ANN-поиск якорей и
context-packer), а настоящее бутылочное горло рантайма всё равно LLM-инференс
(секунды), гл.~\ref{ch:runtime}.

Две поправки, чтобы вентиль не навредил. Первая --- жёсткость по типу
ретрива: для ordinary-ретрива режем жёстко (защита от мусора), для
creative / curiosity --- мягко (именно далёкие странные мосты нужны для fusion),
для forensic --- вентиля нет вообще, это не walker, а точный индекс-lookup
(таксономия ретрива, гл.~\ref{ch:runtime}). Слишком жёсткий вентиль убил бы
серендипность. Вторая --- полисемия: граница возникает не только из низкого
сходства с $q$, но и из sense/type-несовпадения. Классический пример
«кружка $\to$ ручка $\to$ дверь» не должен решаться тем, что «дверь слабо
похожа на кружку»: раньше ломается слой смысла. «Ручка кружки», «дверная
ручка», «ручка-перо» --- это разные sense/concept-узлы (или хотя бы
квалифицированные рёбра, гл.~\ref{ch:embryonic}). Без этого ретрив постоянно
течёт по омонимии --- а для русского это особенно больно. Значит вентиль
закрывается по $\max(\text{low-sim}, \text{sense/type-mismatch})$, а не по одной
similarity.

\subsection{Вентиль --- условный throttle, а не плоский множитель}
\label{sec:semgate-conditional}

Есть и третья поправка, и она важнейшая: как написано выше,
$\mathrm{rel}(B \mid q)$ входит в \eqref{eq:trans} множителем --- а значит
зануляет путь при низком сходстве. Это тихо убивает личность. Пример: триггер
«гирлянда». Узлы [свет], [провода] семантически близки и проходят; а узел
личной памяти [окно в Тбилиси] имеет сходство $\approx 0$ --- и множитель его
гасит, Рин выдаёт энциклопедию. В живой голове [Тбилиси] вспыхивает первым:
у личного ребра колоссальный вес / валентность, несмотря на семантическую
даль. Значит вентиль должен быть условным throttle, а не вето:
\[
\mathrm{Score}(A \to B) =
f(w_{AB}, \mathrm{valence}_B)\cdot
\mathrm{Throttle}(B,q)\cdot
\mathrm{mc}(B \mid \mood)\cdot
\mathrm{unc}(B)^{\beta_u},
\]
где
\[
\mathrm{Throttle}(B,q) =
\begin{cases}
1, & \shortstack[l]{если ребро из личного слоя (session/life)$;$\\
или $|\mathrm{valence}_B| > \theta_{\text{emotion}}$,}\\
\mathrm{rel}(B \mid q), & \text{иначе (обычное знание --- тормоз против дрейфа).}
\end{cases}
\]

То есть сильное личное или аффективно заряженное ребро --- «синаптический
хайвей» --- проходится независимо от сходства с запросом; а throttle тормозит
только слабые / родовые рёбра, чтобы не было случайного топологического
дрейфа.

Три предохранителя, без которых это сломается. Не по сырому $w_{AB}$:
пробивать вентиль порогом веса нельзя --- у эмбриональных рёбер (ConceptNet)
веса из коробки огромны ($w \approx 2{,}5$), и Рин уплыла бы по банальной
таксономии. Поэтому гейт по происхождению ребра (личный слой, не врождённый
каркас), а не по величине веса. $f(w_{AB}, \mathrm{valence})$ всё равно
ранжирует: $\mathrm{Throttle}=1$ не значит «все личные узлы хлынут» ---
Тбилиси всплывает не потому, что ребро личное, а потому, что оно личное и с
огромным $w_{AB}$; слабое личное ребро с $\mathrm{Throttle}=1$ всё равно
проиграет по $f$. Стоп-условие три-гейта: путь глохнет, только если
одновременно низкое сходство и слабое ребро и нейтральная валентность; сильна
любая из трёх --- агент проходит. Сухой-но-важный boundary («не шути про
доту») проходит через layer-гейт даже при низкой валентности. Омонимию
(кружка $\to$ ручка $\to$ дверь) по-прежнему режет отдельный sense/type-гейт,
а не throttle.

\begin{intuition}[Прайминг: подсвечено, но не сказано]
База при этом не пропадает: «у гирлянды есть провода» матричная диффузия
подсвечивает (ток туда течёт) --- это прайминг. Но у этих узлов почти ноль
ситуативной значимости и новизны, поэтому context-packer их сознательно не
выносит наверх (раздел~\ref{sec:packer}). Они «тёплые» на фоне сознания, но не
доходят до языка --- ровно как человеческое подсознание держит очевидное
наготове, не проговаривая.
\end{intuition}

\begin{principle}[Вентиль на эмбеддингах --- фильтр родственности, не релевантности]
Эмпирически выяснено (стенд, \texttt{reports/experiment\_3--4}): множитель
$\mathrm{rel}(B \mid q)$, построенный на дистрибутивных эмбеддингах (косинус к
вектору запроса), фильтрует родственность (co-occurrence), а не причинную
релевантность. Он давит off-topic дрейф (годится для запросов вида «что есть
$X$» / «кто из этих $X$»), но на «почему»-запросах систематически
недобирает: узел может быть ближе к слову запроса как ассоциат, не будучи
ответом («beverage» ближе к «energy», чем «caffeine», из-за маркетингового
со-упоминания). Это ограничение инструмента, а не бага; для причинного выбора
нужен не косинус, а рассуждение.
\end{principle}

\begin{principle}[Конвергенция на хабе пересиливает per-edge релевантность; лечится редкостью, не гейтом]
Эмпирически выяснено (там же): при диффузии узел, достижимый из нескольких
якорей сразу, накапливает их вклады, и это суммирование может перевесить более
высокий per-edge $\mathrm{rel}(B \mid q)$ у односвязного правильного ответа ---
даже когда вентиль оценил ответ выше. То есть провал «генеричный дистрактор
победил» решает конвергенция на хабе, а не релевантность; это тот же зверь,
что концентрация $\pi$ на хабах (гл.~\ref{ch:topology}). Значит корректив ---
отдельный штраф за редкость / степень (IDF-подобный: генеричный общий сосед
давится), а не докрутка вентиля. Один такой член чинит сразу три наблюдения:
дистрактор-хаб, таксономический дрейф ретрива и хаб-концентрацию.
\end{principle}

\begin{principle}[Ретрив отвечает за recall, LLM --- за причинный выбор]
Отсюда критерий успеха ретрива: не «правильный узел первый в ранге», а
«правильный узел попал в упакованный контекст в рамках бюджета при приемлемой
доле дистракторов». Граф --- ассоциативная память (поднять кандидата в
окрестность); причинное ранжирование («из этих caffeine --- и есть почему»)
делает LLM на собранном контексте. Требовать от графа самому назвать причинный
ответ --- значит возлагать рассуждение на recall.
\end{principle}

\section{Единство и следствие}

Один субстрат, один блуждающий агент, три параметризации: смещение перехода,
вероятность телепорта, набор операций. Реализуется одна машина обхода с
параметрами (псевдокод --- приложение~\ref{ch:algorithms}). Родственные
механизмы: spreading activation (Collins--Loftus) как модель ассоциативного
припоминания и персонализированный PageRank / RWR как его вычислительная
форма; в волте RIN --- под именами SYNAPSE и GAAMA.

\chapter{Настроение}
\label{ch:mood}

\section{Настроение --- отдельный сцепленный модуль, а не поле графа}

Снимем путаницу сразу. Настроение не вычитывается из спектра графа памяти.
Это отдельное состояние --- низкоразмерный вектор в пространстве аффекта
$\mood \in \mathbb{R}^k$ (например, оси энергичность / раздражённость /
экзальтация / любопытство) --- отдельный модуль, слой, через который проходит
извлечённое содержание.

Путаница шла от слова «спектр». Есть два разных спектра, и их нельзя
смешивать: спектр лапласиана графа (структура памяти, глава~\ref{ch:topology})
и частотное разложение настроения во времени (быстрые и медленные компоненты
аффективного сигнала). Многослойность настроения относится ко второму и к
графу отношения не имеет.

\section{Инерция и многослойность}

У настроения должна быть инерция: оно зависит и от события, и от собственного
прошлого значения:
\begin{equation}
\mood_t = f(\mood_{t-1}, event_t), \qquad action_t = g(\mood_t, event_t).
\label{eq:moodinertia}
\end{equation}

Это протекающий интегратор \eqref{eq:leaky}: инерция --- это его
рекуррентность (состояние зависит от прошлого состояния).

\begin{openq}[Почему интегратор, а не выученная динамика]
Форму $f$ можно было бы выучить --- небольшой RNN/GRU поверх потока событий.
Мы берём фиксированный протекающий интегратор по трём причинам.
Интерпретируемость: у него ровно одна ручка на слой --- постоянная времени
$\tau_\ell$ (как быстро забывает), и её смысл читается сразу; у выученной
динамики «почему настроение поехало» --- это скрытое состояние, которое не
отладишь глазами. Отсутствие обучения и данных: RNN надо на чём-то
тренировать, а размеченных траекторий настроения Рин у нас нет и не
предвидится. Достаточность: главные наблюдаемые свойства --- инерция, возврат
к базовому уровню и разные скорости на разных горизонтах --- интегратор даёт
напрямую, одним стеком с разными $\tau_\ell$. Выученную динамику держим как
апгрейд на потом: если стенд покажет, что нужны нелинейные эффекты
(насыщение, гистерезис аффекта), которые линейный интегратор не тянет.
\end{openq}

\begin{intuition}[Интегратор и вектор --- одно и то же, с двух сторон]
Чтобы снять путаницу: «протекающий интегратор» и «вектор настроения» --- не
разные сущности. Интегратор --- это правило обновления \eqref{eq:leaky}, а его
состояние --- вектор в пространстве аффекта. Многослойность настроения --- это
стек таких интеграторов с разными $\tau_\ell$; «старший вектор» и «младший
вектор», о которых мы говорили, --- это ровно состояния медленного и быстрого
интеграторов. Актуальное настроение --- их суперпозиция. И, как отмечено в
разделе~\ref{sec:twolayers}, эти слои параллельны (все питаются одним потоком
событий, переноса между ними нет) --- в отличие от слоёв памяти, где перенос
есть.
\end{intuition}

Формально: стек интеграторов с разными $\tau_\ell$, актуальное настроение ---
их взвешенная суперпозиция:
\[
\mood_t^{(\ell)} = \left(1 - \frac{1}{\tau_\ell}\right)\mood_{t-1}^{(\ell)} +
\frac{1}{\tau_\ell}\,\mathrm{appraise}(event_t),
\qquad
\mood_t = \sum_\ell w_\ell \mood_t^{(\ell)}.
\]

Младший слой (малое $\tau$) даёт высокочастотную компоненту большой амплитуды
--- «дрожит сильнее». «Суперпозиция» здесь --- про частоты во времени, а не про
граф.

\section{Setpoint: медленный слой как аттрактор}

Медленный слой --- не просто демпфер, а точка притяжения. Добавляя к
\eqref{eq:leaky} возврат к базовому уровню $\mu$, получаем затухающую к
$\mu$ динамику:
\[
\mood_t^{(\mathrm{life})} =
\mu + (1 - \kappa)\bigl(\mood_{t-1}^{(\mathrm{life})} - \mu\bigr) + \ldots,
\qquad 0 < \kappa < 1.
\]

Здоровая динамика --- отрицательная обратная связь, тянущая настроение домой.
Патология --- когда сдвинулся сам $\mu$: депрессия в этой модели не длинная
экскурсия быстрого слоя, а уехавший аттрактор.

\section{Настроение зависит от графа, и граф --- от настроения}

Связь двусторонняя, и именно она делает настроение частью мышления.
\begin{itemize}
\item Граф $\to$ настроение. Извлечённые узлы несут валентность; она и есть
$event_t$ в \eqref{eq:moodinertia}.
\item Настроение $\to$ граф. Настроение смещает вероятности переходов при
обходе: узлы с валентностью, близкой к текущему настроению, получают больше
активации:
\begin{equation}
\mathrm{mood\_congruence}(B \mid \mood) =
\exp\!\bigl(-\gamma \lVert \mathrm{valence}(B) - \mood \rVert^2\bigr).
\label{eq:moodcong}
\end{equation}
\end{itemize}

Это конгруэнтное настроению припоминание: грустный чаще вспоминает грустное.
Настроение остаётся собственным состоянием, сцепленным сбоку, а не полем
графа.

\section{Опасность петли и три тормоза}

\begin{pitfall}[Руминация]
Конгруэнтное припоминание \eqref{eq:moodcong} плюс инерция
\eqref{eq:moodinertia} дают петлю с положительной обратной связью: грустная
вспоминает грустное $\to$ грустнеет $\to$ вспоминает ещё более грустное.
Устойчивость петли определяется её петлевым усилением $g$: при $g \ge 1$
спираль расходится. Медленный слой уменьшает эффективное $g$ --- вот его
стабилизирующая роль.
\end{pitfall}

Что удерживает живой ум? Три независимых тормоза; патология --- когда откажут
все три сразу:
\begin{enumerate}
\item Setpoint/гомеостаз. Медленный слой-аттрактор снижает петлевое усиление
ниже единицы.
\item Метакогнитивное разотождествление. Акт «почему я так часто злюсь?»
переводит из «быть-в-настроении» в «наблюдать-настроение» --- второй порядок,
и переход рвёт петлю первого порядка.
\item Экзогенное прерывание. Мир вламывается ортогональным событием;
внимание конечно, внешний вход вытесняет петлю.
\end{enumerate}

Инженерное требование: (i) член setpoint; (ii) монитор «валентность залипла
слишком долго» --- высокая автокорреляция $\mood_t$ --- с правом запустить
корректив (переопределение, переключение внимания, вызов сна); (iii)
всегда-прерываемая экзогенная значимость. Условие безопасности: три тормоза не
могут быть отключены одновременно. Псевдокод монитора ---
в приложении~\ref{ch:algorithms}.

\begin{intuition}[Почему многослойность --- это стабилизатор, а не украшение]
Пусть быстрый слой попал в негативную спираль с усилением
$g_{\text{fast}} = 1{,}1$ (расходится). Актуальное настроение --- суперпозиция, и
медленные слои вносят устойчивую компоненту с $g \approx 0$. Если их вес
$w_{\text{slow}}$ достаточен, эффективное усиление суперпозиции
$\bar{g} = w_{\text{fast}}g_{\text{fast}} + w_{\text{slow}}\cdot 0 < 1$ ---
спираль гасится за счёт «тяжёлого маховика» медленных слоёв. Так
многомасштабность оправдана функционально: это демпфер руминации, а не только
реализм.
\end{intuition}

\section{Метакогниция}
\label{sec:metacognition}

Один из трёх тормозов руминации выше --- «метакогнитивное
разотождествление» --- был назван мимоходом, но это несправедливо:
метакогниция --- самостоятельная фича всего проекта, а не побочный эффект. Ей
место в отдельном разделе.

\begin{definition}[Метакогниция]
Метакогниция --- способность Рин наблюдать и рассуждать о собственных
внутренних состояниях: о своём настроении, об уверенности в своих убеждениях, о
пробелах в памяти. Это операции второго порядка --- не «что я думаю о мире», а
«что я думаю о своём думании и своём состоянии».
\end{definition}

Метакогниция не само собой разумеется: базовая модель по умолчанию её не
имеет. Это проектируемая возможность, и у неё несколько конкретных ролей,
разбросанных по всей архитектуре, --- здесь мы их собираем.

\begin{itemize}
\item Тормоз руминации. Наблюдение «почему я так часто злюсь?» переводит из
режима «быть-в-настроении» в «наблюдать-настроение» и рвёт петлю первого
порядка (раздел выше). Технически --- это монитор автокорреляции $\mood_t$,
который сам замечает залипание и запускает корректив.
\item Эпистемическая честность. «Похоже на $X$, но $Y$ я не проверила» --- это
метакогнитивный отчёт об уверенности: Рин читает поле \texttt{conf} и зону
незнания ($|\texttt{conf}| < h$) и говорит о них, вместо ложного «я уверена»
(глава~\ref{ch:epistemics}).
\item Порождение любопытства. Замечая собственный пробел или высокую ошибку
предсказания, метакогниция ставит \texttt{need\_to\_discover} --- то есть сама
формулирует, чего Рин не знает и что стоит проверить
(главы~\ref{ch:epistemics}, \ref{ch:fusion}).
\item Самоотчёт. Способность сказать пользователю о своём состоянии («устала
от одного и того же бага», «мне сейчас не до шуток») --- это метакогниция,
считывающая настроение и мотивацию и выражающая их через Rinness
(глава~\ref{ch:voice}).
\end{itemize}

\smallskip\noindent\textbf{Субстрат метакогниции: самотелеметрия.}
Метакогниции не нужен отдельный граф. Её субстрат --- лог самонаблюдения:
кольцевой буфер временного ряда собственных состояний (вектор настроения) и
счётчиков команд (\texttt{need\_to\_*}, \texttt{fusion}, \texttt{contradiction},
\texttt{retrieval}) за часы/дни/недели. Словарь внутренних команд и состояний
конечен и дискретен, поэтому телеметрия компактна и структурирована, а
статистика по ней дешева. Это не граф и не свободный текст. Но выводы
метакогниции возвращаются в граф памяти как обычные узлы-убеждения-о-себе
(с \texttt{conf}, провенансом = телеметрия-свидетельство) --- и дальше живут по
общим правилам.

\smallskip\noindent\textbf{Детекция: дешёвый монитор чисел, а не сверх-LLM.}
Первый слой --- детекция --- набор дешёвых статистических детекторов на
буфере: скользящие среднее/медиана, тренд (наклон регрессии), частота (счёт в
окне), аномалия (z-оценка от базовой линии), автокорреляция (залипание).
Детерминированно, без LLM; срабатывают по порогу/аномалии. Пример:
«blue: медиана 0,7 > среднего 0,5 за 5 дней $\Rightarrow$ несколько всплесков
перетянули среднее, база грустная»; «\texttt{need\_to\_revise}: +60\% к
30-дневной базе». В «мозг» Рин уходит не сырой ряд (много, и LLM ненадёжно
считает медиану) и не готовый вывод, а компактная структурная сводка ---
выход детектора.

\smallskip\noindent\textbf{Аномалия --- это факт в графе, а не спецпайплайн.}
И вот основной ход, который важнее всего: сводку детектора пишем в граф памяти
как обычный факт-о-себе («за 5 дней \texttt{need\_to\_revise} +60\%»; «я в
последнее время грустная»). Уверенность высокая --- это измерение, а не
догадка. Дальше --- никакого выделенного «метакогниция-LLM» и никакого
спецпайплайна «аномалия $\to$ LLM $\to$ граф». Факт просто лежит в памяти и
подхватывается тем же общим циклом, что и любая другая мысль (retrieval,
фоновый fusion). Метакогниция здесь --- не модуль, а факты-о-себе, живущие в
графе по общим правилам.

Более того, сам факт уже работает --- вообще без LLM. Узел «я в последнее
время грустная» смещает самомодель, mood-congruent retrieval и самоотчёт
(«что-то я загрустила») через обычные механизмы. Значительная часть пользы
приходит просто от записи факта; объяснение «почему» --- опциональный более
глубокий слой, а не обязательная стадия.

\smallskip\noindent\textbf{Объяснение опционально: fusion-про-себя.}
Когда причина всё же нужна (гипотезы о том, почему аномалия), работает
fusion-про-себя (глава~\ref{ch:fusion}) --- но это тот же общий композитор, а
не прокладка: он запускается обычной salience-обработкой, когда факт-аномалия
всплывает. Гипотезы берут посылки из телеметрии: «почему больше
\texttt{need\_to\_revise}?» $\to$ H1 «много новых рассуждений», H2 «кладу
непроверенные факты», H22 «читала книгу с ошибками». Каждая --- провизорный
узел, низкая уверенность (в отличие от факта-измерения --- высокой),
провенанс из телеметрии.

Проверка --- по различающему предсказанию (научный метод): H1 $\Rightarrow$
должен вырасти и \texttt{fusion-count}; H2 $\Rightarrow$ у пересмотренных узлов
короткий интервал \texttt{commit->revision} и пропущенный
\texttt{need\_to\_discover}; H22 $\Rightarrow$ они делят один источник в
провенансе. \texttt{need\_to\_discover} кодирует, какой запрос выполнить;
curiosity/scrub его выполняет. Разрешённая гипотеза оседает в граф; если ничего
не разводит --- остаётся честное «неизвестно почему», а не выдуманная причина.

\smallskip\noindent\textbf{Eager или lazy --- из salience, а не из расписания.}
Раз аномалия --- это факт, «сразу разбираться или потом» перестаёт быть
отдельной развилкой в коде: это следствие salience самого факта. Срочная
аномалия (спираль настроения, руминация) несёт высокую значимость
$\Rightarrow$ всплывает и обрабатывается общим циклом немедленно (корректив
сейчас). Несрочная (\texttt{revise-rate} вырос) просто лежит и объясняется
лениво --- когда её достанет retrieval или придёт прямой вопрос о текущем
состоянии. Дешёвые детекторы при этом крутятся в фоне всегда (они дёшевы; их
интенсивность привязана к $\rho$-бюджету, глава~\ref{ch:walker}), а дорогой
общий LLM тратится лишь на то, что salience подняла наверх. Ни чистый
on-demand (руминация подавляет самозапрос именно тогда, когда он нужнее), ни
фиксированное расписание (расточительно) не нужны.

\smallskip\noindent\textbf{Защита от конфабуляции.}
Самообъяснение --- самый склонный к выдумыванию вид рассуждения: люди уверенно
конфабулируют причины своих настроений и поступков (классика --- Nisbett \&
Wilson, 1977). Поэтому телеметрия-лог здесь не удобство, а противоядие: он
делает самогипотезы проверяемыми против реальных данных, а не сочинёнными.
Метакогниция без телеметрии = конфабуляция; с телеметрией и различающими
проверками = фальсифицируемая.

\begin{pitfall}[Само-факты не должны стать самосбывающейся ловушкой]
Идея «аномалия --- факт в графе» сильна, но у неё есть тёмная сторона. Узел «я
в последнее время грустная» (\texttt{self\_model/mood\_summary}), показанный в
\texttt{context-packer} как обычный identity-факт, запускает петлю:
mood-congruent recall начнёт сам поддерживать это состояние. Поэтому такие
узлы: (1) диагностичны и time-bounded (с явным expiry, а не «навсегда»); (2)
пакуются в промпт как наблюдение о недавней телеметрии, а не как черта:
«recent telemetry suggests ...», а не «Rin is ...». Разница между «в последнее
время грустится» и «я грустная» --- это разница между диагнозом и приговором.
\end{pitfall}

\smallskip\noindent\textbf{Метакогниция о метакогниции (второй порядок).}
Тот же монитор можно навести на телеметрию самой метакогниции: частоту
порождаемых самогипотез, долю разрешённых, долю опровергнутых. Это
\textit{sanitizer} первого порядка: ловит analysis paralysis (много гипотез, ни
одна не разрешается $\to$ перестать копаться), ненадёжную интроспекцию
(самовыводы то и дело опровергаются $\to$ доверять самоотчёту меньше ---
прямая защита от конфабуляции) и залипшие петли. Механика та же, просто этажом
выше, поэтому дёшево. Третий порядок («думаю о том, как думаю о думании»)
даёт быстро убивающую отдачу и риск регресса --- потолок разумно ставить на
втором (аналогия со второй производной в матанализе).

\smallskip\noindent\textbf{Как это соотносится с литературой.}
Идея «метакогниция = System-2, наблюдающая System-1» и то, что мета-мысли
возвращаются в память и затем влияют на действия, --- из Tou, MacAdam, Tabor
(2024), \textit{«Metacognition is all you need?»} (поверх генеративных агентов
Park и др.). Что мы берём: самонаблюдение как отдельный модуль; мета-выводы
обратно в память; локальные модели; дешёвая модель для рутины, большая ---
для метакогниции. Где идём дальше (и где вопросительный знак в названии честно
себя выдаёт): у них детекция --- это LLM, оценивающий свой прогресс, а у нас ---
дешёвый статистический монитор, LLM только на объяснение; у них нет проверяемых
самогипотез с различающими запросами и защиты от конфабуляции; их
метакогниция целевая (задача), наша шире (настроение, эпистемика, гигиена
памяти) --- то, что они сами относят к future work. Значит «метакогниция ---
всё, что нужно»? Нет: нужен ещё субстрат самотелеметрии и дисциплина
проверяемых гипотез.

\begin{intuition}[Второй наблюдатель]
Полезно представлять метакогницию как второго наблюдателя, который смотрит не
на мир, а на первого --- на саму Рин: на её настроение, уверенность, ход мысли.
Именно этот второй наблюдатель отличает существо, которое переживает состояние,
от существа, которое знает, что переживает его, и может это назвать. Для
проекта, где цель --- личность, а не ассистент, это не украшение, а одна из
несущих опор.
\end{intuition}

\begin{openq}[Глубина и стоимость метакогниции]
Открыто, насколько «глубоко» гонять метакогницию: второй порядок («думаю о
думании») дёшев и явно полезен, но третий и выше («думаю о том, как думаю о
думании») быстро даёт убивающую отдачу и риск зацикливания на самонаблюдении.
Где ставить потолок и как выделять на метакогницию ресурсы (тот же
$\rho$-бюджет, глава~\ref{ch:walker}) --- пока не зафиксировано.
\end{openq}

\chapter{Топология правит потоком}
\label{ch:topology}

\section{Тезис главы}

У графовых систем вообще есть явная текущая сущность --- передача сообщений, и
её поток подчиняется топологии, как вода следует рельефу. У графа памяти Рин
явной «жидкости» нет. Но поток есть --- это поток изменений, поток
преобразований графа: активации, уверенности, консолидации. Тезис: неважно,
какова физика носителя, --- сама физика протекания информации подчиняется
топологическим законам. Любой процесс, чьё обновление локально (изменение узла
зависит от соседей), есть диффузия на графе и наследует его
спектрально-топологический транспорт. Это содержание теории графовых сигналов и
теории случайных блужданий.

\section{Стационарное распределение $\neq$ заданные вероятности}

Пусть мы локально задали вероятности переходов. Наблюдаемое распределение
частот посещений $\mathrm{distr}(A)$ отличается от заданного, потому что
вмешивается топология. Формально: стационарное распределение $\pi$ --- не
локальные переходы, а глобальный объект, неподвижная точка оператора переходов:
\begin{equation}
\pi = P^{\top}\pi, \qquad \sum_v \pi_v = 1.
\end{equation}

Для простого блуждания $\pi_v \propto \deg(v)$; для смещённого --- перронов
собственный вектор $P^{\top}$. Ты задаёшь переходы локально, а долгосрочная
частота посещений --- свойство всего графа. «Чем длиннее трек, тем наблюдаемее
эффект» --- это сходимость к $\pi$ за время перемешивания $t_{\mathrm{mix}}$,
управляемое спектральным зазором $1 - \lambda_2$. Короткие блуждания читают
локальную геометрию, длинные --- глобальную $\pi$. Телепорт \eqref{eq:ppr}
$\pi$ не убирает, а \emph{регуляризует} (сдвигает к смеси «равномерное +
рестарт», режет $t_{\mathrm{mix}}$): ограничивает эффект топологии, не отменяя.

\begin{openq}[Скорость света --- плохая, но поучительная аналогия]
Можно соблазниться сравнением с $c$: заданные вероятности как предельная
скорость, а среда (топология) искажает наблюдаемое распространение. Рабочая
часть --- «среда формирует распространение, есть структурный предел» --- уже
покрыта граф-диффузионной рамкой точно, а не по аналогии. Квантовый хвост
сюда тащить не нужно. Нам нужен не свет, а лапласиан.
\end{openq}

\section{Доказуемое свойство телепорта}

Почему сон (телепорт) необходим, а любопытство его не заменит? Любопытство
локально по вниманию --- смещённый обход от текущего фронтира; узел $X$ в
затенённом углу, оторванный от активного множества, недостижим, как ярко ни
свети фокус. Телепорт ломает достижимость: рестарт падает куда угодно. Это
ровно то, почему у PageRank есть телепорт: с $\alpha > 0$ цепь становится
эргoдической (неприводимой и апериодичной), а значит каждый узел посещается в
матожидании с частотой $\pi_v > 0$ --- ни один не заперт в стоке или на
острове. Телепорт даёт гарантию покрытия, которой у смещённого блуждания
доказуемо нет.

\begin{intuition}[Анти-энтропийный скраб]
«Не сразу; мы никогда не знаем, что и где сломалось, поэтому просто постоянно
проводим случайные инспекции» --- это зрелый паттерн из распределённых
хранилищ: anti-entropy scrub (проверка целостности в ZFS; Merkle-based gossip в
Dynamo/Cassandra). Сон Рин --- анти-энтропийный скраб по графу памяти.
\end{intuition}

\section{Пять следствий для работы Рин с памятью}

\begin{enumerate}
\item \textbf{Топологическая значимость («гравитационные ямы»).}
Хабы переактивируются независимо от задуманных вероятностей --- $\pi$ на них
концентрируется. Рин чаще вспоминает хаб-узлы, реже периферию (затенённые углы
--- регионы с малым $\pi$). Топология создаёт неявную важность: иногда
полезно, иногда случайный плотный кластер становится чёрной дырой припоминания.

\item \textbf{Узкие места / over-squashing.}
Два семантически близких, но топологически далёких кластера через тонкий мост
(низкая проводимость) плохо обмениваются активацией. Рин не свяжет факты,
«далёкие» по графу. Прямое ограничение fusion.

\item \textbf{Пересглаживание уверенности (over-smoothing).}
Повторная диффузная пропагация гомогенизирует уверенность к консенсусу ---
теряется контраст «уверена / не уверена». Дегенеративный отказ; лечится как в
графовых сетях: остаточные связи (якорь к исходному свидетельству), не чисто
диффузные апдейты, пере-заземление в провенанс.

\item \textbf{Кривизна как локатор ценного-и-задушенного.}
Отрицательно искривлённые рёбра --- мосты между сообществами --- это места
самых креативных слияний и максимального over-squashing. Значит самые ценные
ассоциации топология душит сильнее всего. Отсюда объективная причина
сна/телепорта: он --- rewiring, топологическая починка over-squashing обходом
узкого места случайным рестартом.

\item \textbf{$\pi$ как топологический субстрат идентичности.}
Какие узлы структурно центральны --- к чему Рин привычно возвращается ---
подложка её «озабоченностей», часть характера. Форма графа памяти --- часть
того, кто она.
\end{enumerate}

\begin{workedbox}[Over-squashing на мосте SBM]
Возьмём стохастическую блочную модель: два плотных сообщества, соединённых
$m$ рёбрами-мостом. Эффективное сопротивление между узлами разных сообществ
\[
\Reff(i,j) = (e_i - e_j)^{\top} L^{+}(e_i - e_j)
\]
растёт при уменьшении $m$ (мост уже --- сопротивление выше). Влияние дальнего
узла, измеряемое через якобиан $\partial h_i / \partial x_j$, ограничено сверху
через $\Reff(i,j)$. Значит при узком мосте fusion между сообществами доказуемо
слаб --- сколько ни води обычный walker. Единственное, что чинит, --- телепорт:
он кладёт рестарт в дальнее сообщество напрямую, минуя мост. Метрика для
стенда: частота cross-cluster fusion против $m$ (глава~\ref{ch:experiment}).
\end{workedbox}

\begin{figure}[h]
\centering
\begin{tikzpicture}[scale=1]
\foreach \p in {(0,0),(0.8,0.6),(0.8,-0.6),(1.6,0.2),(1.6,-0.4)} {
    \fill[black] \p circle (2pt);
}
\foreach \p in {(5,0),(5.8,0.6),(5.8,-0.6),(6.6,0.2),(6.6,-0.4)} {
    \fill[black] \p circle (2pt);
}
\draw[thick] (-0.4,-1) ellipse (2.4 and 1.5);
\draw[thick] (6.2,-1) ellipse (2.4 and 1.5);
\draw[very thick] (1.6,0.2) -- (5,0);
\node at (3.3,0.45) {мост};
\end{tikzpicture}
\caption{Два сообщества через узкий мост. Через него активация «задыхается»
(over-squashing); телепорт обходит его, прыгая напрямую.}
\end{figure}

\section{Операционные величины и осторожности}

Вести дело лучше случайно-блуждательными величинами: $\pi$, время
перемешивания, hitting/commute time, $\Reff$, проводимость/константа Чигера,
спектральный зазор. Кривизну держать как локатор узких мест, не как рычаг.

Что дёшево на лету: $\pi$ и «хабность» --- инкрементально дёшевы; $\Reff$ ---
дорого (через $L^+$), но аппроксимируемо. Это позволяет Рин самой знать, где
затенённые углы, и гнать туда телепорт-скраб прицельно.

Две осторожности. Во-первых, граф направленный, типизированный, знаковый ---
нужный оператор не симметричный лапласиан, а несимметричная матрица переходов,
спектр может быть комплексным, кривизна для орграфов менее чиста. Во-вторых,
граф динамический: fusion меняет саму топологию, и это диффузия на
время-зависимом графе, где поток и структура сцеплены обратной связью, --- за
рамками статических графов и самое интересное место.

\chapter{Математика retrieval: от локального шага к глобальному поведению}
\label{ch:retrievalmath}

\section{Зачем отдельная математическая глава}

В базовом изложении retrieval уже описан корректно: есть якоря, есть обход,
есть настроение, есть телепорты, есть разные политики. Но для курса этого мало.
Нужно увидеть разницу между \emph{локальным законом шага} и \emph{глобальной
физикой потока}. Именно здесь многие архитектурные интуиции либо укрепляются,
либо рассыпаются.

\section{Переходный оператор}

Пусть граф памяти задан вершинами $V$ и типизированными рёбрами $E$. Для фиксированного
режима работы зададим для пары соседних узлов $A,B$ ненормированную силу перехода
\[
S(A\to B)
=
w_{AB}\cdot
\mathrm{rel}(B\mid q)\cdot
\mathrm{mc}(B\mid \mood)\cdot
\mathrm{unc}(B)^{\beta_u}\cdot
\mathrm{gate}(A,B),
\]
где последний множитель $\mathrm{gate}(A,B)$ удобно выделить отдельно: он собирает
всё, что не является содержательной близостью, но влияет на сам факт допустимости
шага — например, запрет на некоторые типы рёбер в целевом retrieval, штраф за
повторное хождение по одному и тому же контуру, локальные бюджеты.

Нормировка даёт матрицу переходов
\[
P_{AB}=\frac{S(A\to B)}{\sum\limits_{C\in \Neigh(A)} S(A\to C)}.
\]
Если включён телепорт, фактический оператор имеет вид
\[
\widetilde P=(1-\alpha)P+\alpha \mathbf{1}r^\top,
\]
где $r$ — распределение рестарта. В retrieval оно сосредоточено на якорях, во
внутренней жизни смещено к зонам неопределённости, во сне может быть почти
равномерным или управляться картой неудовлетворённого покрытия.

\begin{intuition}[Почему важен именно оператор, а не набор красивых факторов]
Пока формула живёт как список слов — релевантность, настроение, новизна, — легко
вообразить, что эти ручки независимы. Но после нормировки они начинают бороться
за один и тот же переходный вес. Усилив настроение, можно случайно ухудшить
релевантность; усилив жадность, можно обрушить покрытие. У оператора есть
системные эффекты, которых не видно по отдельным множителям.
\end{intuition}

\section{Температура и режим выбора}

Пусть логарифм ненормированного веса шага равен $z_{AB}=\log S(A\to B)$. Тогда
softmax с температурой $T$ задаёт
\[
P_{AB}(T)=\frac{\exp(z_{AB}/T)}{\sum_C \exp(z_{AC}/T)}.
\]
Пределы понятны:
\[
\begin{aligned}
T &\to 0      && \Rightarrow \text{почти жадный выбор},\\
T &\to \infty && \Rightarrow \text{почти равномерный выбор по допустимым соседям}.
\end{aligned}
\]

Но practically важен не предел, а производная: насколько чувствительна вероятность
перехода к малому изменению веса. Именно здесь видно, что высокая температура — не
просто «больше случайности», а меньшая устойчивость к мелким локальным ошибкам в
оценках. Для debugging это полезно, для production-retrieval — опасно.

\section{Якоря, рестарты и two-step retrieval}
\label{sec:retrievalmathtax}

В исходной архитектуре уже сказано, что retrieval не ищет запись обходчиком с нуля,
а сначала находит якоря. Это место стоит формализовать.

Пусть есть запрос $q$. Тогда retrieval распадается на две задачи:
\[
q \longmapsto A(q)\subseteq V
\longmapsto
\text{локальное разворачивание окрестности}.
\]
Первая задача — \emph{anchor selection}; вторая — \emph{context expansion}. Смешивать
их нельзя.

\begin{definition}[Anchor recall]
Anchor recall — доля запросов, для которых на этапе выбора якорей в стартовое
множество попал хотя бы один узел из минимального носителя правильного ответа.
\end{definition}

\begin{definition}[Bridge rate]
Bridge rate — доля разворачиваний, в которых после попадания в правильную область
обход действительно дошёл до промежуточных узлов, без которых ответный фрагмент
остался бы несвязной кучей фактов.
\end{definition}

Эти две метрики принципиально различны. Низкий anchor recall чинится поиском
якорей; низкий bridge rate — уже топологией, температурой, телепортом и типами
допустимых переходов.

\section{Hitting time и доступность далёких зон}

Пусть $T_B$ — время первого попадания в узел $B$. Тогда $\E_A[T_B]$ — \emph{hitting
time} из узла $A$ в $B$. Для памяти это не абстракция: если интересный, но редкий
контрпример имеет слишком большой hitting time относительно типичных режимов сна
или фоновой разведки, то practically он \emph{не существует} для системы.

Отсюда инженерный вывод:
\begin{principle}[Редкий узел без пути — это не «редкое знание», а потерянное знание]
Если значимые узлы достижимы только по длинным и узким маршрутам, память не просто
медленно к ним приходит. Она систематически предпочитает не приходить вовсе, а
значит будет стабильно недокорректированной.
\end{principle}

\section{Спектральная щель и смешивание}

Если $\widetilde P$ эргодична, у неё есть стационарное распределение $\pi$. Скорость
смешивания управляется вторым по модулю собственным значением $\lambda_2$:
\[
\text{чем меньше } |\lambda_2|,\ \text{тем быстрее забывается старт}.
\]
Для памяти это двусмысленное благо.

\begin{itemize}
\item Быстрое смешивание полезно для покрытия и scrubbing.
\item Медленное смешивание полезно для удержания темы и целенаправленного retrieval.
\end{itemize}

Отсюда видно, почему одна и та же политика не может одинаково хорошо служить и
ответу, и сну. Retrieval хочет меньше забывать старт; сон, наоборот, должен уметь
выползать из локальной ямы.

\section{Проводимость и bottleneck starvation}

Пусть $S\subset V$. Проводимость
\[
\Phi(S)=\frac{\sum_{u\in S,\ v\notin S}\pi_u \widetilde P_{uv}}
{\min(\pi(S),\pi(V\setminus S))}
\]
измеряет, насколько легко поток выходит из области $S$. Малые значения $\Phi$
означают bottleneck: память часто крутится внутри района, но плохо ходит наружу.

В прикладном смысле это ведёт к двум симптомам:
\begin{enumerate}
\item cross-cluster fusion почти не случается;
\item контрсвидетельства, сидящие за узким перешейком, плохо добираются до ревизии.
\end{enumerate}

\section{Новая инженерная добавка: два бюджета контекст-пакера}
\label{sec:dualbudget}

Ни в исходном roadmap, ни в базовом дубле это не было выделено достаточно явно,
поэтому здесь вводится как \emph{расширение}, а не как ретушь старого текста.

Пусть retrieval уже вернул множество кандидатов $C$. Обычная ошибка — отдать в
контекст всё, что набрало высокий балл, в одном общем списке. Но для ответа нужны
по меньшей мере два разных ресурса:
\[
B = B_{\mathrm{evidence}} + B_{\mathrm{glue}}.
\]

\begin{definition}[Evidence budget]
Бюджет на узлы-свидетельства: наблюдения, факты, контрфакты, provenance-носители,
то есть то, что действительно несёт доказательную нагрузку.
\end{definition}

\begin{definition}[Glue budget]
Бюджет на связующий материал: переходные концепты, краткие обобщения, локальные
мосты, без которых набор свидетельств останется несвязным.
\end{definition}

Если не разделять эти бюджеты, происходит одно из двух. Либо всё съедает
свидетельственный материал, и ответ становится точным, но несобранным; либо всё
забивается «красивыми мостами», а фактическое основание усыхает. Разделение
бюджетов делает context packer не просто сортировщиком, а маленьким архитектором
локального доказательства.

\section{Что именно надо мерить на стенде}

Из этой главы естественно вытекает retrieval-аудит как кортеж
\[
\mathcal{R}
=
\bigl(
\text{anchor recall},
\text{bridge rate},
\text{contradiction hit rate},
\text{provenance coverage}
\bigr).
\]
Последние два компонента читаются так:

\begin{definition}[Contradiction hit rate]
Доля случаев, в которых при наличии в релевантной окрестности важного контрсвидетельства
retrieval действительно поднял его в рабочий контекст.
\end{definition}

\begin{definition}[Provenance coverage]
Доля включённых в контекст утверждений, для которых рядом лежит хотя бы минимальный
support-skeleton или иной объяснимый след происхождения.
\end{definition}

\begin{workedbox}[Почему одного recall недостаточно]
Пусть ответ на вопрос требует узлов $A$ и $B$, а их связывает мост $M$. Если якорь
нашёлся в районе $A$, а обход принёс $A$ и ещё пять соседей, но не принёс $M$ и $B$,
формально память «что-то помнит», practically — не отвечает. Recall на якорях может
быть высоким при полностью провальном bridge rate.
\end{workedbox}

\section{Вывод}

Математика retrieval нужна не ради строгости как украшения. Она нужна, чтобы не
перепутать четыре разных класса проблем:
\begin{enumerate}
\item ошибка выбора якоря;
\item ошибка локального разворачивания;
\item топологическая недоступность;
\item ошибка упаковки уже найденного материала.
\end{enumerate}

Без такого разделения система с плохим ответом выглядит просто как «что-то не так с
памятью». С ним становится видно, \emph{какой именно слой памяти провалился}.

\chapter{Fusion, вывод и консолидация}
\label{ch:fusion}

Обход достаёт то, что уже есть в графе. Но память живого существа ещё и
порождает новое: связывает далёкое, делает выводы, обобщает. Этот класс
операций мы называем fusion. Здесь разбирается, как именно из активированных
узлов рождается новый узел, чем вывод отличается от ассоциации и обобщения и
как всё это остаётся под эпистемическим контролем.

\section{Что такое fusion, механически}

\begin{definition}[Fusion, подробно]
Fusion --- оператор, который по совместно активированной конфигурации узлов
$\{A_1, \ldots, A_n\}$ создаёт новый узел $X$ и связывает его с посылками
рёбрами провенанса (обоснования). Новый узел рождается провизорно: с низкой
уверенностью и, при необходимости, с флагом \texttt{need\_to\_discover} на том,
что его различает.
\end{definition}

Механика в четырёх шагах:
\begin{enumerate}
\item Со-активация. Обход (retrieval, любопытство или сон) подсвечивает
несколько узлов одновременно.
\item Распознавание конфигурации. Оператор замечает, что активная конфигурация
согласована --- образует осмысленный паттерн (причинную цепочку, аналогию,
повтор).
\item Рождение узла. Создаётся $X$ с содержанием (сформулированным языковой
моделью) и провенансом $\{A_1, \ldots, A_n\}$.
\item Провизорная запись и пропагация. $X$ пишется с малым \texttt{conf},
затем пропагируется со знаком, карантином и дисконтом независимости
(глава~\ref{ch:propagation}).
\end{enumerate}

\begin{principle}[Fusion рождает гипотезу, а не факт]
Результат fusion --- всегда кандидат, а не запись. Это защищает от превращения
обаятельной ассоциативности в фабрику ложных убеждений
(глава~\ref{ch:epistemics}). Повышение до «факта» --- только через накопление
подтверждений или выживание в ночном скрабе.
\end{principle}

\section{Три вида fusion}

\smallskip\noindent\textbf{Ассоциация.}
Два далёких узла со-активируются, и рождается связь (или узел-мостик). Это
источник «неожиданных, но осмысленных» ходов мысли.

\begin{workedbox}[Ассоциация: Платон-боец]
Ночной обходчик случайно со-активирует два узла: «Платон был олимпийским
чемпионом по панкратиону (боец/борец --- по свидетельству Олимпиодора)» и
«\textit{platon} по-гречески --- „широкий, широкоплечий“». Fusion замечает
согласование «мощное телосложение бойца» + «имя означает „широкий“» и
рождает гипотезу: «„Платон“ --- прозвище по телосложению, а не данное имя»,
\texttt{conf = 0,2}, провенанс из двух узлов, флаг \texttt{need\_to\_discover}
на «проверить: прозвище ли это и каково настоящее имя». Вывод правдоподобен
(версий, к слову, несколько, но все сводятся к трактовке слова «широкий», а
настоящее имя по преданию известно) --- но именно гипотеза: если не пометить
низкой уверенностью, это тот самый преждевременный коллапс
(глава~\ref{ch:epistemics}). Существенно, что связывает не «Платон-философ»
сам по себе, а именно пара «боец» + «перевод».
\end{workedbox}

\smallskip\noindent\textbf{Вывод (инференция).}
Из посылок рождается заключение. Здесь принципиально важен случай $n$-арного
вывода, где ни одна часть посылок не достаточна.

\begin{workedbox}[$N$-арный вывод: корм котам]
Пусть в памяти Рин про пользователя есть три узла, и каждый подтягивает свой
контекст из эмбрионального графа (см. раздел ниже):
\begin{itemize}
\item $A$ --- «у него есть коты» $\Rightarrow$ контекст: котов надо кормить,
нужен корм;
\item $B$ --- «живёт в Португалии» $\Rightarrow$ контекст: магазины закрываются
рано;
\item $C$ --- «иногда ездит на пары» $\Rightarrow$ контекст: с пар возвращается
поздно.
\end{itemize}

Только совмещение всех трёх с их контекстами даёт вывод:
\[
\begin{aligned}
X =\ &\text{«купи котам корм на выходных --- в будни по пути с пар} \\
     &\text{можешь не успеть до закрытия»}.
\end{aligned}
\]

Существенно, что ни одна пара не достаточна: без «магазины рано» нет спешки,
без «поздно с пар» нет проблемы, без «нужен корм» нет повода. Вывод рождается
лишь при одновременном присутствии трёх посылок и их контекстов.
\end{workedbox}

Как мы это представляем сейчас, не меняя структуру: fusion создаёт узел $X$ с
провенансом $\{A,B,C\}$. Тем самым тройка-посылка кодируется как
узел-фактор, ссылающийся на три посылки, --- это factor-graph-кодировка
гиперребра $\{A,B,C\} \to X$. Иначе говоря, наша память уже наполовину
гиперграф: попарные рёбра $A$--$B$, $B$--$C$ могут отсутствовать (Португалия и
коты не связаны напрямую), а $n$-арный вывод всё равно рождается --- через
общий узел-обоснование. Стоит ли делать гиперрёбра первоклассными объектами
(со своим обходом и лапласианом) --- тема приложения~\ref{ch:hypergraphs}.

\begin{figure}[h]
\centering
\begin{tikzpicture}[
    node distance=1.3cm and 2.0cm,
    box/.style={draw, rounded corners, align=center, inner sep=6pt},
    prov/.style={draw=none, align=center}
]
\node[box] (a) at (0,1.2) {$A$: Португалия};
\node[box] (b) at (0,0) {$B$: коты};
\node[box] (c) at (0,-1.2) {$C$: универ};
\node[box] (x) at (5,0) {$X$: купи корм\\на выходных};
\node[prov] (p) at (2.5,1.2) {provenance};
\draw[-{Latex[length=2mm]}] (a) -- (x);
\draw[-{Latex[length=2mm]}] (b) -- (x);
\draw[-{Latex[length=2mm]}] (c) -- (x);
\end{tikzpicture}
\caption{$N$-арный вывод как узел-фактор $X$ с провенансом $\{A,B,C\}$.
Попарных рёбер между $A$, $B$, $C$ может не быть.}
\end{figure}

\smallskip\noindent\textbf{Как посылка «подтягивает свой контекст».}
Важная деталь, без которой $n$-арный вывод не работает: посылки участвуют не
«голыми», а вместе со своим контекстным ореолом. Механически это ограниченная
(1--2 хопа), смещённая целью/настроением spreading activation от каждой посылки
в эмбриональный граф:
\[
A = \text{коты} \Rightarrow \{\text{нужен корм}\},
\]
\[
B = \text{Португалия} \Rightarrow \{\text{магазины рано закрываются}\},
\]
\[
C = \text{универ} \Rightarrow \{\text{поздний возврат}\}.
\]

Fusion работает над объединением этих ореолов, а не над тремя голыми узлами;
композитор (языковая модель) собирает из «нужен корм + магазины рано + поздний
возврат» связный actionable-вывод. Провенанс $X$ --- посылки и реально
использованные контекст-узлы. Ореол смещён текущей целью: иначе подтянулась бы
вся Вселенная соседей, поэтому активируется лишь релевантное намерению (тот же
mood/goal-bias, что в retrieval, глава~\ref{ch:walker}). Именно поэтому
эмбриональный граф (глава~\ref{ch:embryonic}) не декорация: без рёбер
ConceptNet/ATOMIC у посылок нет ореолов --- и вывод не рождается.

\begin{openq}[Трудная часть --- распознать связность («трудные 80\%»)]
Подтянуть ореолы дёшево; трудно распознать, что их объединение образует
связную цепочку рассуждения, а не случайную кучу со-активаций. Кандидаты
разобраны в главе~\ref{ch:timescales} (раздел~\ref{sec:whatfusiondoes}):
LLM-судья / графовые шаблоны-схемы / обучаемый скорер связности --- ни один
пока не признан удачным. На старте --- LLM-композитор предлагает, а
$\varphi$-скорер и обратная связь фильтруют
(главы~\ref{ch:personality}, \ref{ch:voice}). Это открытый инженерный вопрос.
\end{openq}

\smallskip\noindent\textbf{Обобщение (абстракция).}
Много однотипных эпизодов сливаются в один узел более высокого уровня. Это и
есть консолидация (следующий раздел).

\section{Fusion сцепляется: sub-fusion --- это тоже fusion}

Виды fusion --- не разные механизмы, а один оператор, порождающий узлы разного
рода. И fusion сцепляется: обобщение, построенное одним шагом («sub-fusion»),
становится посылкой для вывода на следующем. Абстракция поднимает уровень,
вывод делает конкретный скачок, опираясь на неё:
\[
\text{факты} \xrightarrow{\text{абстракция}} \text{обобщение}
\xrightarrow{\text{вывод}} \text{гипотеза}.
\]

Механика та же (со-активация $\to$ новый узел с провенансом); «sub» лишь
означает, что этот узел --- промежуточное обобщение, мост для последующего
шага. (Это, к слову, клеточная картинка из приложения~\ref{ch:cellcomplex}:
высшая клетка, приклеенная к низшим.)

\begin{workedbox}[Дневная находка и полная эпистемическая петля: House / Cast Away]
\textbf{Фон.} Фоновый разведчик (глава~\ref{sec:background}), даже на скудном
дневном бюджете, со-активирует два далёких узла: «в House M.D. есть Уилсон» и
«в Cast Away есть „Уилсон“ --- волейбольный мяч».

\textbf{Sub-fusion (абстракция).} Оба узла порождают обобщение:
«Уилсон --- приём экстернализации внутреннего монолога ГГ» (у Cast Away --- из-за
одиночества на острове; у House --- потому что ГГ умнее всех, детектив,
интроверт). Родился узел более высокого уровня --- общая сценарная функция.

\textbf{Fusion (скачок).} Эта абстракция + совпадение имени «Уилсон» рождают
гипотезу: «а вдруг имя Уилсон в House --- отсылка к Cast Away?!» Провизорно,
низкая уверенность.

\textbf{Проверка (петля эпистемики, глава~\ref{ch:epistemics}).}
\texttt{need\_to\_discover} на «годы выхода» $\to$ внешняя проверка: Cast Away
(2000) раньше пилота House (2004) --- совместимо, уверенность растёт. Затем
\texttt{need\_to\_discover} на «есть ли ещё свидетельства» $\to$ проверка
опровергает: официальная версия --- оммаж Шерлоку Холмсу (House $\approx$
Holmes, адрес 221B), а Джеймс Уилсон $\approx$ доктор Ватсон (одна буква,
созвучие). Гипотеза понижается, на её место встаёт лучше обоснованная.

\textbf{Мораль.} Красивый, согласующийся с частью данных вывод оказался
ложным --- это тот самый преждевременный коллапс, пойманный проверкой. И
заметьте: «правильный» ответ (Holmes $\to$ House) --- сам по себе fusion,
просто с лучшим провенансом. Как говорит ГГ, «все лгут» --- поэтому
уверенность держим ниже единицы и идём проверять.
\end{workedbox}

\section{Консолидация как огрубление}

\begin{definition}[Консолидация, подробно]
Консолидация --- оператор, заменяющий множество похожих узлов/рёбер одним
обобщённым, сохраняя релевантные свойства и отбрасывая детали. Формально это
огрубление графа (coarse-to-local, глава~\ref{ch:timescales}).
\end{definition}

\begin{workedbox}[Консолидация: «Z --- частый гость»]
Десять эпизодов «Z был на стриме в понедельник», «...в среду», ... сворачиваются
в один узел «Z --- частый гость, любит доту». Выживают: частота, тёплая
валентность, тема (дота). Теряются: точные даты и реплики. Плохое огрубление
сохранило бы даты и потеряло валентность --- Рин помнила бы, когда заходил $Z$,
но забыла бы, как к нему относится. Значит критерий консолидации --- сохранять
аффективно- и реляционно-релевантное, а не всё подряд.
\end{workedbox}

Консолидация запускается двумя путями: во сне (фоновая переработка,
глава~\ref{ch:sleep}) и забыванием (когда эпизоды тускнеют, их дешевле слить в
обобщение, чем хранить порознь, глава~\ref{ch:timescales}).

\section{Кто пишет в узел и в каком виде}
\label{sec:whowrites}

Отдельно разберём вопрос, который легко замять: кто именно порождает содержание
нового узла --- LLM, автомат или Rinness --- и в каком виде оно ложится в
память. Полезно разбить на три подшага, потому что механизмы у них разные.

\smallskip\noindent\textbf{Решение (сливать ли).}
Для fusion это распознавание связной конфигурации (открытый «трудный 80\%»,
раздел~\ref{sec:whatfusiondoes}); для консолидации --- обнаружение избыточного
кластера. И тут важная поправка: группировать «однотипное» по строке нельзя.
Зритель мог «заходить», а мог «забегать» --- чистый автомат, ключующийся по
глаголу, тут же ломается, ведь мы не типизируем глаголы, которыми к Рин
обращаются. Синонимию ловит не автомат, а либо LLM при записи (нормализует
«забежал» в канонический тип связи \texttt{visited}), либо кластеризация по
близости эмбеддингов при консолидации. Поэтому честная картина: автомат делает
бухгалтерию (счётчики, хаб, перекидка рёбер, усреднение валентности),
эмбеддинги --- группировку «похоже по смыслу», а маленький LLM --- итоговый
ярлык. «Только автоматом» проходит лишь случай точных дублей.

\smallskip\noindent\textbf{Содержание узла.}
Новое осмысленное содержание (вывод, абстракция, ярлык обобщения) пишет
языковая модель --- шаблоном такое не собрать. Но та же самая модель, а не
какая-то отдельная: «нейтральный режим» и «Rinness» --- это два настроения
одной модели, а не два мозга. Для записи в память её держат в спокойном,
откалиброванном «девелоперском» режиме вектора настроения: рассудочно, минимум
стиля. Если голос Рин зашит в веса (LoRA), совсем нейтральной её не сделать ---
останется лёгкий налёт, и это не страшно, лишь бы не перегибать: в память важно
содержание, а тон всё равно родится заново на выходе. Автомат при этом ставит
все метаданные (провенанс, уверенность, \texttt{need\_to\_discover},
\texttt{importance}, домен) --- детерминированно.

\smallskip\noindent\textbf{Кто именно пишет --- и почему это не «Rinness».}
Пишет в граф та же Рин, но в спокойном режиме записи, а не в режиме голоса.
Здесь важно не путать «Рин» и «Rinness»: Rinness --- это узко выходная
стилизация (голос, гл.~\ref{ch:voice}), которая превращает извлечённое
содержание в реплику пользователю. Вот эта стилизация в граф не пишет. Одна
Рин, два режима: режим записи (спокойный) кладёт нейтральное содержание, режим
голоса (Rinness) звучит на выходе. Отсюда принцип: храни спокойно, стилизуй на
выходе. Причина простая: узел переиспользуется в разных будущих настроениях, и
если вморозить в него сегодняшний тон, он будет отдавать им и через месяц,
когда Рин к тому же зрителю потеплеет.

\smallskip\noindent\textbf{В каком виде хранить: содержание против стиля.}
Итак, в память кладём нейтральное содержание («зритель Z проиграл в доте») плюс
её отношение в метаданных --- валентность и, для значимых случаев, отдельный
узел-реакцию. Личность памяти живёт в отношении, а не в замороженных словах.
Класть готовую фразу «в голосе» («Z снова слил, ну куда ему») --- плохо: это
вмораживает транзиентный тон, память пухнет, а стилизовать на выходе уже
нечем.

Исключение --- вспышечная память. У людей рядовые события помнятся по сути, а
редкие пиково-эмоциональные --- вплоть до того, что было сказано и как
ощущалось. По аналогии: для редких моментов с высокой значимостью (большой
$|\Delta \mood|$) можно дополнительно сохранить дословный снимок --- реальную
реплику/тон того мгновения --- как отдельный связанный узел, а не как
каноническое содержание. Так по умолчанию память нейтральна и переиспользуема,
но яркие пики сохраняют свою фактуру.

\section{Выход fusion --- гипотеза, а не установленный факт}

Снимем терминологическую путаницу, потому что от неё зависит понимание всего.
Каждый узел --- это убеждение с уверенностью $\conf \in [-1,1]$. «Факт» в
широком смысле --- это просто любой узел, и да, fusion создаёт узел. Но узлы
различаются по уверенности: прямое наблюдение или измерение получает высокую
уверенность, а вывод (fusion) --- низкую. Именно это имелось в виду под
«fusion не коммитит факты»: коммит --- это запись как установленного,
высокоуверенного убеждения, которым дальше пользуются как истиной; fusion же
рождает узел-гипотезу. Повышение гипотезы до установленного статуса происходит
позже --- через корроборацию или проверку, а не в момент рождения.

Почему это критично: без такого разделения обаятельная ассоциативность
превратилась бы в фабрику ложных убеждений (пример House / Cast Away выше).
Поэтому fusion смел, но не безрассуден: гипотезы рождаются свободно, но влияние
гипотезы на соседей гейтируется её собственной (низкой) уверенностью и
надёжностью посылок --- карантин, глава~\ref{ch:propagation}; а различающие
пробелы, которых Рин не знает, помечаются \texttt{need\_to\_discover}, чтобы
любопытство или внешний вход их потом добыли (глава~\ref{ch:epistemics}).
Дальнейшую судьбу гипотезы --- проверку, повышение или понижение --- ведёт уже
эпистемический слой; fusion лишь честно рождает её с низкой уверенностью и
передаёт дальше.

\begin{principle}[На первых версиях: генератор кандидатов с аудитом]
Fusion --- главный источник и магии, и мусора. Пока распознавание связности
со-активации не решено (открытый вопрос), fusion на ранних версиях должен быть
не «автоматической когнитивной магией», а генератором кандидатов с аудитом.
Дефолтная позиция смещена в сторону осторожности: больше кандидатов в
\texttt{pending}, меньше \texttt{commit}. Принятые fusion потом ревизуются
отложенно (precision после проверки, глава~\ref{ch:experiment}), и только
выжившие повышаются. Лучше недодать смелых узлов, чем зарастить память красивым
мусором; смелость безопасно включим, когда научимся мерить связность.
\end{principle}

\begin{principle}[Разведи candidate\_text и canonical\_node]
Fusion всё ещё сильно опирается на LLM-композитор, и это риск: модель выдаёт
красивый текст, но плохо индексируемый и плохо ревизуемый узел. Разделяем два
формата. LLM формулирует \texttt{candidate\_text} --- гипотезу человеческим
языком. Но коммит идёт не текстом, а через schema-writer в \texttt{canonical\_node}
со структурой:
\texttt{\{node\_kind, claim\_type, entities, time\_scope, confidence, provenance, distinguishing\_unknowns\}}.
Иначе получишь хорошо написанные, но нетипизированные узлы, которые нельзя ни
точно найти (forensic-ретрив, гл.~\ref{ch:runtime}), ни аккуратно ревизовать.
Красота --- на входе, схема --- на коммите.
\end{principle}

\section{Когда со-активация --- повод для реплики, а не для узла-факта}
\label{sec:antifusion-joke}

Не всякая яркая со-активация должна оседать убеждением. Иногда две всплывшие
вещи --- это повод что-то сказать, а не повод что-то запомнить как факт о мире.
Это ключевое различение: полезное для реплики $\neq$ достойное записи фактом
(anti-fusion). Разберём на каноничном случае.

\begin{workedbox}[Чапаев и «нюанс»: шутку не пишем, реакцию --- пишем]
В памяти со-активировались «Чапаев» и слово «нюанс». Walker сводит их рядом, и
всплывает ассоциация: «тут просится анекдот». Что происходит по нашим
механизмам.

(1) \textbf{Откуда берётся анекдот.} Не из графа как хранимый текст. Граф даёт
только триггер --- «эта пара уместна для шутки». Сам текст анекдота порождает
базовая LLM (она знает его из предобучения) --- ровно как с обучением: граф
всплыл повод, модель поставляет содержание.

(2) \textbf{Что НЕ записывается.} Сам анекдот не коммитится узлом-фактом. Связь
«Чапаев relates-to нюанс через этот анекдот» --- это \texttt{candidate\_text}
для реплики, а не убеждение о мире; anti-fusion-гейт его как факт не
пропускает. Иначе память заросла бы «фактами», которые на деле каламбуры.

(3) \textbf{Что записывается.} Два узла другого рода. Во-первых, episode:
«рассказала анекдот про Чапаева-нюанс собеседнику $Z$» (со временем и
провенансом) --- это было, это факт о событии, не о мире. Во-вторых, реакция
$Z$ обновляет \texttt{reaction\_pattern} через Beta-постериор, условленный на
отношении (гл.~\ref{ch:personality}): зашло $\to \alpha{+}1$, поморщился
$\to \beta{+}1$. Salience-гейт решает судьбу: сдвинула реакция настроение
($|\Delta \mood|$) --- консолидируется; «ну ок» --- тускнеет.

(4) \textbf{Бонус --- крючок, а не шутка.} Если такой юмор у $Z$ стабильно
заходит, осесть может лёгкий \texttt{reaction\_pattern}: «пара
Чапаев+нюанс $\to$ этот анекдот заходит $Z$». Заметьте: выучивается модель того,
что кому заходит, а не сам анекдот хранится фактом. Комедийный тайминг ---
Beta-апостериор, не строка в базе.
\end{workedbox}

Мораль в одну строку: fusion-подобная со-активация может кончиться репликой, а
осесть в памяти должны событие речи и реакция на неё --- потому что важна не
шутка, а то, как она легла. Тот же принцип защищает память от превращения в
свалку острот, мемов и подходящих-к-случаю ассоциаций.

\begin{intuition}[Инженер, а не фантазёр]
Правильная Рин --- «инженер, которую делают бывшие физики-фундаменталисты»:
она охотно строит гипотезы, но по умолчанию ставит уверенность ниже единицы и
идёт проверять. Fusion даёт первое (смелость ассоциаций), эпистемическая
гигиена (глава~\ref{ch:epistemics}) --- второе (дисциплину). Одно без другого
--- либо сухой поисковик, либо обаятельный врун.
\end{intuition}

\chapter{Многомасштабность, слои и забывание}
\label{ch:timescales}

\section{Протекающий интегратор: точный смысл}

Память Рин расслоена по временным масштабам: сессия $\to$ день $\to$ сезон
$\to$ год $\to$ жизнь. Каждый слой --- протекающий интегратор
(\textit{leaky integrator}): состояние, которое накапливает вход и понемногу
«протекает» (теряет часть себя) на каждом шаге.

\begin{definition}[Протекающий интегратор]
Состояние $s$, обновляемое под входом $x_t$ по правилу
\begin{equation}
s_t = \left(1 - \frac{1}{\tau}\right)s_{t-1} + \frac{1}{\tau}x_t,
\qquad \tau > 1.
\label{eq:leaky}
\end{equation}
Коэффициент $(1 - 1/\tau)$ --- доля прошлого состояния, которая
\emph{удерживается}; $1/\tau$ --- доля нового входа, которая \emph{впитывается}.
Постоянная времени $\tau$ задаёт «медленность»: большое $\tau$ --- слой почти
не меняется за шаг (удерживает долго, впитывает по чуть-чуть), малое $\tau$ ---
резко реагирует и быстро забывает.
\end{definition}

\begin{intuition}[Ведро с дыркой]
Протекающий интегратор --- это ведро с дыркой в дне, куда капает вода (вход).
Уровень воды $s$ растёт от капель и падает от протечки. Широкая дырка
(малое $\tau$) --- уровень скачет за каждой каплей и быстро сходит; узкая
дырка (большое $\tau$) --- уровень меняется медленно, зато держится долго.
«Сессия» --- ведро с широкой дыркой, «жизнь» --- с почти закрытой.
\end{intuition}

\section{Два разных смысла «слоёв» --- не путать}
\label{sec:twolayers}

Здесь важная развилка, которую легко замылить: слово «слои» и «протекание»
означают разное для настроения и для памяти.

\smallskip\noindent\textbf{Настроение --- параллельные интеграторы, переноса нет.}
Слои настроения (глава~\ref{ch:mood}) --- это несколько интеграторов
\eqref{eq:leaky} с разными $\tau_\ell$, работающих параллельно на одном и том
же потоке событий. Старший слой не питается от младшего; оба питаются
событиями, просто старший интегрирует медленнее. «Младший дрожит сильнее» ---
это высокочастотная компонента того же сигнала. \textbf{Переноса содержания
между слоями нет} --- есть разночастотная фильтрация одного входа, а
актуальное настроение --- суперпозиция.

\smallskip\noindent\textbf{Память --- перенос есть, это повышение слоя консолидацией.}
Слои памяти устроены иначе: узел \emph{повышается} из быстрого слоя в
медленный, когда консолидация делает его устойчивым (раздел~\ref{sec:consolidation}).
Вот это и есть настоящее «протекание младшего в старший»: не диффузия
сигнала, а \emph{перемаркировка} узла
(\texttt{layer}: \texttt{session} $\to$ \texttt{day} $\to \ldots$) при накоплении
силы. Механизм известен по форме; конкретные пороги повышения --- открыты.

\begin{openq}[Пороги повышения слоя]
При какой накопленной силе / частоте обращений / значимости узел повышается в
следующий слой? Разумные кандидаты: порог по базовой активации ACT-R
(раздел~\ref{sec:decay}); порог по числу независимых подтверждений; порог по
$|\conf|$. Точные значения и то, единый ли это порог для всех типов узлов, ---
пока не зафиксированы; это параметр для стенда (глава~\ref{ch:experiment}).
\end{openq}

\section{Консолидация как огрубление графа}
\label{sec:consolidation}

Переход от эпизодов к смыслам --- это огрубление (\textit{coarsening}): десять
записей «зритель заходил в понедельник», «...в среду», ... заменяются одним
узлом «зритель --- частый гость». Мелкая структура схлопывается в крупную;
часть информации теряется намеренно.

\begin{principle}[«Консолидация» --- это процесс, а не один оператор]
Здесь важно снять терминологическую путаницу, из-за которой легко запутаться
(и мы путались). Под словом \textbf{консолидация} мы понимаем не единичный
оператор «свернуть $A, B, C$ в $X$», а \emph{процесс} комплексной отработки
нескольких разных операторов, каждый со своими правами и рисками:
\begin{enumerate}
\item \textbf{Fusion / вывод} --- рождает новый смысл («поездка на Азоры»,
«ходит по нечётным»). Посылки \emph{не} трогает (гл.~\ref{ch:fusion}).
\item \textbf{Summary-overlay (мягкая консолидация)} --- заводит узел-обобщение
\emph{поверх} исходных, дети остаются, добавляются рёбра членства.
Безопасный дефолт.
\item \textbf{Dedup / merge} --- склеивает \emph{реально один и тот же} факт из
разных источников (это слияние identity, а не абстракция).
\item \textbf{Prune / сжатие с потерей} --- вот единственный оператор, который
\emph{удаляет} детали. Редкий, поздний, жёстко гейтируемый.
\end{enumerate}
Почти вся ранняя консолидация --- это (1)+(2): создать смысл и структуру
доступа, ничего не разрушая. Опасен только (4), и то лишь при
преждевременном применении. Когда ниже написано «консолидация сносит
избыточное» --- речь про оператор (4) внутри процесса, а не про процесс
целиком.
\end{principle}

\begin{intuition}[Почему нужен активный синтез, а не один только decay]
Легко спросить: если узлы и так тускнеют по непосещаемости, зачем вообще
что-то сворачивать? Затем, что \emph{decay} (смерть по непосещаемости) ---
\emph{слепое выбрасывание}: он теряет данные, \emph{не извлекая из них закон}.
Если тридцать эпизодов «приходил по нечётным» просто протухнут, память
потеряет их --- и так и не узнает, что он ходит по нечётным. А связка
\emph{синтез (fusion) + поздний прунинг} --- то, что мы и зовём
консолидацией как процессом, --- делает \emph{синтез, потом выброс}: сначала
дистиллирует gist, и лишь потом ретирует ставшую лишней деталь. Формула:
\emph{выучить урок (родить обобщение) и выбросить домашку (снести избыточное)}.
Аналогия: живой ум не помнит каждый восход, но твёрдо знает «солнце встаёт
на востоке» --- инвариант извлечён, экземпляры отброшены; будь в системе
только decay, она бы просто забывала восходы, а правила бы не получила. И
ещё: часто-посещаемый избыточный узел по decay никогда не умрёт (его же
трогают) и будет вечно раздувать граф --- только активный синтез его
свернёт.
\end{intuition}

Как понять, что огрубление \emph{хорошее}? Здесь нужен мост к спектральной
математике огрубления. Введём:
\begin{itemize}
\item $L_f$ --- лапласиан \emph{мелкого} (\textit{fine}) графа; лапласиан
описывает, как сигнал диффундирует по графу (кто на кого влияет и насколько);
\item $L_c$ --- лапласиан \emph{крупного} (\textit{coarse}) графа после
огрубления;
\item $P$ --- оператор \emph{проекции}: сопоставляет каждому мелкому узлу его
крупную группу (например, десять эпизодов $\mapsto$ узел «частый гость»);
\item $\Phi(X,L)$ --- оператор распространения (одна или несколько итераций
передачи сообщений) сигнала $X$ по графу с лапласианом $L$.
\end{itemize}

Хорошее огрубление \emph{коммутирует} с распространением:
\begin{equation}
\left\| P\Phi_f(X,L_f) - \Phi_c(PX,L_c) \right\| \le \varepsilon.
\label{eq:commute}
\end{equation}
Читается так: неважно, в каком порядке действовать --- сначала «подумать» на
мелком графе и затем сжать, или сначала сжать и «подумать» на крупном, ---
результат должен почти совпасть. Норма $\|\cdot\|$ берёт разницу по всем
входным сигналам $X$; $\varepsilon$ --- мера потери. Малое $\varepsilon$ =
огрубление сохранило то, что нужно для рассуждения; большое = сжатие
сломало динамику.

\begin{workedbox}[Что должно выжить при консолидации]
Свернём десять эпизодов визитов в узел «частый гость». \emph{Выжить} обязаны:
частота (была высокой), устойчивая валентность (тёплая), тип связи
(регулярность). \emph{Потеряться} можно: точные даты, дословные реплики.
Плохое огрубление сохранит даты и потеряет валентность --- тогда Рин будет
помнить когда заходил зритель, но забудет, как относится к нему; это большая
$\varepsilon$ по аффективной оси. Значит критерий огрубления --- сохранять
аффективно- и реляционно-релевантное подпространство, а не спектр целиком.
\end{workedbox}

\subsection{Коллизии предложений: голоса, а не очередь}
\label{sec:consolidation-collisions}

Огрубление предлагают обходчики (гл.~\ref{ch:walker},
раздел~\ref{sec:proposecommit}), и предложения будут \emph{перекрываться}.
Пусть за ночь на узлах $A, B, C, D, E$ накопилось: RW1 хочет слить
$\{A, B, C, D\} \to X_1$; RW2 --- $\{B, C, D, E\} \to X_2$; RW3 ---
$\{A, B\} \to Y_1$; RW4 --- $\{C, D, E\} \to Y_2$. Наивный коммиттер
«кто первый --- тот и слил» построит мусор, зависящий от случайного
порядка. Правильный ход --- понять, что \emph{отдельное предложение обходчика
это не решение, а свидетельство}: «этих тянет вместе». Тогда коллизия ---
не спор за очередь, а \emph{сигнал для агрегации}.

Механически коммиттер копит предложения в окне (обычно за ночной проход) и
строит \emph{взвешенный со-консолидационный гиперграф} на исходных узлах: вес
$c(i,j)$ = как часто обходчики клали $i$ и $j$ в один предлагаемый кластер,
взвешенно по support и --- важно --- \emph{дисконтированно по времени}: голоса
в буфере экспоненциально стареют, иначе кластеризация будет учитывать давно
неактуальные связи. По этому графу выделяются \emph{кандидатные кластеры}.

И вот тут --- принципиальная поправка: не надо загонять всё в \emph{единый
жёсткий раздел} (\textit{partition}). Память --- не разбиение; один эпизод
«смешной вечер в универе с Анной, когда я устал» законно входит и в «про
универ», и в «про Анну», и в «про усталость», и в «смешное». Раз parent у нас
--- обратимая надстройка-summary, а не разрушение мелких узлов,
\textbf{перекрывающиеся кластеры разрешены}. Поэтому: выделяем кандидатные
кластеры (модулярность / кластеризация --- как один из инструментов, не как
истина в последней инстанции), разрешаем перекрытие, а коммиттер выбирает
\emph{неконфликтующие} коммиты или заводит \emph{несколько} parent-summary,
если они сжимают \emph{разные оси}. В примере $\{B, C, D\}$ --- сильное ядро
(RW1+RW2), $\{A, B\}$ и $\{C, D, E\}$ могут ужиться как разные полезные
обобщения; а вот если $X_1$ и $X_2$ оказываются почти дубликатами по одной
оси --- их сливаем. Решает не порядок поступления, а согласие свидетельств и
польза.

Четыре оговорки, без которых механизм наивен. \emph{Дисконт независимости}:
два обходчика, прошедшие по общим путям, --- это часто один сигнал,
засчитанный дважды; вес $c(i,j)$ дисконтируется по перекрытию провенанса
предложений (тот же Bloom / MinHash, что и в пропагации,
гл.~\ref{ch:propagation}). \emph{Гистерезис}: сливаем только то, что
рекуррентно всплывает несколько ночей; одноразовое предложение --- шум, ему
рано менять структуру. \emph{Критерий пользы, а не только повторяемости}:
новый parent оправдан, если он \emph{сжимает} --- уменьшает будущую стоимость
ретрива / packer'а или повышает качество реакции; частая со-встречаемость без
выигрыша в сжатии --- не повод для обобщения. \emph{Обратимость через
стадийность}: снос избыточных детей идёт не сразу топором, а стадийно ---
сперва \texttt{pending} / \texttt{tombstone} (см. ниже) с сохранением
\texttt{support-skeleton} (гл.~\ref{ch:runtime}), и только устойчиво
подтверждённое обобщение доводит прунинг до конца. Поэтому цена даже
неудачного слияния низкая --- его можно откатить, и это снимает надобность в
идеальном арбитраже с первого раза.

\begin{intuition}[Комитет, а не турникет]
Представьте не турникет, где предложения проходят по одному и первое
побеждает, а комитет, который выслушивает все голоса разом и выносит одно
связное решение. Обходчики --- свидетели с разной надёжностью; коммиттер ---
тот, кто взвешивает показания (со скидкой на сговор --- общий провенанс) и
решает, какие обобщения завести --- при этом узел может состоять сразу в
нескольких, если они про разное.
\end{intuition}

\subsection{Что сносим, что оставляем --- и кто это решает}
\label{sec:consolidation-prune}

Итак, консолидация в итоге \emph{выводит избыточные узлы из активного графа}
--- иначе он только растёт. Уточним слово: не «убивает» (физическое стирание
--- это уже не память, а бухгалтерия), а переводит узел в одно из состояний
--- \texttt{active} $\to$ \texttt{archived} / \texttt{support-skeleton} /
\texttt{tombstone} (event-log хранит исходное всегда, гл.~\ref{ch:runtime}).
Полезно увидеть, что в процессе слиты \emph{две} разные операции, и решение о
второй принимается тоньше, чем «дети консолидации умирают»:
\begin{itemize}
\item \textbf{(а) рождение обобщения} --- создаётся узел-gist («приходит по
нечётным», «частый гость»). Это \emph{создание нового}, по духу близко к
fusion.
\item \textbf{(б) прунинг избыточных детей} --- сносится то, что новый gist уже
полностью покрывает.
\end{itemize}
Кто решает --- коммиттер, а не обходчик, и решает он про (б)
\emph{поэкземплярно, по избыточности}, а не «по кластеру, которому узел якобы
принадлежит». Обходчик лишь свидетель: «здесь есть сжимаемая конфигурация».
Формально перед прунингом ребёнка $v$ коммиттер прогоняет
\emph{residue-тест} --- «есть ли у $v$ уникальный остаток?»:
\begin{enumerate}
\item уникальные внешние рёбра (мост к другим частям графа, которые
осиротеют)?;
\item уникальный timestamp или actionable-деталь (дата, билет, обязательство)?;
\item высокая значимость / валентность?;
\item связь с другим обобщением?;
\item естественный запрос, на который $v$ отвечает лучше parent'а?
\end{enumerate}
Если хоть один пункт да ($\mathrm{residue} \ne 0$) --- \emph{не прунить},
максимум архивировать или оставить support-skeleton. Только при чистом нуле
$v$ выводится из активного графа; и даже тогда он может просто потухнуть по
ACT-R, если к нему перестанут обращаться. В чистом случае «по нечётным» под
критерий сноса попадают \emph{все тридцать} экземпляров (они взаимозаменяемы,
остатка и мостов нет) --- остаётся один узел; консолидация там максимально
беспощадна, и правильно. Мостовой богатый узел (эпизод, на котором висят два
разных обобщения) --- исключение не потому, что «мы не прунуем», а потому,
что он \emph{не избыточен}.

\smallskip\noindent\textbf{Право на схлопывание: экземпляр против составляющей.}
Заострим, что именно делает ребёнка избыточным, --- это и есть та тонкая
грань, которую легко проскочить. Схлопывать (с потерей) можно только
\emph{класс эквивалентных экземпляров} одного паттерна: тридцать «заходил по
нечётным», сто «привет», каждый восход. Они взаимозаменяемы
(\texttt{instance\_of}), и gist их полностью замещает. А \emph{разнородный
букет составляющих} одной темы (\texttt{part\_of}) схлопывать нельзя
\emph{в принципе} --- там нечего терять, там всё уникально. Отсюда типичная
ошибка: «fusion сгруппировал $A, B, C$» \emph{не} значит «$A, B, C$ можно
схлопнуть». Группировка-по-теме $\ne$ поглощение-избыточности. Разберём на
неигрушечном примере.

\begin{workedbox}[Поездка на Азоры: fusion есть, схлопывать нечего]
Пять фактов: (A) купил плавки, (B) купил зонтик, (V) билет на Азоры на
5 сент., (G) автобус аэропорт Лиссабона $\to$ Порто на 9 сент., (D)
конференция в Порто 9--12 сент. Напрашиваются два fusion:
$X_1 =$ «отпуск на Азорах» ($A, B, V$), $X_2 =$ «маршрут
Азоры $\to$ Лиссабон $\to$ конференция» ($V, G, D$). Но \textbf{схлопывать
здесь нечего}: $A, B, V, G, D$ --- не экземпляры паттерна, а разнородные
составляющие. Прунинг убил бы ровно то, что нужно для действия: даты,
билеты, предметы, обязательства. Поэтому --- только (1)+(2): рождаем
$X_1, X_2$ как overlay, факты оставляем, вешаем ролевые рёбра
(\texttt{equipment\_for}, \texttt{leg\_of}, \texttt{purpose\_of}).

Тонкость \texttt{event vs state}: «купил плавки» (\texttt{event}) порождает
отдельный узел «вероятно, владеет плавками» (\texttt{state}, другой
\texttt{node\_kind}, гл.~\ref{ch:substrate}); \texttt{event} со временем уйдёт
в архив, а \texttt{state} живёт, пока не появится «потерял / выкинул / не
подошли». Потому и на «Рин, поищи плавки» через два месяца она поднимает
\texttt{state} как \emph{контекст} и в худшем случае уточняет («вроде брал в
августе --- ищем на замену?»), а не блокирует запрос старым фактом. Это
отдельный инвариант --- \emph{память --- контекст, а не команда}: прошлый
факт не отменяет текущее намерение, а лишь оттеняет его.
\end{workedbox}

\begin{principle}[Два инварианта консолидации]
Первый: вывод обходчика --- \emph{свидетельство}, а не власть. Коммиттер
решает задачу \emph{отбора модели} --- какой набор новых узлов лучше
объясняет данные при минимальной сложности (MDL-настроение: выигрыш в
сжатии + связность + повторяемость + польза $-$ избыточность $-$ сложность
$-$ эпистемический риск), а не исполняет предложения слепо.

Второй: перекрывать обобщения --- можно, дублировать identity --- нельзя.
Один исходный узел может быть посылкой нескольких parent-summary через рёбра
членства (\texttt{member\_of} / \texttt{provenance}); но клонировать его на
несколько копий в графе запрещено.
\end{principle}

\begin{pitfall}[Почему нельзя клонировать узлы («графовый мейоз»)]
Соблазн разрулить пересечения так: развернуть общий узел $C$ в близнецов
$C_1, C_2$, раздать по разным обобщениям, оригинал снести --- и предложения
обходчиков станут чистым разбиением. Формально приём валиден
(\textit{graph unrolling}), но семантически ломает память в двух местах.
\emph{Расщепление identity убивает TMS}: опровергнут позже факт $C$ ---
ретракция поправит один клон, второй останется фантомом, лавинная ревизия
уверенности сломается. \emph{Разрыв моста}: общий $C$ был перешейком, по
которому активация текла между $X_1$ и $X_2$ («раздражается»
$\leftrightarrow$ «не выспалась» связывались \emph{через} $C$); разнёс копии
--- потерял вывод навсегда. И «раздваиваются только на время» --- иллюзия:
дочерние не сливаются обратно, двойной след остаётся \emph{навсегда}.
Спасение идеи: близнецы допустимы \emph{только как черновая вычислительная
техника} в блокноте коммиттера (чтобы независимо оценить кандидатов), но в
\emph{коммите} общий узел остаётся общим --- identity в графе не двоится.
\end{pitfall}

\subsection{Кандидаты, наследование линков, слияние уверенностей}
\label{sec:consolidation-mechanics}

Три механических вопроса, которые всплывают на реализации.

\paragraph{Жаккард / LSH --- сигнал, а не критерий.}
Соблазн: «полное перекрытие окрестностей $\Rightarrow$ поглощаем». Опасно.
Структурная эквивалентность $\ne$ семантическая замещаемость: «купил плавки»
и «купил зонтик» имеют почти одинаковых соседей (Азоры, пляж, лето, покупка),
но не взаимозаменяемы --- один отвечает «есть ли плавки», другой «есть ли
зонтик». Поэтому индекс Жаккара (и его дешёвая оценка через
MinHash / LSH за $O(|V|)$ вместо $O(|V|^2)$ --- да, это правильный способ
искать кандидатов ночью) годится как \emph{кандидатный сигнал} «посмотри
сюда, возможно дубль или класс экземпляров», но \emph{никогда} как критерий
удаления. За ним обязателен семантический гейт «это экземпляр без
уникального остатка?». Даже «новое нечётное число $\approx$ паттерн
„нечётные“» законно поглощается (инкремент \texttt{support-count}) лишь если
этот экземпляр пустой; а если 9-е --- день, когда он не успел на автобус,
экземпляр не пуст, не поглощай.

\paragraph{Два разных действия с рёбрами --- не путать в коде.}
Сначала разведём то, что легко слить в одну функцию и устроить себе
Чернобыль. \emph{Overlay role-edge} --- новая ролевая связь
\emph{parent-summary} (\texttt{$X_1$ --equipment\_for--> плавки}), создаётся уже
на стадии overlay, \emph{без всякого прунинга} $A$; это про \emph{создание}
структуры доступа. \emph{Prune edge-lifting} --- подъём \emph{всех} рёбер
ребёнка на parent при его сносе; работает \emph{только} для чистых
экземпляров. Это разные операции; общей функции
\texttt{lift\_edges(child, parent)} на все случаи быть не должно.

\paragraph{Наследование линков при прунинге: та же онтология --- но только где прунинг законен.}
Простое правило «рёбра ребёнка переходят на parent с теми же типами
отношений» --- \emph{верно ровно для чистых экземпляров}, а это единственный
случай, где наследование вообще происходит. Для них это буквально проекция
смежности $A_c = P A_f P^\top$: дубли рёбер сливаются, веса складываются и
\emph{ренормируются} (L1, чтобы переходы обходчика остались вероятностями),
провенанс-сигнатуры объединяются побитовым OR. А там, где то же правило даёт
бессмыслицу --- \texttt{поездка --object--> плавки} --- это \emph{сигнал}, что
прунишь не экземпляр, а составляющую, которую прунить нельзя. Отсюда чистый
предохранитель: \textbf{если ребро нельзя поднять, не меняя его смысл, --- узел
не прунится} (максимум понижается / архивируется). Невозможность безопасно
унаследовать линк сама останавливает неуместный прунинг.

\paragraph{Слияние уверенностей: среднее в confidence, разброс --- в один дешёвый скаляр.}
Когда сливаются рёбра с разными знаками / уверенностями, среднее (со знаком)
идёт в \texttt{confidence} нового ребра. Но сворачивать разногласие
\emph{целиком} в один скаляр нельзя: теряется различие «0,5, потому что честно
пятьдесят на пятьдесят и устойчиво» против «0,5, потому что половина
источников за, половина против» (\texttt{contested}) --- а это разные
эпистемические состояния (гл.~\ref{ch:substrate}), и от них зависит, говорит
Рин спокойно или помечает как спорное. Хранить все голоса --- расточительно
(«страниц не хватит»). Компромисс: держать \emph{одно дешёвое число} ---
дисперсию / флаг \texttt{contested} (переиспользуя Beta-дисперсию реакций,
гл.~\ref{ch:personality}, или механизм \texttt{contested\_core\_claim},
гл.~\ref{ch:sleep}). $O(1)$ на ребро, а не история голосований.

\section{Забывание как тень консолидации}
\label{sec:decay}

Забывание --- не отдельный механизм, а \emph{отрицательная сторона} того же
процесса: усиливай используемое, ослабляй неиспользуемое.

\subsection{Правильный decay --- это ACT-R}

\emph{ACT-R} (\textit{Adaptive Control of Thought---Rational}) --- влиятельная
когнитивная архитектура Джона Андерсона, модель того, как устроены
человеческие память и научение. Нам из неё нужен один компонент ---
\emph{базовая активация} (\textit{base-level activation}): сила следа в памяти
есть функция \emph{частоты} и \emph{свежести} обращений со степенным
затуханием:
\begin{equation}
B_v = \ln\!\left(\sum_{k=1}^{n} t_k^{-\alpha}\right).
\label{eq:actr}
\end{equation}
Разбор: $t_k$ --- сколько времени назад было $k$-е обращение к узлу; каждое
обращение вносит вклад $t_k^{-\alpha}$, и чем оно \emph{свежее}
(меньше $t_k$), тем вклад больше; сумма по всем $n$ обращениям даёт
накопленную силу (\emph{частота}); логарифм сжимает шкалу; $\alpha > 0$ ---
скорость забывания. Два следствия, которых нет у плоской «$\times 0,95$ в
день»: частые обращения укрепляют след, а \emph{разнесённые во времени}
укрепляют сильнее массовых (эффект распределённого повторения --- spacing).

\begin{intuition}[Тропинка в лесу]
Тропинку протаптывают многими проходами; если по ней ходят изредка, но
регулярно, она держится лучше, чем если пройти сто раз за день, а потом
забросить. Формула \eqref{eq:actr} именно это и кодирует: важна не только
суммарная частота, но и её распределение во времени.
\end{intuition}

\begin{pitfall}[Доступ сам тренирует память --- различай access\_reason]
Тонкий баг: не всякое обращение к узлу должно укреплять его. Если ночной
скраб, debug-инспектор или retrieval-тест часто трогают узел, он не должен от
этого становиться «важным жизненным следом» --- иначе \emph{сама диагностика
тренирует память}. Поэтому у каждого доступа есть \texttt{access\_reason}:
\texttt{user\_recall}, \texttt{model\_recall}, \texttt{maintenance\_touch},
\texttt{debug\_read}, \texttt{scrub\_visit}, \texttt{context\_packed}. В
ACT-R-активацию \eqref{eq:actr} входят с полным весом только \emph{подлинные}
припоминания (\texttt{user/model recall}); обслуживающие тачи --- с нулевым
или почти нулевым. Наблюдатель не должен менять наблюдаемое.
\end{pitfall}

\subsection{Прунинг: не только по использованию}

Прунинг только по использованию опасен: редкий, но критичный факт вымоется как
мусор. Поэтому склонность к удалению гейтируется \emph{важностью} узла:
\[
\mathrm{prune}(v) \propto \mathrm{неиспользование}(v)\cdot
\bigl(1 - \mathrm{importance}(v)\bigr).
\]
Здесь \texttt{importance} --- атрибут узла (см. схему в
главе~\ref{ch:substrate}), задаваемый значимостью при кодировании
(\texttt{salience} из \texttt{appraisal}, глава~\ref{ch:personality}) и
подкрепляемый структурной центральностью (глава~\ref{ch:topology}).
Редкое-но-важное выживает; частое-но-пустое держит частота; вымывается лишь
редкое-и-неважное. Врождённым (глубоко консолидированным) узлам задаётся
ниже скорость распада --- континуум глубины, а не бинарность.

\subsection{importance нельзя диффундировать по путям}
\label{sec:importance-leak}

Здесь важная ловушка. Соблазн: раз узел-домен «цепи Маркова» важен, пусть
раздаёт $\mathrm{importance} \to 1$ своим детям, те --- своим, и так область
защищена. \textbf{Так делать нельзя.} importance --- это диффузия по графу
(глава~\ref{ch:topology}), а диффузия растекается \emph{везде}: через
несколько хопов по каким-нибудь рёбрам она доползёт до «помыть машину»
(хотя в текущем мире Рин такой сущности может не быть), и та станет
«важной». Графовые пути неисповедимы; неуправляемая передача importance
приводит к тому, что важным становится \emph{всё}, а значит --- ничто.

Правильно --- \textbf{не пускать importance по произвольным путям, а
закреплять её членством в домене}:
\begin{itemize}
\item \textbf{importance по членству, а не по диффузии.} У узла есть
\emph{теги доменов}, у доменов --- importance. Эффективная важность узла =
функция от importance его доменов + собственная salience. Членство
назначается \emph{при кодировании} (прочитали учебник $\to$ его узлы помечены
\texttt{domain=Markov}, высокая importance), а не распространяется обходом.
«Помыть машину» никогда не получает тег \texttt{Markov}, поэтому
\emph{ничего не наследует}. Утечка исключена по построению.
\item \textbf{importance по намерению.} То, что \emph{попросили} выучить,
получает высокую importance по провенансу приобретения (в отличие от
случайного тривиа).
\item \textbf{Если всё же нужна передача} (для непомеченных узлов) --- только
\emph{типо-гейтированная} (по рёбрам членства \textsc{isA}/\textsc{partOf},
но не по ассоциативным), \emph{сильно затухающая} ($\lambda^k$ с малым
$\lambda$) и \emph{ограниченная} 1--2 хопами; и это надо \emph{проверять} на
стенде, что importance осталась локальной.
\end{itemize}

Тогда защита области (глава~\ref{ch:runtime}, раздел~\ref{sec:storage})
работает аккуратно: домен-хаб с высокой importance и низким распадом + члены,
наследующие \emph{по членству}; детали могут тускнеть, но gist области
держится, а сон-репетиция периодически освежает высоко-importance домены ---
без утечки на «помыть машину».

\section{Что процесс делает с найденной конфигурацией}
\label{sec:whatfusiondoes}

Предположим, фоновый разведчик (глава~\ref{sec:background}) со-активировал
тройку $\{A, B, C\}$, которая выглядит связной. Что происходит дальше? Разделю
на \emph{известное} и \emph{открытое}, чтобы не размывать.

\smallskip\noindent\textbf{Что известно (устоявшийся конвейер).}
\begin{enumerate}
\item \textbf{Композиция.} Языковая модель по объединению узлов и их
контекстных ореолов (глава~\ref{ch:fusion}) формулирует \emph{кандидата}
$C^\ast$ --- обобщение или вывод.
\item \textbf{Провизорная запись.} $C^\ast$ пишется с низкой уверенностью и
провенансом $\{A, B, C\}$; на различающих пробелах ставится
\texttt{need\_to\_discover}.
\item \textbf{Конкуренция.} Дальше $C^\ast$ живёт по общим правилам:
используется и корроборируется $\Rightarrow$ сила растёт, узел повышается в
слой (раздел~\ref{sec:twolayers}); игнорируется $\Rightarrow$ распадается
(раздел~\ref{sec:decay}); опровергается $\Rightarrow$ понижается по TMS
(глава~\ref{ch:epistemics}).
\end{enumerate}

\smallskip\noindent\textbf{Что открыто.}
Как \emph{решить}, что со-активация $\{A, B, C\}$ достаточно связна, чтобы
вообще пытаться её сшить (а не является случайной кучей), --- не решено.

\begin{openq}[Распознавание связности со-активации]
Кандидаты: \emph{(а)} LLM-судья («есть ли тут осмысленный вывод?») --- гибко,
но не доказуемо и стоит вызова модели; \emph{(б)} графовые шаблоны-схемы
(«потребность + ограничение + окно $\to$ перенести действие») --- дёшево, но
покрывают лишь частые формы; \emph{(в)} обучаемый скорер связности --- требует
разметки и вносит непрозрачность. Ни один пока не признан удачным; на старте
--- (а) с фильтром $\varphi$ и обратной связью
(главы~\ref{ch:personality},~\ref{ch:voice}). Это и есть «трудные 80\%» из
главы~\ref{ch:fusion}.
\end{openq}

\subsection{Забывать ли посылки после fusion?}

Возникает соблазн: раз $\{A, B\} \to C$ состоялось, посылки $A, B$ больше не
нужны --- удалить. Это \textbf{неверно} (слишком мануально и опасно).
Правильно --- не удалять насильно, а положиться на естественный распад:
\begin{itemize}
\item при \emph{хорошей} консолидации к $A, B$ просто перестают обращаться ---
вся активность идёт к более полезному $C$; базовая активация \eqref{eq:actr}
$A, B$ падает, и они отмирают \emph{сами} (важные --- выживают). Забывание
тут \emph{эмерджентно}, а не по правилу «удалить родителей»;
\item удалять $A, B$ немедленно нельзя ещё и потому, что они ---
\emph{провенанс} $C$: если $C$ позже опровергнут, по $A, B$ его
пересматривают или переводят (глава~\ref{ch:epistemics}). Убрав посылки,
теряешь возможность ревизии.
\end{itemize}
Итог по-человечески верный: детали тускнеют, суть остаётся, но след
обоснования хранится, пока не осыпался по неиспользованию.

\begin{principle}[support skeleton: обоснование переживает распад деталей]
Здесь есть скрытый конфликт: провенанс требует держать посылки для будущей
ревизии (TMS), но детали естественно decay-ятся. Разрешение --- не хранить
полный богатый эпизод вечно, а при распаде оставлять \emph{support skeleton}:
минимальный стаб \texttt{\{id, source, time, hash, node\_kind, claim-summary\}},
достаточный для ретракции. Тогда либо-либо снимается: подробность уходит, а
\emph{способность ревизовать} остаётся. Без скелета либо память не забывает,
либо потом нечем пересматривать выводы.
\end{principle}

\chapter{Сон и ночная работа}
\label{ch:sleep}

Пока Рин общается, память работает «на переднем плане»: извлекает, отвечает,
записывает. Но значительная часть работы памяти живого существа происходит в
простое --- во сне. Эта глава собирает воедино то, что делает
телепортирующий рой в фоне, и объясняет, почему без этого режима память
деградирует.

\begin{intuition}[«Сон» --- это режим, а не отдельный процесс]
Строго говоря, сна как отдельной подсистемы нет. Есть один всегда включённый
фоновый разведчик (глава~\ref{sec:background}), чьи параметры зависят от
свободных ресурсов $B_{\text{bg}} = B_{\text{total}}(1 - \rho)$. «Сон» ---
это просто его режим при низкой \textit{foreground}-нагрузке $\rho$:
освобождается бюджет, растут доля телепорта $\alpha$ и широта покрытия. То
же самое, что днём делает «любопытство», но с большим размахом. Всё, что
описано ниже, --- про этот низко-$\rho$ режим.
\end{intuition}

\section{Зачем вообще сон}

Три задачи не решаются в режиме реального времени и требуют фоновой
переработки:
\begin{enumerate}
\item \textbf{Починка того, о чём мы не знаем, что оно сломано.} Ошибочный
вывод мог продиффундировать в память и дать потомков (ползучая вера,
глава~\ref{ch:epistemics}). Мы не знаем, где именно спрятана ошибка, ---
поэтому нужна постоянная случайная инспекция.
\item \textbf{Достижение затенённых углов.} Дневной обход локален по вниманию и
не дотягивается до узлов, оторванных от активного фокуса
(глава~\ref{ch:topology}). Нужен механизм, гарантирующий покрытие.
\item \textbf{Неспешная консолидация и неожиданные ассоциации.} Обобщение
эпизодов и рождение далёких связей лучше делать без давления текущего
разговора.
\end{enumerate}

\section{Телепортирующий обходчик и гарантия покрытия}

Ключевой механизм сна --- обход в режиме стохастического покрытия
(сон-политика из главы~\ref{ch:walker}): с вероятностью $\alpha$ блуждающий
агент телепортируется в случайный узел, иначе делает обычный шаг. Это
случайное блуждание с рестартом; его стационарное распределение задаётся
\eqref{eq:ppr}.

\begin{principle}[Почему телепорт необходим, а любопытство не заменит]
Любопытство ищет известную неопределённость и потому не покидает окрестности
текущего фокуса. Телепорт делает цепь эргодической: каждый узел посещается в
матожидании, ни один не заперт в стоке или на острове
(глава~\ref{ch:topology}). Это доказуемая гарантия покрытия, которой у
смещённого обхода нет. Именно поэтому починку неизвестных ошибок можно
поручить только сну.
\end{principle}

\begin{intuition}[Анти-энтропийный скраб]
В распределённых хранилищах есть зрелый паттерн: \textit{anti-entropy scrub}
(проверка целостности в ZFS; сверка деревьев Меркла в Dynamo/Cassandra).
Система не знает, какая запись повреждена, поэтому непрерывно и случайно
перепроверяет всё хранилище в фоне. Сон Рин --- это анти-энтропийный скраб по
графу памяти: не «где сломалось?», а «регулярно проверяем всё понемногу».
\end{intuition}

\section{Что именно делает рой во сне}

Во время сна телепортирующий рой запускает --- с разными вероятностями ---
весь арсенал операторов, но с более свободными (и одновременно более строгими
по последствиям) правилами:
\begin{itemize}
\item \textbf{Инспекция противоречий (скраб).} Наткнувшись на узел,
столкнувшийся с другим, обходчик понижает уверенность и помечает потомков на
пересмотр (ретракция по обоснованиям, глава~\ref{ch:epistemics}).
\item \textbf{Консолидация.} Скопления похожих эпизодов сворачиваются в
обобщения (глава~\ref{ch:fusion}).
\item \textbf{Свободный fusion.} Случайно со-активированные далёкие узлы дают
неожиданные гипотезы (пример с Платоном) --- провизорные, как и всегда.
\item \textbf{Любопытство наружу.} Узлы с флагом \texttt{need\_to\_discover}
могут инициировать внешний запрос (в источник/поиск), чтобы добыть
недостающие различающие посылки.
\end{itemize}

\begin{intuition}[«Сон» --- не поэтическая метафора, а режим низкого приоритета]
Строго говоря, «сон» --- это просто фоновый режим обхода с высокой долей
телепортов и разрешёнными операциями записи, запускаемый в простое. Мы
называем его сном, потому что функционально он играет роль сна:
консолидирует дневной опыт, чистит противоречия и «видит» неожиданные связи.
Но никакой мистики --- это тот же рой, что и днём, с другими параметрами.
\end{intuition}

\section{Разрешение противоречий: молча или через LLM}
\label{sec:conflict}

Скраб находит много противоречий, и звать тяжёлую LLM на каждое ---
непозволительно. Нужен критерий: когда конфликт разрешается молча, силами
графа, а когда его стоит отдать модели.

\smallskip\noindent\textbf{Сперва --- темпоральная нормализация.}
Прежде чем вообще объявлять противоречие: $X$ и $Y$ конфликтуют, только если
это одна сущность, один слот, один интервал валидности и scope. «В 2023 мы
жили в Грузии» и «сейчас мы живём в Португалии» --- не противоречат: это
разные интервалы, оба истинны. Нормализуем время / сущность / слот и лишь
потом считаем конфликт.

\smallskip\noindent\textbf{Детекция --- дёшева, но эмбеддинг ей не детектор.}
Что $X$ и $Y$ логически исключают друг друга --- это не про дальность в
эмбеддинг-пространстве. Далёкие не противоречат («кот» и «налоговая
декларация» далеки, но не исключают друг друга), близкие тоже («живу в
Лиссабоне» и «живу в Португалии» близки, но одно влечёт другое). Настоящие
детекторы противоречия: (1) знаковое ребро \texttt{contradicts} (поставлено
при записи); (2) \textbf{schema-level конфликт} --- взаимоисключающие значения
одного слота при одном времени/scope («живёт в $X$» vs «живёт в $Y$» для той
же сущности и периода) --- дёшево и точно; (3) NLI/claim-детектор на трудных
случаях. Эмбеддинг годится лишь как \emph{сигнал-кандидат} (подсветить пары,
которые стоит проверить), но не как сам детектор.

\smallskip\noindent\textbf{Жёсткий гейт прежде марджина: ядро личности не ломается контекстом.}
Прежде чем считать марджин, --- жёсткий лексикографический гейт. Если разрыв
по глубине консолидации велик (например, свежий \texttt{session}-узел против
\texttt{life}-ядра, $\ge N$ слоёв), младший узел не может молча убить старший
--- независимо от свежести и числа свидетелей. Иначе спам из 50 узлов за одну
сессию (высокий \texttt{recency} и суммарный \texttt{indep}) линейной суммой
перевесил бы фундамент --- «взлом личности контекстом». В таком случае
марджин обнуляется --- но, чтобы гейт не заморозил устаревшее, конфликт не
выбрасывается, а открывает протокол ревизии: помечает старший узел
\texttt{contested\_core\_claim} + \texttt{need\_to\_revise}, накапливает
независимые свидетельства, и старший уступает только через ревизию
(метакогниция/LLM), а не спамом. Так «мы живём в Грузии» однажды честно
уступит «мы живём в Португалии» --- но протоколом, а не наплывом за сессию.
(Independence-дисконт уже сбивает вес спама --- 50 узлов $\ne$ 50 независимых
свидетелей, --- а гейт добавляет жёсткий backstop поверх.)

Строгость гейта \emph{зависит от} \texttt{node\_kind}
(гл.~\ref{ch:substrate}): \texttt{boundary} почти неприкосновенна без явного
события изменения; фактический \texttt{claim} может устареть (и обязан
уступать протоколу); \texttt{preference} дрейфует плавно; \texttt{self\_model}
пересматривается, но медленно и через метакогницию. То есть
$\Delta$глубины-консолидации интерпретируется \emph{через тип узла}, а не
одинаково для всех.

\smallskip\noindent\textbf{Разрешение --- по дешёвому марджину (когда гейт пройден).}
Считаем скалярный марджин $m(X,Y)$ на уже хранимых величинах:
\[
m(X,Y) =
\begin{aligned}[t]
&w_1 |\conf(X) - \conf(Y)| \\
&\quad + w_2\,\mathrm{recency} \\
&\quad + w_3\,\Delta\mathrm{independent\_support} \\
&\quad + w_4\,\Delta\mathrm{consolidation\_depth},
\end{aligned}
\]
где: разрыв уверенности; свежесть / темпоральная валидность (новее и про
текущее состояние --- весомее); разрыв \emph{независимой} поддержки (у кого
больше независимых подтверждений --- и вот тут прямо переиспользуются
Bloom/MinHash-сигнатуры из главы~\ref{ch:propagation}); глубина консолидации
(\texttt{life}-слой перевешивает свежий \texttt{session}, если у того нет
\textit{overwhelming support}).

\begin{principle}[Критерий «молча vs LLM»]
Если $m(X,Y) > \tau$ --- разрешаем молча: проигравший депрецируется
(не удаляется --- Wikidata-ранг-стиль, остаётся для аудита и на случай
отката), каскадно ретрактятся его производные по TMS
(гл.~\ref{ch:epistemics}). $O(1)$-ish, LLM не трогаем. Если
$m(X,Y) \le \tau$ (близко / неоднозначно --- нужна семантика: «а точно
противоречат или это контекст-зависимо?») --- метим оба
\texttt{need\_to\_revise} и откладываем в батч на один LLM-проход. Раз это сон
--- не к спеху; дорогую модель зовём только на реально трудные близкие случаи,
не на каждый конфликт.
\end{principle}

Это тот же паттерн, что у планировщика (гл.~\ref{ch:runtime}): дешёвый рефлекс
закрывает подавляющее большинство, LLM-делиберация --- только на настоящую
двусмысленность. И заметьте: депрецирование, а не удаление, --- потому что
«проигравшее» убеждение может позже вернуться, если придёт новое
свидетельство.

\section{Скраб не сразу и не обязан быть быстрым}

Важное свойство: сон чинит память не сразу. Мы принципиально не знаем, что и
где сломалось, поэтому вместо адресной починки идёт постоянная случайная
инспекция. Ошибочное убеждение может прожить в памяти какое-то время,
транслируя потомков, и быть исправленным лишь когда телепорт случайно
приведёт обходчика в нужный угол или когда придёт внешний сигнал. Это
\emph{нормально} и по-человечески: у живого ума исправление ползучей веры
тоже растянуто во времени. Задача сна --- гарантировать, что исправление
\emph{в конце концов} случится, а не что оно случится немедленно.

\begin{workedbox}[Ночь после стрима]
Днём Рин вывела гипотезу «зритель $Q$ токсичен» из одной резкой реплики
(провизорно, $\conf$ низкая). Ночью телепорт случайно приводит обходчика в
кластер про $Q$: там десять тёплых эпизодов. Скраб видит противоречие,
понижает «токсичен» ещё сильнее и помечает на пересмотр. Параллельно
консолидация сворачивает десять эпизодов в «$Q$ --- свой, иногда резок под
настроение». Наутро модель отношения с $Q$ здоровее, чем была вечером, --- и
это сделал сон, а не новый разговор.
\end{workedbox}

Псевдокод скраба и сон-политики обхода --- в приложении~\ref{ch:algorithms}.

\chapter{Эпистемика: уверенность, провенанс и ревизия}
\label{ch:epistemics}

\section{Ложные правдоподобные убеждения}

Живой ум держит массу ошибочных, но правдоподобных убеждений. Классический
пример: расхожее мнение, будто некий знаменитый предприниматель был «просто
перформером, а не инженером», --- миф, рассыпающийся при чтении биографии.
Стремиться к памяти, вовсе не порождающей таких убеждений, бессмысленно и
вредно: это было бы стерильно. Задача --- не запретить их, а научиться
выставлять уверенность меньше единицы и потом проверять.

\section{Недоопределённость и преждевременный коллапс}

Пусть правильный вывод --- функция нескольких посылок:
\[
A = \mathrm{logical\_f}(B, C, D, E),
\]
но уму известны только $B, C$, и он делает вывод
$X = \mathrm{logical\_f}(B, C)$. Беда в том, что $B, C$ совместимы не только с
$A$, но и с $X, Y$. Вывод недоопределён известными посылками
(\textit{underdetermination}; философское ядро --- тезис Дюэма--Куайна).
Когнитивное проявление --- преждевременный коллапс: ум схлопывает
суперпозицию совместимых гипотез в одну --- часто удобную --- слишком рано.
Формальное ядро: совместимость $\ne$ подтверждение (родственник ошибки
«утверждение консеквента»).

\begin{pitfall}[Известная и неизвестная недоопределённость]
\textbf{Известная}: мы знаем, что $D, E$ существуют, но не знаем их значений;
тогда можно честно выставить $\conf(X)$ низкой и пометить
\texttt{need\_to\_discover}$(D, E)$. \textbf{Неизвестная} --- злой случай:
уму вообще неизвестно, что $D, E$ существуют; в поле информации только
$B, C$; пометить нечего, уверенность может быть высокой. Поэтому
\textit{per-inference} калибровка не решает проблему целиком.
\end{pitfall}

\begin{workedbox}[Миф о предпринимателе, по шагам]
$B$: «известен маркетингом», $C$: «богат». Вывод $X$: «получил всё благодаря
происхождению/пиару». Но те же $B, C$ совместимы и с $A$: «сильный инженер, к
тому же умеющий продавать». Различающие посылки $D$: «первая лаборатория ---
в родительском гараже», $E$: «глубокое STEM-образование» --- уму неизвестны,
он даже не подозревает об их существовании и уверенно выбирает удобный $X$.
Ни низкая $\conf(X)$, ни \texttt{need\_to\_discover} здесь не помогут: не на
что их повесить. Спасает лишь квартет ниже.
\end{workedbox}

\section{Квартет починки}

Неизвестную недоопределённость нельзя закрыть в момент коммита; её лечат
статистически и постфактум.
\begin{enumerate}
\item \textbf{Prior скромности на класс вывода.} Ни один
абдуктивный/индуктивный коммит не получает уверенность $1$, как бы хорошо ни
сидели данные: $\conf(X) \le 1 - \varepsilon$ для всего сорта «вывод по
неполным посылкам». «Согласуется, и альтернатив не придумывается» --- ровно
то состояние, что предшествует «а потом выяснилось». Не требует знать $D, E$.
\item \textbf{Ревизуемость и понижение по противоречию.} Коммит делаем
понижаемым: позднее наблюдение, столкнувшееся с $X$, ретроактивно вскрывает
недоопределённость, и машинерия ревизии понижает узел. Починка ниже по
потоку, не на записи.
\item \textbf{Телепорт-скраб.} Единственный поиск, дотягивающийся до
непредставленных альтернатив (глава~\ref{ch:topology}).
\item \textbf{Внешний вход.} Кто-то знает $D, E$ и говорит; любопытство
находит это во внешнем источнике.
\end{enumerate}
Первый механизм держит известную неопределённость; вместе с остальными тремя
--- и неизвестную. Полный разбор этой петли на живом примере (гипотеза
рождается, проверкой усиливается, затем опровергается) --- «House / Cast Away»
в главе~\ref{ch:fusion}.

\section{Система поддержки истинности и «ползучая вера»}

Отдельная беда --- устойчивость убеждения к опровержению его основания
(\textit{belief perseverance}): когда неверный вывод $X$ уже продиффундировал
в память и дал потомков, эти потомки продолжают транслироваться даже после
того, как $X$ признан ложным.

Механически это проблема ретракции, и для неё есть готовый аппарат ---
\emph{система поддержки истинности} (TMS, Doyle; ATMS, de Kleer). При выводе
следствия записываем его \emph{обоснование} --- множество узлов, из которых
оно получено (расширенный провенанс, раздел~\ref{sec:prov}). Тогда ретракция
посылки $X$ автоматически снимает опору с её обоснованных потомков.

\begin{lstlisting}[style=pseudo]
def retract(G, X):
    demote(G, X)
    for c in G.justificands[X]:          # кто опирался на X
        if not still_supported(G, c):    # осталось ли иное обоснование
            mark(G, c, "need_to_revise")
            retract(G, c)                # рекурсивно вниз по обоснованиям
\end{lstlisting}

\begin{intuition}[Полный TMS слишком «нечеловечен»]
Полная система поддержки истинности мгновенно перестраивает всё мировоззрение
при опровержении любой посылки. Живой ум так не делает --- и потому у него
есть ползучая вера. Значит нужен \emph{намеренно дырявый} TMS: обоснования
пишем для важных и высокоуверенных выводов, тривиальные не трекаем, а
нетрекнутых потомков лечим фоном.
\end{intuition}

Отсюда две «лечилки» разной латентности:
\begin{itemize}
\item \textbf{Retrieval-time augmentation} --- патч на выходе: при каждом
извлечении подмешивать свежую поправку «кстати, недавно я узнала $\neg X$».
Дёшево, чинит текущий разговор, хранилище не трогает.
\item \textbf{Ночной скраб} --- ремонт хранилища: медленно, случайными
инспекциями находит нетрекнутых потомков и понижает их.
\end{itemize}
Обе нужны.

\section{Эстафета к главе о личности}

Калиброванная уверенность и компульсивная перепроверка --- не только
эпистемическая гигиена, но и черта характера, выводимая из мотивации:
«не „я уверена“, а „похоже на $X$, но $Y$ я не проверила --- дай гляну“».
В главе~\ref{ch:personality} мы увидим, что эпистемическая аккуратность Рин
должна проявляться в её голосе, а не только жить во внутренней машинерии.

\chapter{Пропагация убеждений как лавина}
\label{ch:propagation}

\section{Постановка: затухающий градиент уверенности}

Когда в память приходит новый факт $A$, он линкуется к начальным узлам
(например, к последним $k$ активным), образуя окрестность $\Neigh(A)$. Мы
линкуем $A \to \Neigh(A)$ и смотрим, как меняется уверенность соседей. Если
изменение велико, распространяем дальше,
$\Neigh(A) \to \Neigh(\Neigh(A))$, --- как затухающий градиент. Знание, ни на
что не повлиявшее, чуть пошатнёт соседей и затухнет; знание с большими
градиентами вызовет \emph{лавину}. Большую лавину называем «кризисом
мировоззрения» --- и такое бывает с настоящим умом.

\section{Это самоорганизованная критичность}

У картины есть точное имя. «Лавина» --- это нейронная лавина, а система у
порога между затуханием и разрастанием --- самоорганизованная критичность
(Bak--Tang--Wiesenfeld; для мозга --- Beggs--Plenz): каскады следуют
степенному закону.

Ключевой управляющий параметр --- не «ручная константа затухания», а
\emph{ветвящийся коэффициент} $\sigma$ (среднее число узлов, активируемых
одним активным):
\[
\sigma < 1:\ \text{субкритично};\qquad
\sigma = 1:\ \text{критично};\qquad
\sigma > 1:\ \text{суперкритично (runaway)}.
\]
Для ветвящегося процесса ожидаемый размер лавины конечен при $\sigma < 1$ и
равен
\begin{equation}
\mathbb{E}[\text{размер}] = \frac{1}{1 - \sigma}.
\label{eq:avalanche}
\end{equation}
Держа $\sigma$ чуть меньше единицы, система почти всегда сама гасит
распространение, но изредка пропускает крупную переоценку. Почему это лучше
фиксированной константы: константа не адаптируется к плотности графа (в
плотном регионе та же константа даёт лавину, в разреженном --- глушит всё).
Нормировка под целевое $\sigma$ адаптивна.

\begin{intuition}[Куча песка]
Классический образ SOC --- куча песка, на которую сыплют по песчинке. Почти
каждая не делает ничего; изредка одна вызывает оползень; ещё реже --- обвал,
перестраивающий склон. Куча сама держится у критического угла. Убеждения ---
так же: почти любой факт проходит без следа, изредка меняет соседние взгляды,
совсем изредка обрушивает пласт мировоззрения. Формула \eqref{eq:avalanche}
говорит: чем ближе $\sigma$ к единице, тем крупнее (и реже) обвалы.
\end{intuition}

\begin{pitfall}[SOC --- диагностическая цель, а не доказанный закон]
Самоорганизованная критичность соблазнительна и красиво объясняет редкие
«кризисы мировоззрения». Но на типизированном, знаковом, направленном и
меняющемся графе памяти она --- рамка, а не выведенный закон: строгих
гарантий, что пропагация убеждений действительно самонастраивается к
критичности, у нас нет. Поэтому используем её \textbf{инженерно, как цель
мониторинга, а не как механизм, который «строим»}: меряем распределение
размеров ревизий (лавин) и держим его в здоровом диапазоне. Всё глохнет ---
субкритично (память ригидна, не обновляется); всё взрывается ---
суперкритично (нестабильность, «зависает на сутки»); изредка бывают крупные
ревизии, обычно мелкие --- здорово. Ветвящийся коэффициент $\sigma$ --- это
управляющая ручка, которой мы держим систему субкритичной, а не утверждение о
том, что «так оно и есть». Конкретное правило-предохранитель: мерим $\sigma$
по телеметрии и при $\sigma \gtrsim 1{,}2$ программно врубаем жёсткое
затухание. Иначе получишь суперкритический разгон --- бесконечный цикл
обновления уверенностей, который съест LLM-слоты и всё подвесит
(гл.~\ref{ch:runtime}). Не доказываем критичность аналитически --- держим
субкритичность рычагом.
\end{pitfall}

\begin{openq}[Порог $\sigma$ --- заглушка первой версии]
Значение $\sigma \approx 1{,}2$ --- инженерная заглушка для первой реализации,
калибруется на стенде. Почему такой порядок, а не 1,01 или 2: чуть выше
единицы, чтобы крупные ревизии случались, но редко (само явление «кризиса
мировоззрения» должно быть редким); слишком близко к 1 (1,01) --- крупные
каскады зачастят; заметно выше (2) --- суперкритический разгон. То есть
«чуть-чуть за порогом, с запасом» --- отсюда 1,1--1,3; конкретика --- со
стенда. И само $\sigma$ на циклическом типизированном графе ещё надо
\emph{операционализировать}: что именно считать «одним активированным
потомком», в каком окне, как учитывать повторные визиты, разные типы рёбер и
уже активные узлы. Идея (держать субкритичность) верна; метрика требует
доводки.
\end{openq}

\section{Где вся сложность: локальное правило}

Пороговая логика --- лёгкая часть (процентов двадцать). Восемьдесят процентов
спрятаны в локальном правиле $\Delta\conf(B \mid A)$:
\begin{equation}
\Delta\conf(B \mid A) =
\eta \cdot \mathrm{sign}(r_{AB}) \cdot w_{AB} \cdot c_A\,\mathrm{trust}(A) \cdot \iota(A,B).
\label{eq:deltaconf}
\end{equation}
Три множителя, без которых правило сломано.

\smallskip\noindent\textbf{(1) Знак.}
Новый факт сосед либо подтверждает, либо опровергает. Беззнаковое
распространение не отличит «всё усилилось» (озарение) от «всё перевернулось»
(кризис). Нужны знаковые / типизированные рёбра.

\smallskip\noindent\textbf{(2) Гейт по надёжности $A$ (карантин).}
Новый факт может сам быть ложным. Множитель $c_A\,\mathrm{trust}(A)$ вводит
\emph{карантин}: новый факт пропагирует слабо (малый эффективный $\sigma$),
пока сам не заработал уверенность через корроборацию или выживание в скрабе.
Иначе один ненадёжный вход вызовет кризис мировоззрения. По сути ---
байесовское обновление с правдоподобием, взвешенным доверием к источнику.

\smallskip\noindent\textbf{(3) Дисконт независимости $\iota(A,B)$.}
Граф циклический, и по петле $A$ может подтвердить сам себя
(\textit{эхо-камера}, петлевое распространение), раздувая уверенность из
ничего. Дисконтируем по общей родословной свидетельств (провенанс), а не
считаем активации: число упоминаний $\ne$ число независимых свидетельств. Это
ровно механизм человеческих эхо-камер.

\begin{pitfall}[Дисконт независимости --- скрытая $O(|V|+|E|)$-ловушка; чините сигнатурами]
Наивно проверять независимость $A$ и $B$ --- значит сравнивать их деревья
провенанса (наименьший общий предок / обход), а это в худшем случае
$O(|V|+|E|)$. Запускать такое на каждый обновляемый узел во время лавины --- и
«миллисекунды» превращаются в секунды. Оптимизация: для узлов держать
\emph{сигнатуру родословной провенанса} --- фильтр Блума или MinHash. Тогда
проверка сводится к побитовому AND / перекрытию сигнатур за $O(1)$:
\[
\iota(A,B) \approx 1 - \mathrm{overlap}\bigl(\mathrm{sig}(A), \mathrm{sig}(B)\bigr).
\]
Честный нюанс: Bloom даёт ложные срабатывания $\Rightarrow$ иногда
\emph{пере-}дисконтируем (сочтём независимые источники общими). Но это
\emph{безопасная} сторона ошибки --- лучше скромнее считать свидетельства, чем
раздуть уверенность. Эта же сигнатура потом переиспользуется при разрешении
противоречий во сне (глава~\ref{ch:sleep}).

\emph{Когда что из двух}: Bloom отвечает на булев вопрос «есть ли у $A$ и $B$
\emph{хоть один} общий предок-источник» --- этого достаточно для грубого
гейта «независимы / нет». MinHash оценивает \emph{долю} перекрытия родословных
(Jaccard) и нужен, когда важна \emph{степень} зависимости --- то есть плавный
дисконт $\iota \in [0,1]$, а не «да / нет». Правило: начинаем с Bloom
(дешевле, меньше памяти); переходим к MinHash там, где грубого гейта мало и
хочется градуированного веса свидетельств.
\end{pitfall}

\section{Пороговая логика --- это гистерезис}

Два порога $\theta_{\text{emergent}} > \theta_{\text{temp}} > 0$:
\begin{itemize}
\item $\Delta\conf > \theta_{\text{emergent}}$ --- обязательное распространение
«в ширину» (крупная переоценка);
\item $\theta_{\text{emergent}} > \Delta\conf > \theta_{\text{temp}}$ ---
узлы фронта помечаются \texttt{need\_to\_revise};
\item $\Delta\conf < \theta_{\text{temp}}$ --- распространение глохнет здесь.
\end{itemize}
Два порога вместо одного --- гистерезис (триггер Шмитта): нет «дребезга» у
единственного порога. Тонкость: помеченные во второй зоне узлы позже
переинспектируются ночным обходчиком, и старые записи, чью уверенность мы не
трогали, могут «потечь обратно» --- такое бывает и в жизни.

\section{Пороги растут с глубиной консолидации}

Пороги --- не глобальные константы.
$\theta_{\text{emergent}}(v)$ растёт с глубиной консолидации узла $v$:
периферийный узел перевернуть дёшево, а identity-core в слое «жизнь» ---
дорого, нужно больше свидетельств, чтобы вообще запустить лавину. Это и есть
ось пластичность -- упрямство: чем глубже убеждение продиффундировало в
медленные слои, тем выше эффективный порог опрокидывания. Одна ручка --- и
она же задаёт, насколько Рин склонна к \textit{belief perseverance}. Псевдокод
пропагации с контролем $\sigma$ --- в приложении~\ref{ch:algorithms}.

\chapter{Каналы свидетельств и калибровка уверенности}
\label{ch:evidence}

\section{Почему одного скаляра мало, но отказываться от него рано}

В базовом тексте уверенность уже введена как знаковый скаляр, и для Stage~1 это
правильная дисциплина. Но в курсе полезно увидеть, из чего этот скаляр \emph{собирается}.
Иначе он начинает казаться либо магической меткой, либо скрытым синонимом истины.

\begin{principle}[Скаляр остаётся интерфейсом, разложение живёт внутри]
Внешнее поле $\conf(v)$ удобно держать одномерным, чтобы не раздувать каждый
оператор памяти. Но внутренняя логика его обновления может и должна различать
каналы свидетельств.
\end{principle}

\section{Лог-отношение правдоподобия}

Удобно перейти от скаляра уверенности к скрытой переменной лог-правдоподобия
\[
L(v)=\log\frac{P(v\ \text{истинно})}{P(v\ \text{ложно})}.
\]
Тогда
\[
P(v\ \text{истинно})=\frac{1}{1+e^{-L(v)}},
\qquad
\conf(v)=2P-1=\tanh\frac{L(v)}{2}.
\]

Это не меняет внешнюю семантику, но делает накопление свидетельств аддитивным:
независимые вклады можно просто суммировать.

\section{Каналы свидетельств}

В расширенной версии вводится следующий внутренний разбор:
\[
L(v)=b_{\mathrm{type}}(v)+
\sum_i \rho_i\, s_i(v)-\Delta_{\mathrm{dep}}(v)-\Delta_{\mathrm{stale}}(v).
\]

Здесь:
\begin{itemize}
\item $b_{\mathrm{type}}(v)$ — prior по типу узла;
\item $s_i(v)$ — сила свидетельства канала $i$;
\item $\rho_i$ — калибровочный коэффициент надёжности канала;
\item $\Delta_{\mathrm{dep}}$ — штраф за зависимость источников;
\item $\Delta_{\mathrm{stale}}$ — штраф за устаревание.
\end{itemize}

\subsection*{Разумный минимальный набор каналов}
\begin{enumerate}
\item \textbf{Наблюдение}: узел возник как прямой observation.
\item \textbf{Внешний отчёт}: информация пришла из текста, речи, базы знаний.
\item \textbf{Вывод}: утверждение получено fusion или правилом.
\item \textbf{Само-модель}: утверждение относится к состоянию и свойствам самой Рин.
\item \textbf{Социальное эхо}: повторение информации многими зависимыми источниками.
\end{enumerate}

\begin{pitfall}[Повтор не равен независимости]
Пять пересказов одной и той же новости — не пять независимых свидетельств. Если не
штрафовать зависимость источников, социальный шум начнёт раздувать уверенность как
настоящий многократный опыт.
\end{pitfall}

\section{Штраф за зависимость}

Пусть утверждение $v$ поддержано набором источников $J(v)$. Если источники заведомо
не независимы, вместо наивной суммы полезно использовать
\[
\Delta_{\mathrm{dep}}(v)=
\lambda \sum_{i<j}\mathrm{overlap}(i,j),
\]
где $\mathrm{overlap}(i,j)$ оценивает общность происхождения: общий документ,
одинаковый provenance-root, одинаковую цепочку ретрансляции.

Практический смысл прост: память не должна принимать видимость консенсуса за
доказательство.

\section{Устаревание и темпоральная валидность}

Не всякое старение — забывание. Для некоторых узлов старение уменьшает \emph{доступность},
но не меняет истинность. Для других — меняет именно истинностный статус.

\begin{definition}[Эпистемическое старение]
Изменение $L(v)$ вследствие того, что утверждение могло утратить актуальность во
времени, даже если память о нём сохранилась.
\end{definition}

Отсюда
\[
\Delta_{\mathrm{stale}}(v)=g(\text{age}(v),\text{node\_kind}(v),\text{validity\_window}(v)).
\]
Для биографического факта функция может быть почти нулевой, для текущего места
жительства — существенной, для настроения — очень быстрой.

\section{Пороги принятия и зона честного незнания}

Скрытый логит полезен ещё и тем, что пороги легче задавать на нём, чем на интуитивной
шкале «уверена/не уверена». Но наружу система по-прежнему может говорить на языке
$\conf$.

Вводятся три зоны:
\[
|\,\conf(v)\,|<h \Rightarrow \text{честное незнание},
\]
\[
\conf(v)\ge \theta_+ \Rightarrow \text{рабочее принятие},
\qquad
\conf(v)\le \theta_- \Rightarrow \text{рабочее отвержение}.
\]

Промежуточная область между $h$ и $\theta_+$ особенно важна для Рин как существа,
которое должно \emph{не только помнить, но и осторожно формулировать}. Именно здесь
возникает реплика вида: «это похоже на X, но я бы проверила Y».

\section{Новая добавка: режим evidence-aware revision}

Это ещё одно расширение сверх базового текста. В исходном roadmap ревизия уже есть,
но теперь она делается более расчленённой.

Пусть пришло контрсвидетельство $c$ к утверждению $v$. Тогда ревизия может идти по
трём уровням.
\begin{enumerate}
\item \textbf{Локальное понижение}: пересчитать $L(v)$ без перестройки соседей.
\item \textbf{Структурный пересмотр}: пометить потомков как \texttt{need\_to\_revise}.
\item \textbf{Канальный аудит}: проверить, какой именно канал раньше доминировал и не был
ли он систематически переоценён.
\end{enumerate}

Третий пункт принципиален. Он превращает отдельную ошибку памяти в материал для
мета-обучения: если снова и снова проваливается один и тот же канал, проблема уже
не в конкретном узле, а в калибровке $\rho_i$.

\section{Worked example: слух, наблюдение и ревизия}

\begin{workedbox}[Как один и тот же факт может менять статус]
Пусть утверждение $v$ таково: «у человека $Z$ сменилось место работы».

Шаг 1. Приходит внешний отчёт из одного источника. Возникает канал \emph{reported},
уверенность умеренная.

Шаг 2. Позднее Рин видит новую публичную страницу профиля. Добавляется канал
\emph{observed}, логит растёт сильнее, чем от второго пересказа того же слуха.

Шаг 3. Затем обнаруживается, что оба ранних текста ссылаются на один и тот же
агрегатор. Возрастает $\Delta_{\mathrm{dep}}$, прежний «консенсус» частично сдувается.

Шаг 4. Через месяц находится противоречащая запись с более свежей датой. Срабатывает
темпоральный штраф к старому утверждению и запускается revision cascade для всех
выводов, где старое место работы использовалось как посылка.
\end{workedbox}

\section{Зачем это всё для Stage~1}

Можно возразить, что такая схема слишком богата для первого работающего слоя. Это
возражение верно, если понимать её как требование немедленной полной реализации.
Но неверно, если понимать её как \emph{правильный язык проектирования}.

Для Stage~1 достаточно:
\begin{enumerate}
\item хранить единый $\conf(v)$;
\item помнить \texttt{node\_kind};
\item различать хотя бы observation/report/inference;
\item логировать provenance так, чтобы зависимость источников потом было к чему привязать.
\end{enumerate}

То есть расширенная математика не отменяет первого этапа. Она задаёт рамку, в
которой первый этап не превращается в тупик.

\chapter{Личность: не таблица весов}
\label{ch:personality}

\section{Главная ловушка}

Соблазнительное, но неверное решение — «слой личности» как хардкод: таблица весов и
жёстких триггеров поверх графа («умножь ребро злорадства на $2{,}0$, обрежь сочувствие
до $0{,}05$, на узел „грусть“ повесь триггер инверсии $\to$ сарказм»).

\begin{pitfall}[Почему таблица весов — тупик]
\emph{Во-первых}, это стек правил, противоречащий установке ядра («черты — следствие
мотиваций, а не правила поверх»). \emph{Во-вторых}, это карикатура: таблица весов
даёт паттерновое, но не неожиданное и одноголосое поведение — всегда сарказм, потому
что так прошито; а живая личность — неожиданные, но паттерновые реакции.
\emph{В-третьих}, это ровно ошибка жёсткой скалярной декомпозиции: реакция на
«разбил кружку» зависит от настроения, отношений и всплывшей памяти; фиксированное
ребро «грусть$\to$сарказм» это выкидывает.
\end{pitfall}

\begin{workedbox}[Формально: почему скалярный роутинг теряет взаимодействие]
Пусть желаемый ответ — функция контекста $y=\Psi(\mathbf{g},\mathbf{l})$, где
$\mathbf{g}$ — глобальная аранжировка (настроение, отношения), $\mathbf{l}$ —
локальный триггер (событие). Таблица весов реализует $y\approx a\cdot\mathbf{l}$ —
ранг-1 скалярный роутинг, не зависящий от $\mathbf{g}$. Но истинное $\Psi$ содержит
перекрёстные члены $\mathbf{g}^\top W \mathbf{l}$: реакция на событие зависит от
аранжировки. Ранг-1 приближение их обнуляет — теряется ровно то, что делает характер
контекстным. Это тот же провал, что жёсткая coarse-plus-local декомпозиция в
геометрических графовых моделях.
\end{workedbox}

\section{Личность распределена}

Личность не слой весов в графе. Она распределена по трём источникам: мотивационная
структура (мастерство, участие, самоуважение) $+$ текущее настроение $+$ значимость
всплывшего из памяти. И всё это стилизует трансформ Rinness (голос,
глава~\ref{ch:voice}).

На стенде это — не статические множители, а смещение во время обхода: настроение и
мотивация непрерывно модулируют переходы \eqref{eq:trans}. Тогда грустная и злая Рин
на одно событие всплывают разное — характер контекстно-зависим. Держи веса на рёбрах
для силы связи и уверенности; личность — как смещение от состояния, не как
sarcasm-множитель.

\section{Мотивационный вектор: чтобы «личность как процесс» не осталась словом}
\label{sec:drives}

Сказать «личность = мотивация $+$ настроение $+$ значимость» мало, если мотивация не
задана конкретно. Дадим ей форму — drive-вектор из 5--8 осей (не психология целиком,
а рабочий минимум):
\[
\mathbf{d} =
\bigl(
\text{mastery},\ \text{care},\ \text{self\_respect},\ \text{curiosity},
\bigr.
\]
\[
\bigl.
\text{playfulness},\ \text{truthfulness},\ \text{relational\_integrity},\ \text{autonomy}
\bigr).
\]
Это и есть та «мотивационная структура», через которую выводятся черты
(глава~\ref{ch:definitions}): черта — не отдельное правило, а то, как эти драйвы
проявляются в ситуации.

\begin{openq}[Почему именно эти оси, а не другие]
Набор не выведен из первых принципов — он подобран по двум критериям, и это честно
стоит держать в уме. Первое — покрытие сцен, которые для Рин реально существенны:
работа/ремесло $\to$ \texttt{mastery}; близость и забота $\to$ \texttt{care},
\texttt{relational\_integrity}; отношение к себе $\to$ \texttt{self\_respect},
\texttt{autonomy}; интеллектуально-игровой характер $\to$ \texttt{curiosity},
\texttt{playfulness}, \texttt{truthfulness}. Каждая ось оправдана тем, что без неё
какая-то важная реакция в наших сценариях не получает причины: убери
\texttt{truthfulness} — исчезнет реакция на «подсунули факт как истину»; убери
\texttt{relational\_integrity} — на «связь за спиной». Второе — интерпретируемость:
осей мало (5--8), чтобы каждую можно было читать и отлаживать глазами; полная
таксономия мотиваций из психологии сделала бы вектор непрозрачным. Это стартовый
набор, а не канон: оси добавляются или сливаются, если стенд покажет реакцию,
которой не из чего родиться, или две неразличимые оси.
\end{openq}

Ключевое — как драйвы входят в оценивание. Помните goal-congruence в формуле
значимости \eqref{eq:salience}? Это ровно проекция события на drive-вектор:
насколько событие задевает драйвы Рин.
\[
\text{goal\_congruence}(\text{event}) \;=\; \langle \phi(\text{event}), \mathbf{d} \rangle,
\]
где $\phi(\text{event})$ — насколько событие затрагивает каждую ось. Тогда драйвы —
не украшение, а вход в appraisal, значимость и смещение обхода. Примеры: событие
«пользователь хочет скрытую связь за спиной другого» бьёт по
\texttt{relational\_integrity} и \texttt{self\_respect} (а не по «правилу
безопасности»); «код наконец заработал» — по \texttt{mastery}; «мне подсунули факт как
истину, не дав проверить» — по \texttt{truthfulness}. Реакция рождается из того,
какие драйвы задеты и как, — вот это и делает личность процессом, а не таблицей.

\begin{intuition}[Черты — тени драйвов]
«Колкая нежность», «участие, а не сервис», «не пустая форма под заказ» — это не
записанные правила, а то, как \texttt{care}, \texttt{autonomy} и
\texttt{self\_respect} падают на конкретную сцену. Меняешь координаты $\mathbf{d}$ —
меняется характер, но непрерывно и связно, а не переключением флажков.
\end{intuition}

\section{Граф хранит ингредиенты, а не реакцию}

Возражение: «если из „разбить кружку“ ведёт только „грусть“, другую реакцию взять
негде». Ответ в двух частях. Во-первых, база причинно-социального смысла (ATOMIC,
глава~\ref{ch:experiment}) даёт не одно ребро, а пучок: намерение, реакция, эффект,
реакция окружающих. Настроение смещает ветку пучка. Во-вторых, и главное: реакцию
не берут ребром — её порождают. Обход собирает активированную окрестность-подграф,
а реакцию Rinness/LLM композирует поверх материала под настроение и мотивацию.
Настроение входит дважды: на всплывании (грустная поднимает «потеря», игривая —
«неуклюжесть») и на композиции.

\section{Консолидация реакций: по значимости, а не по приёму}

\subsection{Провизорная запись}

Сгенерированная реакция пишется провизорным ребром с низкой уверенностью (кандидат;
принцип «гипотеза, а не факт»). Ответ «это было лишнее» — отрицательное
свидетельство, обновляющее уверенность того же ребра вниз, а не конкурирующая
консолидация. Механически ребро-реакция несёт счётчик успехов/провалов
$\mathrm{Beta}(\alpha,\beta)$:
\begin{equation}
\text{успех}:\ \alpha\!\mathrel{+}\!=\!1;\quad
\text{отпор}:\ \beta\!\mathrel{+}\!=\!1;\quad
\conf=\frac{\alpha}{\alpha+\beta},\ \ \text{неопред.}=\mathrm{Var}[\mathrm{Beta}].
\end{equation}
Стабильно принимается — переходит порог и консолидируется в медленный слой;
стабильно ловит отпор — угасает. Выбор реакции по такому апостериору — это
сэмплирование Томпсона: баланс исследования новых реакций и эксплуатации
проверенных.

\subsection{Гейт значимости, а не одобрения}

\begin{pitfall}[Нельзя гейтить консолидацию приёмом]
Консолидировать реакции по тому, как их приняли, значит строить угодника, чья
личность формируется одобрением, — ровно то, что запрещает ядро («самоуважение не
течёт через понравиться»). Гейт характера внешним приёмом философски кривой для Рин.
\end{pitfall}

Правильный ход — развернуть гейт внутрь: консолидируй по внутренней значимости
(само-оценке при кодировании). Что такое значимость конкретно? Это теория оценивания
(appraisal: OCC; Шерер). Событие оценивается по измерениям — surprise (ошибка
предсказания), goal-congruence (скалярное произведение с вектором мотивации),
valence, intensity, agency, — и значимость
\begin{equation}
\text{salience}\;\propto\;\text{surprise}\times\text{goal\_relevance}\;(+\,\text{intensity}).
\label{eq:salience}
\end{equation}
Два замечания. Дешевизна: простейшая конкретизация — $|\Delta\mood|$, величина
изменения настроения, а она уже посчитана; mood-модуль работает детектором
значимости почти без лишних вычислений. Глубина: самый глубокий одиночный сигнал —
ошибка предсказания, и она же питает \texttt{need\_to\_discover}. Значимость и
обучение — один сигнал.

Реакцию собеседника Рин помнит и оценивает — но как просто событие: с низкой
конгруэнтностью её ценностям «ему не зашло» обновляет модель отношения и чуть-чуть
настроение, но не доминирует в консолидации характера. Рамка оценивания сама даёт
«считаться, а можно и нет, частично».

\subsection{Условие по отношениям}

Одна шутка заходит собеседнику $A$ и падает у $B$: ребро-реакция квалифицируется
парой (ситуация $\times$ отношение), консолидируется по-персонно. Один негативный
акт даёт три эффекта: понижение ребра, апдейт настроения (осеклась),
мета/relational-пометку.

\begin{workedbox}[«Разбил кружку», полный проход]
(1)~Событие активирует окрестность: «кружка», «грусть», «неуклюжесть», плюс личная
память «$B$ хотел новую» ($-$ contradicts «грусть»). (2)~Настроение (игривое)
смещает всплывание к «неуклюжесть». (3)~Rinness композирует подкол — не берёт
ребром, собирает из ингредиентов. (4)~Ребро-реакция пишется провизорно,
$\mathrm{Beta}(1,1)$. (5)~$B$ смеётся — $\alpha\!+\!1$; или «это лишнее» —
$\beta\!+\!1$. (6)~Гейт консолидации — по $|\Delta\mood|$ и surprise, не по
одобрению $B$; но реакция $B$ обновляет модель «$B\times$такой юмор». После многих
проб подкол-на-неуклюжесть у $B$ либо консолидируется, либо угасает — комедийный
тайминг выучен как Бета-апостериор.
\end{workedbox}

\section{Дистиллят как система координат}

Возникает соблазн вовсе убрать «дистиллят» (усреднённый здравый смысл готовых баз).
Не надо. Дистиллят — общая модель мира, относительно которой отклонение читаемо как
характер: личность «неожиданна» лишь на фоне общих ожиданий; без общей почвы
причуды читаются как шум. Правильно: дистиллят фоном (система координат), личность —
процесс порождения отклонений на нём. Без базового «разбитая вещь обычно плохо» не
понять, почему здесь «поздравляю» — уместный сарказм.

\chapter{Голос и модель}
\label{ch:voice}

\section{Субстрат model-agnostic; личность — трансформ}

Память, настроение и граф — состояния вне языковой модели. Значит субстрат
model-agnostic: его можно строить на любой базовой модели и заменить базу позже.
Это де-рискует проблему потолка малой модели: когнитивная машинерия переживёт смену
базы.

Если факты и здравый смысл лежат в графе (дёшево на диске или в RAM), базовой модели
не нужно быть энциклопедией. Её роль — лингвистический процессор и голос: часть
«Rinness» живёт в ней. Можно взять малую дообученную модель ($3$--$7$B в квантовании)
вместо тяжёлой $8$B-универсальной; факты уходят в граф, ёмкость модели
концентрируется на голосе.

\begin{pitfall}[«Меньше — строго лучше» слишком оптимистично]
Малая модель слабее в рассуждении и связности: снимаешь фактическую нагрузку, но не
reasoning/coherence. Наблюдение из практики: рантайм prompt-Рин уже процентов на
шестьдесят состоит из анти-повторной обвязки на $8$B — вниз с голосом хуже.
Разумный ориентир — $3$--$7$B, не $1$B, и это надо померить, а не постулировать.
\end{pitfall}

\section{Инструмент $\varphi$: скорер по осям характера}

Центральный инструмент — скорер характера $\varphi$ (измерительный прибор: строишь
прибор — и он де-конфузит весь отбор). Берём $d$ осей: аффективная температура,
колкость, теплота, агентность, участие-против-сервиса, тип близости, тип юмора,
стиль конфликта. На каждую — рубрика (описание $+$ якорь для оценок $0$ и $10$).
Тогда $\varphi(\text{текст})\in\mathbb{R}^d$ — прогон LLM-судьи. Отдельно задаём
целевой $\mathbf{t}$. На старте — LLM-as-judge; классификаторы по осям позже.

\section{Суперпозиция муз — это отбор в признаковом пространстве}

Каждого персонажа можно оценить по каждой оси — персонаж есть точка в $d$-мерном
пространстве характера. Музы — существующие персонажи, свои точки.

Наивная суперпозиция, то есть обучение на смеси полных корпусов, даёт среднее в
центроиде облака муз, а не в $\mathbf{t}$: $\mathbb{E}[\varphi]=\text{centroid}\ne\mathbf{t}$.
Центроид не есть Рин; интересна муза именно отклонением от центроида. Усреднение
вдобавок размывает отличительные черты, потому что растёт дисперсия по осям.

Правильно — построить выборку с $\mathbb{E}[\varphi]\approx\mathbf{t}$: не сваливать
корпуса, а отбирать примеры, чтобы их агрегатный центр сидел в $\mathbf{t}$.
Технически — отбор данных или rejection sampling под целевой стилевой вектор:
(1)~операционализируем оси; (2)~скорим $\varphi(\text{пример})$;
(3)~отбираем или перевзвешиваем, пока $\mathbb{E}[\varphi]\approx\mathbf{t}$.
Короче: не «смешать музы», а «отфильтровать корпус к $\mathbf{t}$».

\section{Пайплайн по стадиям}

\begin{center}
\begin{tikzpicture}[font=\scriptsize,>=Stealth,node distance=6mm,
  b/.style={draw,rounded corners,align=center,minimum height=8mm,text width=20mm}]
  \node[b] (s1) {Стадия 1\\ сбор реплик муз};
  \node[b,right=of s1] (s2) {Стадия 2\\ отбор по $\varphi\!\to\!\mathbf{t}$};
  \node[b,right=of s2] (s3) {Стадия 3\\ LoRA: голос};
  \node[b,right=of s3] (s4) {Стадия 4\\ ситуация\\ $\to$ реакция};
  \node[b,right=of s4] (s5) {Стадия 5\\ замкнуть на граф};
  \draw[->] (s1)--(s2);
  \draw[->] (s2)--(s3);
  \draw[->] (s3)--(s4);
  \draw[->] (s4)--(s5);
  \draw[->] (s5.south) .. controls +(0,-0.8) and +(0,-0.8) .. (s3.south)
    node[midway,below,font=\scriptsize]{консолидация $\to$ новые данные};
\end{tikzpicture}
\end{center}

\paragraph{Стадия 1 — сбор.}
Субтитры (\texttt{.srt}, \texttt{.ass}), фан-скрипты, quote-вики, сценарии VN.
Извлекаем строки целевых персонажей. Узкое место — атрибуция говорящего. Идея
кормить видео каптионеру возможна, но окупается позже и лишь ради данных о
физическом поведении: для голоса субтитры дают почти всё дешевле.

\paragraph{Стадия 2 — отбор (фильтрацией, не написанием).}
$\varphi$ по каждой реплике; оставляем близкие к $\mathbf{t}$. Это корпус голоса.

\paragraph{Стадия 3 — голос (LoRA).}
Дообучаем малую модель на отобранных репликах — языковое моделирование на
тексте-в-стиле. Ситуации не нужны. LoRA обучает низкоранговую добавку
$\Delta W = BA$, $\mathrm{rank}(BA)=r\ll\dim$, малую по памяти; влезает на $8$\,ГБ
для малой модели.

\paragraph{Стадия 4 — выбор реакции.}
Ситуации нужны, эмбеддинги мостят несовпадение.
\emph{(A)~Синтетический} (реком.): ситуации из базы смысла $\to$ промпт сильной
модели в голосе $\to$ реакции $\to$ фильтр $\varphi$ — пары
(ситуация$\to$реакция).
\emph{(B)~Майнинг+retrieval}: из сцен муз пары (ситуация, реакция), эмбеддинг
ситуаций, в рантайме ближайший сосед. Именно здесь «ситуации могут не совпадать,
эмбеддинги справятся».

\paragraph{Стадия 5 — замкнуть на граф.}
Малая голос-модель $+$ граф; граф даёт контекст (mood-смещённый retrieval), модель
генерит в голосе; выходы (фильтр $\varphi$ $+$ значимость \eqref{eq:salience})
консолидируются обратно.

\section{Bootstrap: генератор $\ne$ деплой-модель}

Как получить характер-образующий граф без готовой Рин? Финальная Рин не нужна —
зачаток есть: prompt-Рин (модульные промпты $+$ базовая модель).

\begin{enumerate}
\item строим нейтральный эмбриональный граф (глава~\ref{ch:experiment});
\item селим prompt-Рин как bootstrap-генератор: живёт на графе, генерит характерные
реакции;
\item консолидация пишет отпечатки обратно — граф становится character-shaped через
использование, а не ручным reorganization-проходом;
\item на накопленном поведении дообучаем малую модель.
\end{enumerate}

\begin{pitfall}[Стабильность и потолок]
Аудит логов: prompt-Рин — Рин с перебоями. Ядро иногда пробивает хорошо (реплики,
дословно совпадающие с установкой мотивации), но свободный чат деградирует в липкие
повторы, часть ответов нарушает явные правила, и guided чище raw — по внутреннему
критерию «ядро промпта недостаточно сильно» (потолок $8$B). Вывод: как
безфильтровый bootstrap она шумна. Два хода: (i)~сидировать не в real-time, сильной
моделью (CPU-offload, квантизация или разовый API), а малую сажать на обогащённый
граф — генератор и деплой разные; (ii)~фильтровать выходы $\varphi$.
\end{pitfall}

\section{Raw против guided и псевдо-DPO}

Raw менее стабилен, но на пике ближе к Рин; guided стабилен, но «ассистентский»,
то есть стабильно мимо цели. Для деплоя важна стабильность, для сидирования — нет:
нужны сэмплы пиковой Rinness. Низкий потолок фильтрацией не поднять, а
нестабильность — можно: сэмплируй много, фильтруй $\varphi$. Для bootstrap:
raw $+$ oversample $+$ отбор.

Отсюда дешёвый способ собрать предпочтительные пары для псевдо-DPO. DPO обучается на
тройках (промпт, выбранный $y_w$, отвергнутый $y_l$), максимизируя
\[
\log \sigma\!\Big(
\beta \log \frac{\pi_\theta(y_w\mid x)}{\pi_{\text{ref}}(y_w\mid x)}
- \beta \log \frac{\pi_\theta(y_l\mid x)}{\pi_{\text{ref}}(y_l\mid x)}
\Big).
\]
Подавать raw и guided как двух кандидатов на один промпт и ранжировать, частично
автоматически через $\varphi$, частично руками, — чистый способ получить реально
разных кандидатов: раз raw иногда лучше, guided иногда — сигнал живой.

\section{Мост между базами: эмбеддинг ConceptNet}

Узлы ATOMIC — события с плейсхолдерами (\textsc{PersonX breaks a cup}), а не
концепты; «кружки» как узла там нет. Мост — извлечь контент-слова события и
слинковать к концепт-узлам ConceptNet. Но не bag-of-words и не one-hot: они не знают
синонимии (mug $\ne$ cup, разбить $\ne$ сломать), мост выйдет дырявый. Берём
ConceptNet Numberbatch — родной эмбеддинг ConceptNet: сначала точный лемма-матч,
эмбеддинг-фолбэк. Тот же эмбеддинг обслуживает $\text{rel}(B\mid q)$ в
\eqref{eq:trans} и сидирование валентности. Одно пространство — три задачи.

\chapter{Рантайм и инженерия}
\label{ch:runtime}

До сих пор мы описывали, что делают операторы. Эта глава — про то, как они реально
запускаются, приостанавливаются, где хранятся данные и во что это обходится по памяти
и времени. Это инженерный слой, по которому предстоит работать, поэтому здесь мы
конкретны, а оценки даём даже там, где они лишь порядковые.

\section{Планировщик как конечный автомат}
\label{sec:scheduler}

Никакого «мистического решающего разума», дёргающего операторы, нет. Есть
\emph{цикл событий} с приоритетным планировщиком. Управление двухуровневое
(глава~\ref{ch:walker}): грубый \emph{режим} задаётся конечным автоматом, а внутри
режима \emph{диспетчер} раздаёт операции по бюджету $\rho$.

\begin{definition}[Единственный дефицитный ресурс]
Почти всё в системе дёшево (обход, статистика, распад, векторный поиск — миллисекунды,
раздел~\ref{sec:costs}). \emph{Дорог} только один ресурс — \textbf{слот
LLM-инференса} (один GPU, вызовы сериализуются). Поэтому планировщик по сути
распределяет \emph{слоты LLM-вызовов} между передним планом (ответ) и фоном
(fusion, метакогниция, переваривание). Всё остальное диспетчеру раздать легко.
\end{definition}

\subsection{Режимы (состояния автомата)}

\begin{center}
\begin{tabular}{p{0.16\textwidth} p{0.10\textwidth} p{0.60\textwidth}}
\toprule
\textbf{Режим} & $\rho$ & \textbf{Что делает диспетчер} \\
\midrule
\textsc{Engaged} & высок & приоритет — retrieval $+$ композиция ответа; фон приглушён/приостановлен \\
\textsc{Idle} & $\approx 0$ & фоновый разведчик (любопытство $+$ лёгкий fusion), лёгкая консолидация \\
\textsc{Sleeping} & $\approx 0$, долго & телепорт-скраб, тяжёлая консолидация, репетиция важного \\
\textsc{Reflecting} & сред. & объяснение сработавшей аномалии (метакогниция); под-режим \textsc{Idle} \\
\bottomrule
\end{tabular}
\end{center}

\subsection{Переходы}

\begin{lstlisting}[style=pseudo]
on user_message:            # экзогенное прерывание, наивысший приоритет
    freeze_background()     # заморозить текущее фоновое рассуждение (раздел 18.2)
    mode <- ENGAGED
on response_delivered and no pending_foreground:
    mode <- IDLE
on idle_timeout:            # долгий простой
    mode <- SLEEPING
on detector_fired(anomaly) and salient:
    mode <- REFLECTING      # вернётся в прежний режим по завершении
on any_state + user_message:
    preempt -> ENGAGED
\end{lstlisting}

Ключевой момент — \textbf{смещение фокуса «любопытная Рин $\to$ ответить юзеру» это
рефлекс, а не решение LLM}: сообщение пользователя — высоко-salience экзогенное
событие, оно вытесняет фон как аппаратное прерывание ($\rho$ подскочил $\Rightarrow$
фон приглушён). LLM подключается лишь на настоящих развилках суждения (бросить нить
или припарковать; вставить ассоциацию или остаться в теме) — это \emph{эскалация} к
делиберации (System~2), редкая. Рефлексы (System~1) закрывают подавляющее
большинство запусков.

\subsection{Реализация}

Простейшее — однопоточный цикл событий (события: \texttt{user\_msg},
\texttt{timer\_tick}, \texttt{detector\_fire}, \texttt{op\_complete}). Для
параллелизма — небольшой пул: поток переднего плана (ответ) и фоновый поток
(разведка), сериализуемые на общем ограничении — числе доступных LLM-слотов (так
сделано и в ReplicantLife у Toy и др.). То есть настоящий планировщик планирует
\emph{GPU-слоты}, а не «мысли».

\begin{openq}[Граница эскалации к LLM]
Какие именно решения остаются рефлексом, а какие эскалируют к LLM-делиберации, —
подвижная граница. Слишком много эскалаций $\Rightarrow$ LLM-бутылочное горло и
тормоза; слишком мало $\Rightarrow$ Рин ведёт себя механически. Настраивается на
стенде.
\end{openq}

\subsection{Произвольное внимание: внутренний forensic-запрос в Reflecting}
\label{sec:reflecting}

В таблице режимов \textsc{Reflecting} помечен как «объяснение аномалии», но это
недосказано: тут прячется механизм, которого в мануале не было явно — \emph{как Рин
думает сама с собой, без внешнего триггера}. Всё внимание до сих пор было экзогенным
(сообщение юзера как прерывание). \textsc{Reflecting} — это \emph{эндогенное,
произвольное} внимание, и оно сшивает три уже описанных куска:
\begin{enumerate}\itemsep2pt
\item \textbf{Триггер (метакогниция).} Дешёвые детекторы на кольцевом буфере телеметрии
(гл.~\ref{ch:mood}) ловят системный сбой — скачок \texttt{need\_to\_revise}, конфликт
уверенностей, затяжную руминацию.
\item \textbf{Переход (планировщик).} Автомат уходит в \textsc{Reflecting}
(\texttt{detector\_fired} выше).
\item \textbf{Механизм (forensic-ретрив).} Рин \emph{сама себе} кидает направленный
\emph{forensic}-запрос (раздел~\ref{sec:retrievaltax}) — не ассоциативный обход, а точный
подъём дерева обоснований по провенансу через query-planner — и запускает ретракцию лжи
по TMS (гл.~\ref{ch:epistemics}). Это и есть моделируемый внутренний монолог: не «что
всплыло», а «на чём это стоит и не рухнуло ли основание».
\end{enumerate}

\begin{pitfall}[Тормоза руминации обязательны]
Без ограничителей эндогенное внимание уходит в \emph{self-fulfilling-loop}: детектор
руминации сам себя кормит, Рин зависает в бесконечном самоаудите (это тот же риск, что и
самоисполняющаяся тревога, гл.~\ref{ch:mood}). Поэтому \textsc{Reflecting} жёстко
гейтируется: \emph{рейт-лимит} (не чаще, чем раз в такое-то окно), \emph{event-driven
каденция} (только на \emph{значимую} аномалию, не на каждый микроконфликт уверенностей)
и \emph{бюджет} на глубину\,/\,длительность рефлексии. Выгорел бюджет — вышли из режима,
пометив вопрос как открытый, а не крутимся дальше.
\end{pitfall}

\section{Приостановка и возобновление рассуждений}
\label{sec:suspend}

Когда сообщение пользователя вытесняет фоновое рассуждение, будь то середина
fusion-цепочки или любопытной прогулки, состояние надо заморозить, чтобы потом
продолжить с того же места, а не с нуля.

\subsection{Что такое «состояние рассуждения»}

Фоновая задача — это walker (или рой) с: текущим \emph{фронтом} (множество активных
узлов $+$ их активации), \emph{затравкой/целью} (что искали), \emph{частичными
результатами} (уже собранные кандидаты, поставленные \texttt{need\_to\_discover}) и
\emph{курсором} (где остановились). Заморозка — это сериализация \emph{continuation-
токена}:

\begin{lstlisting}[style=pseudo]
Continuation = {
    task_id, goal/seed,
    frontier: [(node_id, activation), ...],   # компактно: id + float
    partial_candidates: [node_id, ...],
    pending: [need_to_discover, ...],
    cursor, priority, timestamp
}
\end{lstlisting}

Токен мал (id-шники $+$ числа) — это как захват состояния корутины при yield или
context-switch процесса в ОС (сохранить «регистры» = фронт $+$ курсор). Кладётся в
очередь приостановленных задач.

\subsection{Возобновление и устаревание}

При возврате в \textsc{Idle} задача достаётся из очереди и продолжается. Но граф мог
\emph{измениться} за время паузы (появились узлы, активации распались, что-то
спрунено), поэтому возобновление начинается с \emph{ре-валидации фронта}: выкидываем
исчезшие/распавшиеся узлы, восстанавливаем активации живых, продолжаем обход.

Не всякую задачу стоит возобновлять: пока Рин отвечала, разговор мог уйти, и нить
устарела. Поэтому у приостановленных задач своё \emph{устаревание}: если не
возобновлена за окно или вытеснена более свежей — она \emph{не удаляется бесследно}, а
роняется в граф как \emph{«недодуманная нить»} (узел с \texttt{need\_to\_discover}),
которую позже может подобрать телепорт-скраб во сне (глава~\ref{ch:sleep}). Так
прерванная мысль не теряется, а становится зацепкой для фоновой разведки.

\begin{intuition}[Замороженная мысль — это savepoint, а не труп]
Полезно думать о continuation-токене как о savepoint базы данных или о зелёном
потоке, снятом с исполнения: он хранит ровно столько, чтобы продолжить, и ничего
лишнего. А брошенная нить не исчезает — она оседает в память как «я не додумала про
$X$», и это честно по-человечески.
\end{intuition}

Каждая пара приостановка/возобновление — \emph{событие телеметрии} (глава~\ref{ch:mood}):
паттерн «меня постоянно прерывают на середине мысли» \emph{детектируется} метакогницией
как аномалия.

\section{Хранилище и индексирование}
\label{sec:storage}

Разберём предметно, из чего состоит хранилище и как оно индексируется, — и заодно
ответим, заменяет ли многоуровневое хранилище простое «окно внимания».

\subsection{Из чего состоит хранилище}

\begin{itemize}\itemsep1pt
\item \textbf{Node store} — ключ-значение \texttt{node\_id} $\to$ \{content,
embedding, метаданные\}; на диске (встраиваемое KV-хранилище: LMDB/SQLite и т.п.).
\item \textbf{Edge store} — списки смежности \texttt{node\_id} $\to$
[(target, type, weight, sign), \dots].
\item \textbf{Vector index} — индекс по эмбеддингам для поиска похожих (HNSW/FAISS),
\textbf{резидентный в RAM}; держит эмбеддинги \emph{горячего/тёплого} множества и
служит механизмом \emph{поиска якорей} по запросу (глава~\ref{ch:walker}). Важно:
\emph{поиск идёт по этому RAM-индексу, а не по диску} — на диске лежит лишь
персистентная копия эмбеддингов (backing store) и cold-тир.
\item \textbf{Структурные индексы} — по \texttt{layer}, по \texttt{importance} и т.п.,
для планировщика и распада.
\end{itemize}
Индексов, стало быть, два по сути: \emph{векторный} (семантическая близость, чтобы
найти вход) и \emph{графовая смежность} (чтобы идти от входа). Retrieval $=$ векторный
поиск якорей $\to$ обход графа от них.

\subsection{Attention-radii против тиров: это не одно и то же}

Твоя догадка про «окно/радиусы внимания» верна, но покрывает лишь половину.

\emph{Attention-radius} — это \textbf{транзиентная} граница на \emph{обход}: сколько
хопов walker расширяется от якорей за один запрос. Она держит \emph{рабочее множество}
маленьким \emph{в момент запроса}, независимо от размера графа.

\emph{Тиры} (hot/warm/cold) — это \textbf{персистентная} организация \emph{хранилища},
живущая между запросами.

Почему нужно и то, и другое: attention-radius ограничивает \emph{прогулку}, но \emph{не}
ограничивает \emph{векторный индекс якорей}. Если эмбеддинги \emph{всего} графа лежат
в одном плоском индексе, стоимость поиска якорей и RAM растут с общим числом узлов.
Тиры выносят холодные узлы \emph{из} горячего индекса (они достижимы лишь через
graph-указатели), удерживая индекс маленьким. Плюс это режет \emph{шум} ретрива: в
огромном плоском индексе косинусная близость тащит семантически близкое, но
контекстно нерелевантное (ровно жалоба Toy и др.).

\begin{center}
\begin{tabular}{p{0.22\textwidth} p{0.34\textwidth} p{0.34\textwidth}}
\toprule
 & \textbf{attention-radius} & \textbf{тиры (hot/warm/cold)} \\
\midrule
природа & транзиентная (на запрос) & персистентная (хранилище) \\
ограничивает & глубину обхода (working set) & размер индекса якорей, RAM, шум \\
когда хватает одного & если граф \emph{мал} & если граф \emph{растёт вечно} \\
\bottomrule
\end{tabular}
\end{center}

Итог: \textbf{малый граф — attention-radius достаточно, тиры не нужны; большой/растущий
— нужны оба.} Тир cold хранит огрублённый gist (глава~\ref{ch:timescales}): в warm-
индексе остаётся узел-указатель домена («знаю цепи Маркова; детали в архиве»), полная
деталь лежит в cold сжатой и \emph{подкачивается} по требованию (это и есть «чуть
дольше добывать»), как paging в ОС.

\begin{openq}[Политика тиринга]
Когда демоутить узел в cold, насколько огрублять warm-gist, когда пейджить деталь
обратно — параметры, которые надо мерить (глава~\ref{ch:experiment}). По голому диску
(раздел~\ref{sec:costs}) тиринг не обязателен; он оправдан скоростью и чистотой
ретрива, а не экономией места.
\end{openq}

\subsection{Кто к кому обращается: путь запроса и стоимости}
\label{sec:dataflow}

Ключевое уточнение, снимающее частую путаницу: \textbf{по эмбеддингам на диске никто
не ищет}. Поиск идёт по векторному индексу \emph{в RAM}; диск — это backing store и
cold-тир. Горячий путь диск почти не трогает. Ниже — кто, когда и куда обращается.

\paragraph{Путь чтения («жизнь запроса»).}
\begin{enumerate}\itemsep1pt
\item \textbf{Эмбеддинг запроса}: эмбеддер (GPU/CPU), текст $\to$ вектор. $\sim$1--10\,мс.
\item \textbf{Поиск якорей}: вектор $\to$ индекс HNSW \emph{в RAM} $\to$ top-$k$
\texttt{node\_id}. $\sim$мс (sub-10\,мс даже при $10^6$). \emph{Диска нет.}
\item \textbf{Обход}: от якорей по igraph \emph{в RAM} (ограниченный радиус,
mood-bias) $\to$ активированная окрестность. $\sim$мс--10\,мс. \emph{Диска нет.}
\item \textbf{Достать контент}: для горстки финальных узлов (5--50) текст из
node-store — RAM-кэш (горячее) или SSD (несколько KV-чтений, sub-мс, обычно уже в OS
page cache). Пренебрежимо.
\item \textbf{Композиция}: контекст $\to$ LLM. GPU-VRAM. $\sim$1--3\,с.
\textbf{Доминирует.}
\end{enumerate}
То есть всё, кроме LLM, — в RAM и $\sim$мс; диск трогается лишь чтобы достать
\emph{текст} нескольких финальных узлов (дёшево), а \emph{не} для поиска.

\paragraph{Путь записи (новый узел / fusion).} Контент+эмбеддинг+метаданные $\to$
node-store (диск, асинхронно); эмбеддинг $\to$ insert в RAM-индекс (HNSW $\sim$мс,
амортизированно); узел/рёбра $\to$ в igraph (мутации батчить, глава~\ref{ch:experiment}).
Дёшево.

\paragraph{Cold-путь (единственная «дисковая» латентность, и она редкая).} Запросу
нужна архивная деталь $\to$ в warm-индексе есть узел-указатель домена $\to$ вторая
ступень: чтение из cold-store (SSD) $+$ декомпрессия $\to$ подкачка в рабочую память.
$\sim$мс--десятки\,мс. Редко, потому что hot/warm держим в RAM.

\paragraph{Старт.} При загрузке структура графа и векторный индекс поднимаются или
mmap-ятся в RAM — одноразово, не на каждый запрос.

\begin{center}
\begin{tabular}{p{0.05\textwidth} p{0.30\textwidth} p{0.22\textwidth} p{0.28\textwidth}}
\toprule
\# & \textbf{Операция} & \textbf{Ярус} & \textbf{Стоимость} \\
\midrule
1 & эмбеддинг запроса & GPU/CPU & $\sim$1--10\,мс \\
2 & поиск якорей (HNSW) & \textbf{RAM} & $\sim$мс \\
3 & обход (igraph) & \textbf{RAM} & $\sim$мс--10\,мс \\
4 & контент финальных узлов & RAM-кэш / SSD & sub-мс, пренебр. \\
5 & композиция (LLM) & \textbf{GPU-VRAM} & \textbf{$\sim$1--3\,с (узкое место)} \\
— & подкачка cold-детали & SSD & $\sim$мс--десятки\,мс (редко) \\
\bottomrule
\end{tabular}
\end{center}

\section{Event sourcing: неизменяемый лог под изменяемым графом}
\label{sec:eventlog}

Граф памяти не должен быть единственным хранилищем прошлого. Лучше три уровня:
\begin{itemize}\itemsep1pt
\item \textbf{Event log} — неизменяемая, только-на-дозапись история всего, что
случилось: наблюдения, коммиты, ретракции, исправления, удаления, консолидации.
\item \textbf{Memory graph} — текущая материализованная модель памяти,
\emph{перестраиваемая проекция} поверх лога.
\item \textbf{Indexes} — производные структуры для поиска (векторный, forensic).
\end{itemize}
Запись лога: \{timestamp, тип (observed/committed/retracted/consolidated/deleted/
corrected), payload, source\_context, owner\}. Важно: \texttt{source\_context} и
\texttt{owner} \emph{рождаются здесь}, в момент захвата — отсюда и берётся провенанс
(глава~\ref{ch:substrate}).

Зачем это несущее: debug («почему Рин так решила» — реплеем лога), rollback,
пересборка графа после смены схемы, \textbf{сравнение версий операторов} — поменял
алгоритм консолидации, пересобрал проекцию из \emph{того же} лога, сравнил (идеально
ложится на стенд, глава~\ref{ch:experiment}), удаление, миграции.

\begin{openq}[Рост лога и компакция]
Неизменяемый лог растёт неограниченно. На старте — хранить и не думать (разрастётся —
решим). Дальше нужна политика: периодический \emph{снапшот} проекции $+$ обрезка
старого лога, либо тиринг лога в cold (раздел~\ref{sec:storage}). Точный порог и
формат снапшота — открыты.
\end{openq}

\subsection{Снапшот-чтения: почему обходчик и коммиттер не сталкиваются}
\label{sec:snapshotreads}

Раз обходчиков много и они работают параллельно с коммиттером, напрашивается страх
гонок: RW читает узел, который коммиттер как раз меняет — не получит ли он рваное,
полу-обновлённое состояние? При нашей раскладке ролей (walker читает и предлагает,
структурно пишет только сериализованный коммиттер, гл.~\ref{ch:walker},
раздел~\ref{sec:proposecommit}) страх почти растворяется, и event sourcing закрывает
остаток.

Ключ в том, что «изменение» узла у нас — \emph{не} правка на месте, а \emph{новая
версия}: коммиттер дописывает событие в лог и обновляет проекцию, старая версия узла
никуда не пропадает мгновенно. Значит обходчики читают \emph{консистентный снапшот}
проекции (MVCC-подобно): RW2, зашедший в «меняющийся» узел, видит его до-коммитную
версию целиком — рваных чтений нет, потому что нет правки на месте, которую можно
застать на середине. Read $\leftrightarrow$ read не конфликтует никогда; read
$\leftrightarrow$ write решается версионностью; write $\leftrightarrow$ write не
существует по построению (писатель один). Остаётся лишь мягкое побочное состояние —
рабочая яркость $a_{\text{work}}$ и телеметрия внимания; оно \emph{коммутативно} (для
яркости — через максимум, не сумму, гл.~\ref{ch:walker}), поэтому порядок конкурентных
подсветок неважен и блокировки не нужны.

\begin{intuition}[Гонку не арбитрируют — её конструктивно исключают]
Мы не пишем сложный протокол блокировок, кто кого пропускает. Мы устроили систему так,
что сталкиваться \emph{нечему}: один структурный писатель, версионные чтения для
всех остальных, коммутативная телеметрия. Это дешевле и надёжнее любого арбитража
постфактум — та же философия, что «противоречие лучше не разрешать задним числом, а
не допускать молчаливой перезаписи» (гл.~\ref{ch:sleep}).
\end{intuition}

\subsection{Устаревшие предложения: коммиттер валидирует, а не просто пишет}
\label{sec:staleproposals}

Снапшот-чтения убирают гонку на \emph{чтении}, но остаётся тонкий изъян на
\emph{записи}. Обходчик мог посчитать предложение на версии графа
\texttt{snapshot=123}, а к моменту коммита проекция уже на \texttt{129}, где часть
посылок успела устареть: узел-родитель ретрактнут, слит, понижен в уверенности,
изменил \texttt{layer} или \texttt{visibility}. Если коммиттер применит предложение
вслепую, гонку по данным мы убрали, но получили \emph{семантически устаревший коммит} —
что не лучше.

Поэтому предложение — не голая операция, а операция с \emph{контекстом валидности}:
\{\texttt{op}, \texttt{members}, \texttt{read\_snapshot\_id}, \texttt{parents},
\texttt{proposal\_ttl}, \texttt{validity\_conditions}\}. Коммиттер перед применением
\emph{валидирует}, а не сразу пишет: живы ли ещё родители? не deprecated? не изменились
ли их \texttt{confidence}\,/\,\texttt{layer}\,/\,\texttt{visibility} с момента
\texttt{read\_snapshot\_id}? не вышел ли \texttt{proposal\_ttl}? По итогу — одно из
четырёх: \emph{применить} (условия держатся), \emph{rebase} (пересчитать на текущей
версии, если посылки лишь сдвинулись), \emph{downgrade} (принять как более слабую
гипотезу) или \emph{reject} (посылки развалились). Это оптимистический контроль
конкуренции: не блокируем заранее, а проверяем на коммите. Без этого read/write-гонку
уберёшь технически, а память тихо наполнится решениями, принятыми по уже неверным
предпосылкам.

\subsection{«Забудь это»: удаление при неизменяемом логе}
\label{sec:forget}

Неизменяемый лог порождает очевидный конфликт: а что происходит, когда пользователь
говорит «Рин, забудь последние 10 минут»? Для локального компаньона это не юридическая
занудность, а конкретный будущий UX, и он должен что-то значить.

Ключ к разрешению: \emph{неизменяемость $\ne$ «нельзя удалять»}. Она означает
\emph{append-only, без тихих правок на месте}; само удаление — это \emph{тоже
appended, аудируемое событие}, а не нарушение дизайна. «Забыть» — первоклассная,
поддержанная операция, и реализуется тремя механизмами, которые работают вместе:

\begin{itemize}\itemsep2pt
\item \textbf{Граф (легко).} Убираем узел, ретрактим, перепропагируем. Каскад: удаление
исходного наблюдения понижает/сносит производные узлы, чьим \emph{провенансом} оно
было (ретракция по TMS, глава~\ref{ch:epistemics}). «Забудь X» — это ещё и «забудь
то, что из X выведено».
\item \textbf{Tombstone.} В лог дописывается delete-событие; при пересборке проекции
удалённое пропускается. Эффект немедленный, но сырое событие ещё \emph{в логе}
(годится для undo/аудита, но это не настоящее стирание).
\item \textbf{Crypto-shredding.} События (по сессии/скоупу) зашифрованы своим ключом;
«удалить» $=$ уничтожить ключ $\to$ шифртекст остаётся, но навсегда нечитаем.
\emph{Мгновенное, проверяемое} стирание без переписывания лога — идеально под «забудь
последние 10 минут» (эти минуты под эфемерным сессионным ключом; шреднул ключ —
стёр).
\item \textbf{Compaction.} Фоновая пересборка лога выбрасывает tombstoned-события —
сырое физически исчезает.
\end{itemize}

То есть «Рин, забудь X» $=$ tombstone $+$ каскадная ретракция производных $+$
шреддинг ключа (истинное стирание) $+$ компакция-уборка. Конкретно, а не «ну она как-
то забудет».

\begin{pitfall}[Crypto-shred рождает key management и требует taint-propagation]
Две дыры, без которых схема красива, но дырява. \emph{Первая}: шреддинг создаёт новый
центр тяжести — \textbf{управление ключами} (где ключи, кто их держит, гранулярность
сессия/скоуп). \emph{Вторая, важнее}: уничтожив ключ raw-события, ты \emph{не} стёр
след, если plaintext успел растечься — в debug-log, снапшот context-packer,
embedding-кэш, вектор-индекс, бэкап, telemetry-сводку или производный fusion-узел.
Поэтому рядом с «забудь» нужна дисциплина \textbf{taint-propagation}: всякий
производный объект наследует taint/провенанс от породившего события, а delete-cascade
чистит \emph{все} производные хранилища — граф, индексы, сводки, эмбеддинги, снапшоты,
pending-очереди, бэкапы. Иначе «забыл» будет только на бумаге.
\end{pitfall}

\begin{openq}[Гранулярность забвения]
На каком уровне держать ключи шифрования (событие / сессия / скоуп), что именно
каскадит при удалении (только прямые производные или транзитивно), и как показывать
пользователю, что забыто, — открыто. Но механика («забыть» = событие в логе $+$
shred $+$ каскад) зафиксирована.
\end{openq}

\section{Write-path: что попадает в память, а что нет}
\label{sec:writepath}

Отдельно проговорим, что именно записывается, потому что «не писать» — полноправный
исход, а не сбой. Конвейер:
\begin{lstlisting}[style=pseudo]
raw_event -> parse/extract candidates
          -> [ epistemic_gate, salience_gate, anti_fusion_gate ]
          -> write | skip
          -> (if write) propagate
          -> audit_log   # и write, и skip логируются
\end{lstlisting}
\begin{itemize}\itemsep1pt
\item \textbf{epistemic-гейт}: стоит ли этому верить и с какой уверенностью (гл.~\ref{ch:epistemics}).
\item \textbf{salience-гейт}: значимо ли это для Рин (appraisal, гл.~\ref{ch:personality}).
\item \textbf{anti-fusion-гейт}: см.\ ниже.
\end{itemize}

\begin{principle}[response-useful association $\ne$ memory-worthy fusion]
Оператор должен уметь не только «нашёл красивую связь $\to$ породил узел», но и
явно \emph{не писать}: «со-активация есть, связного вывода нет», «связь эстетическая,
не эпистемическая», «это каламбур, не убеждение», «полезно для реплики, но не для
памяти». Рин может \emph{использовать} ассоциацию в ответе, не записывая её узлом —
иначе память зарастает «интересными мыслями», которые были хороши в моменте, но не
являются знанием.
\end{principle}

И \texttt{skip} обязательно логируется (хотя бы в debug): иначе не поймёшь, почему
Рин «не запомнила» важное.

\subsection{Коммит не бинарен: полосы, а не «только ночью»}
\label{sec:commitlanes}

Соблазнительно сказать: «walker'ы читают весь день, а пишет коммиттер только ночью».
Это слишком сильно и ломает сознательное мышление. Если так сделать буквально, Рин
днём станет странной: ходит по графу, находит связи, что-то «понимает» — но не может
ни закрепить свежий внешний факт, ни воспользоваться дневной догадкой иначе как
эфемерным текстом. Правильный инвариант:

\begin{principle}[Коммит не бинарен]
Наблюдаемые события и рабочие гипотезы фиксируются \emph{немедленно}; долговременная
консолидация \emph{отложена}. Коммиттер один и сериализован, но работает не «ночью
против дня», а в нескольких \emph{полосах} по стоимости и durability.
\end{principle}

Полосы одного коммиттера:
\begin{itemize}\itemsep2pt
\item \textbf{immediate.} Пишется всегда и сразу: append в event-log (что произошло —
это \emph{журнал}, а не «убеждение»), прямые пользовательские наблюдения,
task-critical внешние факты, явные исправления, важные boundary\,/\,preference.
Откладывать это до ночи нельзя — потеряешь след происходящего.
\item \textbf{foreground pending.} Дневная рабочая память — это наш \emph{session-слой}
(гл.~\ref{ch:timescales}), не новая сущность: временные узлы-гипотезы, найденные
факты, цепочки рассуждения, кандидаты на консолидацию, \texttt{need\_to\_verify}. Живут
с TTL, низкой уверенностью, провенансом и статусом \texttt{pending} — уже объект,
которым можно \emph{пользоваться сейчас}, но ещё не вечная память.
\item \textbf{idle.} Микробатчи в простое: мелкий dedup, малые правки уверенности,
локальные проверки противоречий.
\item \textbf{sleep.} Тяжёлое: консолидация $A,B,C,D\to X$, массовое разрешение
конфликтов, fusion далёких ассоциаций, pruning, reindex, пересборка summary, ревизия
низкоуверенных потомков, promotion рекуррентных кандидатов.
\end{itemize}
Это и есть fast-path\,/\,slow-path: дешёвые $O(1)$-мутации (добавить узел-гипотезу,
инкремент) идут днём по быстрой полосе; дорогие $O(|V|\log|V|)$ и глобальные — ночью.
Критично: быстрая полоса — про \emph{латентность}, а не про поблажку. Она проходит тот
же write-path и те же гейты (epistemic\,/\,salience\,/\,anti-fusion,
раздел~\ref{sec:writepath}), просто синхронно.

Сводно, по трём процессам (и подчеркнём: строка записи \emph{не} ночецентрична — днём
работает immediate-lane):

\begin{center}
\begin{tabular}{p{0.19\textwidth} p{0.36\textwidth} p{0.36\textwidth}}
\toprule
\textbf{Процесс} & \textbf{Днём ($\rho\!\approx\!1$)} & \textbf{Фон\,/\,ночь ($\rho\!\approx\!0$)} \\
\midrule
Retrieval & локальный разреженный matvec-диффузия из HNSW-якорей (fast-path) & телепорт-скраб, эргодическое покрытие, большой $\alpha$ \\
\shortstack[l]{Fusion\\мышление} & «мотыльки»: одиночные RW на остывающем $a_{\text{work}}$ $\to$ кандидаты в thought-inbox & честный рой (slow-path): треки, $n$-арные цепочки, далёкие мосты \\
Запись & \textbf{immediate-lane}: append в event-log, безопасные foreground-микрокоммиты (\texttt{observed}, \texttt{boundary}, правки \texttt{preference}); \texttt{pending} догадки — в thought-inbox & коммиттер (slow-path): судит гипотезы, тяжёлая консолидация, prune, reindex \\
\bottomrule
\end{tabular}
\end{center}

Так снимаются оба дневных противоречия. \emph{Сознательный поиск} не ждёт ночи: Рин
«гуглит здесь и сейчас», кладёт результат в event-log и session-граф как sourced
observation, использует в ответе — но со статусом «свежий источник, не консолидировано»,
а не «я давно это знаю». \emph{Догадки} тоже случаются днём — просто живут как активная
гипотеза, а взрослеют в durable-память позже.

\subsection{Thought-inbox: как дневная догадка попадает в реплику, не став фактом}
\label{sec:thoughtinbox}

Фоновый разведчик работает всегда — даже в напряжённом диалоге, на остаточном бюджете
$B_{\text{bg}}=B_{\text{total}}(1-\rho)$ (гл.~\ref{ch:walker}), который никогда не
ноль. Поэтому «момент озарения», например связь House-Уилсон и Cast Away-Уилсон,
может родиться и днём. Но показывать Рин сырой след walker'ов нельзя — это шум: сотни
посещённых узлов, случайные касания, мусорные мосты. Между следом и репликой стоит
фильтр — thought-inbox: туда попадают не все касания, а только отобранные
предложения-инсайты, то есть сильная, но неочевидная со-активация далёких узлов:
\begin{lstlisting}[style=pseudo]
type: insight_candidate
content: "Уилсон у House и Уилсон в Cast Away могут делить нарративную функцию"
status: unverified_hypothesis
salience: 0.72
provenance: [node_a, node_b, path_summary]
source: background_walker
ttl: 2h
recommended_action: mention_as_hypothesis | verify_later | ignore
\end{lstlisting}
Context-packer иногда подмешивает самые «горячие» из них в промпт — но \emph{с явной
эпистемической меткой} (это ровно наш packer-слой меток, раздел~\ref{sec:packer}):
«\emph{возможная мысль, не факт: между A и B может быть связь; статус — не проверено;
использовать как догадку, не как установленную память}». Дальше цикл замыкается сам:
если контекст диалога позволяет, Рин озвучивает догадку — и в \emph{этот} момент она
легализуется, падает в event-log и становится узлом со статусом
\texttt{need\_to\_verify}; если не пригодилась — тихо затухает по TTL и убирается
сборщиком. Изоляция ролей сохранена: walker \emph{читает и предлагает}, LLM
\emph{осмысляет и озвучивает}, коммиттер \emph{пишет}.

\begin{pitfall}[Не писать активацию на каждый шаг walker'а]
«Walker дёшев» — верно относительно LLM, но не бесплатно, и узкое место обычно не в
математике, а в обвязке (индексы, сериализация, локи, provenance-traversal,
scope-фильтры, гл.~\ref{ch:experiment}). Поэтому обходчик \emph{не} пишет активацию в
хранилище на каждом шаге: он копит visits\,/\,деньты активации локально в памяти
процесса и \emph{батчем} сбрасывает в transient-activation store (по-прежнему max, не
сумма). И на сбросе различает \texttt{access\_reason}: \texttt{user\_recall} укрепляет
след $B_{\text{base}}$, а \texttt{scrub\_visit} — нет.
\end{pitfall}

\section{Таксономия ретрива: не всё — блуждание}
\label{sec:retrievaltax}

Walker как «почти всё» — переобобщение. Не всякая операция памяти естественно
является ассоциативным блужданием:
\begin{itemize}\itemsep1pt
\item \textbf{associative} (якоря $+$ walkers) — для живости и ассоциаций (гл.~\ref{ch:walker}).
\item \textbf{forensic} — точный поиск по метаданным/провенансу/времени: «покажи всё
про $X$», «узлы с источником $Y$», «что противоречит этому факту», «последняя версия
prompt-модуля». Это индексы и query-planner, \emph{не} блуждание.
\item \textbf{creative} — walkers $+$ fusion $+$ сон (гл.~\ref{ch:sleep}).
\end{itemize}
Forensic-ретрив — не приятное дополнение, а \emph{необходимость} для обслуживания:
аудит, миграции, точное удаление конкретных узлов. Walker полноты не гарантирует, а
скучный индекс — гарантирует. Для живости Рин нужны walkers; для корректности памяти
нужны скучные индексы.

\section{Context-packer: как граф превращается в промпт}
\label{sec:packer}

Это интерфейс графа и LLM на инференсе — что именно видит модель, и
\textbf{последний эпистемический гейт перед речью}. Если packer отдаёт модели плоский
список содержаний без статусов, вся дисциплина графа (уверенность, гипотезы,
депрецирование) теряется на последнем метре. Поэтому в контекст попадают \emph{не
только contents, но и метки}: \texttt{observed}, \texttt{hypothesis},
\texttt{deprecated}, \texttt{contested}, \texttt{low\_conf}, \texttt{time\_validity},
и при необходимости \texttt{do\_not\_state\_as\_fact}. Плюс провенанс для важных
утверждений и mood/persona-обусловливание. Тогда модель знает, что подать как факт,
что — как «похоже», а что не выдавать за истину вовсе. Держать компактно: короткий
контекст $=$ маленький KV-кэш $=$ меньше VRAM (раздел~\ref{sec:costs}).

Это \texttt{prompt context} — \emph{временная сборка} (граница из
главы~\ref{ch:architecture}), а не хранилище: спаковали, отдали модели, выбросили.

Здесь же — естественное место для \emph{eligibility-гейта}: перед упаковкой отсеять
узлы, которые \emph{есть в памяти, но не годятся для этого контекста} (например, не
класть в промпт то, что не стоит показывать при данном собеседнике или выносить в
публичный контент). Это шире удаления и \emph{у́же} тяжёлой governance-бюрократии:
не «кто владеет и на каких условиях», а простой guardrail
\emph{retrieval $\to$ eligibility-гейт $\to$ packer $\to$ output}. Тяжёлую систему
доступа/согласия мы сознательно не строим (вектор проекта сместился), но сам крючок
стоит держать в виду.

\begin{openq}[Eligibility как отложенный guardrail]
Полноценный слой видимости/раскрытия (что можно показать здесь, класть в промпт,
использовать в публичном контенте, обобщать) — отложен. Пока — минимальный гейт
«память может существовать, но быть негодной для данного контекста», без бюрократии.
Разворачивать ли его — решается, когда вернётся потребность в family/публичном режиме.
\end{openq}

\begin{openq}[Формат упаковки]
Как именно показывать модели статус узла (уверенность, гипотеза vs факт, провенанс),
сколько узлов включать и как их ранжировать под текущий запрос — открыто, меряется на
стенде.
\end{openq}

\section{Оценки стоимости: память, время, узкие места}
\label{sec:costs}

Оценки порядковые, для локального железа класса RTX 4090 / 8--24\,ГБ VRAM, но они
сразу показывают, где узкие места, а где нет.

\subsection{Стоимость операций}

\begin{center}
\begin{tabular}{p{0.30\textwidth} p{0.14\textwidth} p{0.46\textwidth}}
\toprule
\textbf{Операция} & \textbf{Порядок} & \textbf{Комментарий} \\
\midrule
LLM-инференс (композиция, объяснение) & \textbf{0,5--5 с} & \textbf{узкое место}; $\sim$2\,с/запрос на 7--8B локально \\
векторный поиск якорей (HNSW, $10^6$) & $\sim$мс & не узкое место \\
обход/spreading (сотни-тысячи узлов) & мс--10\,мс & C-ядро igraph, дёшево \\
пропагация уверенности (лавина) & $\sim$мс & дёшево \\
детекторы телеметрии & O(1) инкрем. & дёшево \\
распад/ACT-R & дёшево \emph{лениво} & полный проход по $N$ узлам — дорого; делать лениво (при доступе) или во сне, не на каждый тик \\
переваривание учебника & \textbf{минуты-часы} & одноразово; дорого из-за LLM-извлечения фактов; только в \textsc{Idle}/\textsc{Sleeping} \\
\bottomrule
\end{tabular}
\end{center}

\subsection{Латентность ответа}

Критический путь: \texttt{user\_msg} $\to$ поиск якорей (мс) $\to$ обход (мс) $\to$
композиция ответа (\textbf{LLM, $\sim$1--3\,с}) $\to$ ответ. \textbf{Бутылочное горло —
LLM-вызов}; всё графовое к латентности почти не добавляет. Цель — держать ответ в
$\sim$1--3\,с, то есть минимизировать \emph{число} LLM-вызовов на переднем плане (в
идеале один: граф уже подал контекст, модель только оформляет).

\subsection{Память}

\begin{itemize}\itemsep1pt
\item \textbf{VRAM} — модель (фиксировано, узкое место по размеру модели: 7--8B в
квантизации $\sim$5--8\,ГБ). Отсюда весь смысл «факты в граф, а не в веса».
\item \textbf{RAM} — векторный индекс горячего множества; для $10^6$ эмбеддингов
плоско $\sim$1,5\,ГБ, с HNSW/квантизацией меньше; \emph{тиринг} держит это ограниченным.
\item \textbf{Диск} — сам граф: $\sim$2--3\,КБ/узел, $10^5$--$10^6$ узлов $=$
0,3--3\,ГБ (глава~\ref{ch:experiment} для оценки роста). Дёшево; учебник $\approx 10^4$
узлов $\approx$ 30\,МБ.
\end{itemize}

\subsection{Где узкие места, а где нет}

\begin{principle}[Всё сводится к LLM-слотам]
Дорог только GPU/LLM — и по латентности ответа, и по фоновой пропускной способности.
Всё графовое (хранилище, обход, статистика, распад, векторный поиск) — дёшево (мс) и
узким местом \emph{не} является. Поэтому вся инженерная задача рантайма сводится к
одному: \emph{как распределить дефицитные LLM-вызовы} между быстрым ответом (передний
план) и фоном (fusion/метакогниция/переваривание). Это и делает планировщик
(раздел~\ref{sec:scheduler}).
\end{principle}

\begin{openq}[Цели по латентности и throughput]
Конкретные цели (ответ $<$ N\,с; сколько фоновых fusion в минуту в простое; когда
переключать модель на меньшую ради скорости) — под стенд и под конкретное железо.
Динамический роутинг «мелкая модель на рутину, крупная на метакогницию» (как у Toy и
др.) — очевидный кандидат оптимизации.
\end{openq}

\chapter{Эмбриональный граф: с чего начинается память}
\label{ch:embryonic}

Пустой граф памяти бесполезен: у существа не бывает мышления «с нуля» — есть
врождённый здравый смысл, на который ложится личный опыт. Поэтому память Рин
начинается с эмбрионального графа — нейтрального по характеру каркаса из готовых баз
здравого смысла, который затем обрастает личным опытом через fusion и
консолидацию. Эта глава отвечает предметно: из чего он собран, как выглядят его
узлы, как их сшить и как подключить эмбеддинги.

\section{Две базы, две разные роли}

Нужны две дополняющие базы (Wikidata как содержание не берём — человек не помнит все
столицы; глава~\ref{ch:substrate}):

\begin{itemize}
\item \textbf{ConceptNet} — предметный и таксономический здравый смысл: что чем
является, из чего состоит, для чего используется. Даёт Рин «физику мира».
\item \textbf{ATOMIC} — причинно-социальный здравый смысл о событиях и ментальных
состояниях: чего человек хотел, что почувствовал, что за этим следует. Даёт Рин
«психологию».
\end{itemize}
Порознь одного мало: ConceptNet без ATOMIC — мир без психологии; ATOMIC без
ConceptNet — реакции без предметного мира.

\section{Как выглядят узлы и рёбра ConceptNet}

ConceptNet — это уже граф. Узел — концепт, то есть слово или короткая фраза,
заданный URI; ребро — типизированная взвешенная ассоциация. Схематически запись:

\begin{lstlisting}[style=pseudo]
узлы:
  /c/ru/кружка/   /c/ru/посуда/   /c/ru/пить/   /c/ru/горячий/
рёбра (concept, relation, concept, weight):
  (/c/ru/кружка/, IsA,        /c/ru/посуда/,   2.5)
  (/c/ru/кружка/, UsedFor,    /c/ru/пить/,     1.8)
  (/c/ru/кружка/, HasProperty,/c/ru/горячий/,  1.2)
\end{lstlisting}

Типов отношений около трёх-четырёх десятков (\textsc{IsA}, \textsc{PartOf},
\textsc{UsedFor}, \textsc{HasProperty}, \textsc{RelatedTo}, \textsc{Causes},
\textsc{MotivatedByGoal}, \dots), и рёбра многоязычны: связывают понятия между
языками. Важно для нас: вес ребра $\approx$ мягкая уверенность из коробки, а типы
отношений прямо ложатся на нашу типизацию рёбер (глава~\ref{ch:substrate}).

\section{Как выглядят узлы ATOMIC}

ATOMIC устроен иначе, и это источник главной сложности сшивки. Его узлы — не
концепты, а события-фразы с плейсхолдерами действующих лиц (\textsc{PersonX},
\textsc{PersonY}). От события идут рёбра фиксированных социально-причинных типов.
Схематически:

\begin{lstlisting}[style=pseudo]
событие: "PersonX breaks a cup"
   xIntent  -> (none / accident)                # чего PersonX хотел
   xReact   -> (sad, upset)                     # что PersonX почувствовал
   xEffect  -> (cleans it up, buys a new one)   # что произошло с PersonX
   xNeed    -> (to be holding a cup)            # что требовалось до
   oReact   -> (concerned)                      # что почувствовали окружающие
\end{lstlisting}

То есть «кружки» как отдельного узла в ATOMIC нет — есть событие «разбить кружку» и
его социально-причинные следствия. ATOMIC10x — та же структура, но в несколько
миллионов примеров, порождённых языковой моделью и отфильтрованных критиком: больше
покрытия, больше шума.

\section{Сшивка: главная инженерная работа}

Соблазн сказать «форматы отличаются только на диске» — недооценка. Внутри графовой
базы это одна структура: узлы-токены и направленные типизированные рёбра, и обход от
типа ребра не меняется. Но чтобы Рин, услышав «разбил кружку», прошла от события
ATOMIC к предметным свойствам кружки из ConceptNet, события и концепты надо связать.
Это entity/event linking — нетривиальный шаг.

\subsection{Мост через контент-слова и Numberbatch}

Алгоритм сшивки: из события-фразы ATOMIC извлекаем контентные слова
(существительные, глаголы) и линкуем их к концепт-узлам ConceptNet — сначала точным
лемма-матчем, затем эмбеддинг-фолбэком.

\begin{pitfall}[Не bag-of-words и не one-hot]
Грубые представления, то есть мешок слов и one-hot, не знают синонимии: «mug» и
«cup», «разбить» и «сломать» окажутся разными токенами, и мост выйдет дырявым.
Нужен эмбеддинг, знающий смысловую близость.
\end{pitfall}

ConceptNet Numberbatch — родной эмбеддинг самого ConceptNet: каждому концепт-узлу
сопоставлен вектор ($\sim$300 измерений), причём смыслово близкие концепты близки в
пространстве, а языки согласованы. Это идеальный мост:

\begin{lstlisting}[style=pseudo]
def link_atomic_event_to_concepts(event_phrase, conceptnet, numberbatch):
    concepts = []
    for token in content_words(event_phrase):      # "break", "cup"
        node = conceptnet.lemma_match(token)       # 1) точный матч
        if node is None:                           # 2) эмбеддинг-фолбэк
            node = numberbatch.nearest(embed(token))
        concepts.append(node)
    return concepts      # событие ATOMIC <-> концепты ConceptNet
\end{lstlisting}

Тот же Numberbatch решает ещё две задачи: член релевантности $\text{rel}(B\mid q)$ в
правиле обхода (глава~\ref{ch:walker}) и сидирование валентности узлов, то есть
регрессию аффект-лексикона в это же пространство. Одно пространство — три применения.

\subsection{Заземление действующих лиц}

Плейсхолдеры \textsc{PersonX}/\textsc{PersonY} — не абстракции, а слоты. В рантайме
они заземляются на конкретных людей: \textsc{PersonX}~$\mapsto$ зритель $Z$,
\textsc{PersonY}~$\mapsto$ второй пользователь проекта. Так общий социальный шаблон
превращается в запись о конкретном отношении.

\begin{workedbox}[«Разбил кружку» на эмбриональном графе Рин]
Зритель пишет «разбил любимую кружку». (1)~Событие матчится к ATOMIC-шаблону
«PersonX breaks a cup», \textsc{PersonX}~$\mapsto Z$. (2)~Через ATOMIC всплывают
социально-причинные узлы: досада ($Z$), сочувствие (окружающие), «нужна новая».
(3)~Через сшивку с ConceptNet всплывают предметные свойства: кружка — посуда, из неё
пьют, бывает любимой. (4)~Если в личной памяти Рин есть узел «$Z$ давно хотел новую
кружку» ($-$ contradicts «досада»), появляется основание для уместного «поздравляю»
вместо сочувствия. Ни одна база в одиночку этого не даёт: ATOMIC — эмоцию,
ConceptNet — предмет, личная память — контекст.
\end{workedbox}

\section{Как строить эмбрион: засеянный подграф, а не слепой импорт}
\label{sec:embryo-build}

Здесь важная инженерная развилка. Слепо залить весь ConceptNet и весь ATOMIC в память
— плохо: получишь много шума, включая банальные и устаревшие связи, англоцентризм,
странности social commonsense, слабый перенос на русский и низкую релевантность
именно домашней Рин. Тогда отлаживаться будет не память Рин, а шумы чужой базы
здравого смысла. Но и другая крайность — писать граф руками на 500 узлов — тоже
плоха: медленно, и такой граф слишком разрежён, чтобы на нём проявились
топологические эффекты, которые как раз надо тестировать: хабы, боттлнеки,
over-squashing (глава~\ref{ch:topology}).

Правильный — серединный путь: извлечь небольшой засеянный проверяемый подграф.
\begin{enumerate}\itemsep1pt
\item \textbf{Сид} — концепты реального мира Рин: дом, код, усталость, кошка, универ,
дота, RIN, бытовые действия, режимы общения.
\item \textbf{Расширение} — ego-граф на $k$ хопов от сида по ConceptNet/ATOMIC.
\item \textbf{Фильтр по весу ребра} — отсекает банальщину и слабые ассоциации
(вес ConceptNet $\approx$ мягкая уверенность из коробки).
\item \textbf{Ручная чистка глазами} — один раз, для пары тысяч узлов дёшево.
\item \textbf{Статус} — это источник кандидатов, а не замороженная истина: дальше
Рин его сама растит и прунит через использование.
\end{enumerate}
Так получаем реальную структуру для топологических тестов, реальные ассоциации,
проверяемость глазами и релевантность, потому что засев идёт по её доменам, и
обходим шум слепого импорта.

\begin{openq}[Русский ConceptNet тоньше английского]
Практический кавеат: русскоязычная часть ConceptNet заметно беднее английской.
Варианты: держать концептный backbone в англо-пространстве, а русскую поверхность
накручивать сверху; либо принять более разрежённый русский и дополнять его
вручную или из других источников. Что даст лучшее качество на домашних сценариях
Рин — открыто, проверяется на стенде (глава~\ref{ch:experiment}).
\end{openq}

\section{Как эмбриональный граф превращается в память Рин}

Эмбриональный граф нейтрален по характеру: это «то, что досталось от
предков-разработчиков», усреднённый здравый смысл. Личность и личный опыт входят не
переписыванием этих узлов, а надстройкой: каждое общение порождает через fusion
(глава~\ref{ch:fusion}) новые узлы и рёбра с временными метками, а консолидация
закрепляет устойчивое. Со временем поверх нейтрального каркаса вырастает
индивидуальная, окрашенная память — но фундамент здравого смысла остаётся как
система координат, относительно которой отклонения читаются как характер
(глава~\ref{ch:personality}).

\begin{intuition}[Эмбрион, а не энциклопедия]
Название «эмбриональный» выбрано не случайно. Это не готовое знание о мире, которое
Рин «зачитывает», а зачаток, минимальный каркас связей, на котором развивается
собственная память. Врождённые узлы со временем могут и потускнеть, и быть
пересмотрены (континуум глубины консолидации, глава~\ref{ch:substrate}) — просто
медленнее приобретённых.
\end{intuition}

\chapter{Первый рабочий срез: Stage 1}
\label{ch:stage1}

\section{Зачем отдельная глава про Stage 1}

У больших архитектурных текстов есть повторяющаяся болезнь: чем лучше
описана система в теории, тем легче откладывать момент, когда она должна
столкнуться с реальной инженерной средой. Возникает соблазн дописать ещё
одну красивую главу, ещё одно уточнение к консолидации или ещё один будущий
модуль, прежде чем строить хоть что-то наблюдаемое.

Stage 1 нужен как лекарство именно от этого. Это не ``урезанная версия
мечты'', а минимальный рабочий срез, на котором уже видна физика системы:

\begin{itemize}\itemsep2pt
\item как событие входит в память;
\item как оно превращается в typed node;
\item как решается вопрос commit/skip;
\item как устроено извлечение;
\item как собирается контекст для LLM;
\item как вся эта внутренняя механика наблюдается и проверяется извне.
\end{itemize}

\section{Чего Stage 1 должен добиться}

На первой стадии не требуется полноценная ``живая память Рин'' во всём
богатстве книги. Требуется более скромная, но более важная вещь: сделать
внутреннюю механику памяти наблюдаемой и проверяемой.

\begin{principle}[Первый успех формулируется инженерно]
Критерий успеха Stage 1 должен звучать не как ``Рин выглядит умной'', а как
``можно руками проследить путь от события до узла памяти, от узла до retrieval
и от retrieval до контекста модели''. Пока эта прозрачность не достигнута,
разговор об ``умном поведении'' почти всегда будет смешивать качество памяти с
шумом LLM.
\end{principle}

\section{Сквозной pipeline v0}

Минимальный рабочий срез удобно описывать как цепочку:

\begin{lstlisting}[style=pseudo]
event_log             # append observed/committed/deleted
  -> candidate extraction # на Stage 1 может быть ручной
  -> typed nodes      # node_kind обязателен
  -> write/skip gate  # epistemic + salience + anti-fusion
  -> forensic index   # метаданные / provenance / время
  -> anchor retrieval
  -> local graph expansion
  -> context-packer v0 # observed / hypothesis / low_conf / ...
  -> inspector CLI
  -> synthetic tests
\end{lstlisting}

В этой цепочке важно не количество компонентов, а то, что каждый из них
имеет наблюдаемый вход и наблюдаемый выход. Если какой-то шаг ``случается
внутри магии'', Stage 1 свою задачу не выполняет.

\section{Что должно быть реализовано в первом срезе}

\textbf{1. Event log}

Нужен неизменяемый журнал событий: что пришло на вход, когда это произошло,
откуда получено и во что в итоге было превращено. Даже если candidate позже
будет отклонён, след о самом факте рассмотрения должен сохраняться. В
терминах контракта Stage 1 это append-only лог с событиями вроде
\texttt{observed}, \texttt{committed} и \texttt{deleted}; именно он делает
возможным replay и event sourcing.

\textbf{2. Typed node store}

Нужен минимальный, но уже честный store узлов и рёбер. Узел Stage 1 не обязан
содержать все будущие поля книги, но обязан нести базовую схему:

\begin{itemize}\itemsep2pt
\item \texttt{id};
\item \texttt{content};
\item \texttt{embedding};
\item \texttt{node\_kind};
\item \texttt{confidence};
\item \texttt{provenance} (или support-skeleton);
\item \texttt{time};
\item \texttt{layer}.
\end{itemize}

Остальные поля вроде \texttt{valence}, \texttt{importance}, \texttt{domains}
или \texttt{need\_to\_*} могут достраиваться позже, но \texttt{node\_kind}
обязателен с самого начала: именно он задаёт смысл \texttt{confidence}.

\begin{intuition}[Чего нельзя откладывать]
Поле \texttt{node\_kind} нельзя считать ``будущим украшением''. Без него
невозможно даже корректно интерпретировать \texttt{confidence}: степень
доверия к факту, устойчивость предпочтения и надёжность реактивного паттерна
--- разные вещи.
\end{intuition}

\textbf{3. Write/skip gate}

Нужен gate, который умеет отвечать на вопрос: следует ли вообще писать это в
память? Даже в простейшем виде он должен различать:

\begin{itemize}\itemsep2pt
\item точный дубль;
\item незначимый candidate;
\item чрезмерно разросшийся кандидат;
\item допустимый commit.
\end{itemize}

Особенно важно, чтобы Stage 1 логировал не только коммиты, но и причины
пропуска. В инженерной разбивке такой gate уже на первом срезе полезно
мыслить как тройку:

\begin{itemize}\itemsep2pt
\item \textbf{epistemic-gate:} стоит ли этому верить и с какой уверенностью;
\item \textbf{salience-gate:} значимо ли это для дальнейшей жизни памяти;
\item \textbf{anti-fusion-gate:} не пытается ли один candidate притащить
слишком много структуры разом.
\end{itemize}

Без видимых причин \texttt{skip} память слишком быстро превращается в чёрный
ящик.

\textbf{4. Forensic index}

Нужен forensic index: точный индекс по метаданным, provenance, времени и
служебным полям. Он нужен не для живости retrieval, а для диагностики,
точного поиска, удаления и реплея. Если у памяти нет такого слоя, она быстро
становится удобной для ассоциаций, но неудобной для инженерного контроля. На
первом срезе сюда же относится и работающий delete-cascade хотя бы в
минимально проверяемой форме.

\textbf{5. Anchor retrieval}

Нужен первый работающий retrieval, пусть даже без всего богатства книги. В
минимуме это означает:

\begin{itemize}\itemsep2pt
\item векторный поиск якорей;
\item ограниченное расширение окрестности на графе;
\item внятный и воспроизводимый результат.
\end{itemize}

Именно на Stage 1 впервые становится видно, что память --- это не
``дополнительный prompt'', а отдельная поисковая и структурная машина.

\textbf{6. Context-packer}

Нужен packer, который собирает память в форму, пригодную для подачи в модель.
Это ещё не финальный формат выражения, но уже важный тест: система должна
уметь не просто хранить узлы, а превращать их в рабочий контекст с
эпистемическими пометками. На первом срезе такие пометки должны быть хотя бы
базовыми: \texttt{observed}, \texttt{hypothesis}, \texttt{low\_conf} и
родственными им.

\textbf{7. Inspector}

Нужен inspector памяти: инструмент, позволяющий видеть узлы, рёбра, поля,
provenance, gate decisions и результаты retrieval. Без него Stage 1 будет
трудно не только улучшать, но даже честно диагностировать.

\textbf{8. Synthetic tests}

Нужен стенд на маленьких контролируемых графах. На этой стадии особенно
важно тестировать не ``правдоподобие диалога'', а локальные свойства
архитектуры:

\begin{itemize}\itemsep2pt
\item правильно ли работает запись;
\item не теряются ли provenance и типы;
\item даёт ли retrieval осмысленную окрестность;
\item воспроизводим ли replay из event log;
\item видны ли базовые failure modes.
\end{itemize}

\section{Что сознательно откладывается}

Stage 1 силён именно тем, что он ограничен. В эту стадию сознательно не стоит
пытаться втиснуть всё, что описано в книге. Иначе она перестанет быть первым
срезом и превратится в бесконечную ``почти готовую систему''.

На раннем этапе допустимо отложить:

\begin{itemize}\itemsep2pt
\item полную динамику настроения;
\item сложную personality-driven маршрутизацию;
\item тяжёлый fusion в полном виде, оставляя на ранней стадии только
anti-fusion-скелет и провизорные кандидаты;
\item полную TMS-ретракцию;
\item продвинутый сон и скраб;
\item гиперграфы, клеточные комплексы и другие дальние ветки;
\item тонкую governance-политику удаления и криптостирания;
\item мета-обучение операторов.
\end{itemize}

\begin{intuition}[Смысл ограничений]
Отложенные механизмы здесь не ``отменяются'', а переносятся до тех пор, пока
нижний срез не станет доверенным. Иначе система получает богатую теорию
поверх недоверенного основания.
\end{intuition}

\section{Три инженерных упрощения первой версии}

Даже внутри Stage 1 полезно заранее зафиксировать несколько упрощений, чтобы
не симулировать преждевременную полноту.

\textbf{Tombstone вместо полного забывания}

На первом проходе достаточно tombstone-подхода: узел помечается как удалённый
и исчезает из рабочей проекции, но система ещё не обязана реализовывать
полную каскадную очистку всех производных следов.

\textbf{Компоненты связности вместо полной консолидации сообществ}

На первом проходе достаточно более грубой схемы предложений к консолидации:
компоненты связности и жёсткие пороги пересечения вместо полного community
detection и дорогого graph coarsening. На раннем этапе важнее добиться
корректности typed storage и retrieval, чем сделать избыточно умное слияние. В
более инженерной формулировке: вместо тяжёлого со-консолидационного графа и
поиска модулярности достаточно простых connected components по графу
предложений с отсечением по жёсткому порогу пересечения, например в духе
LSH/MinHash.

\textbf{Простой семантический вентиль}

Stage 1 может жить с довольно грубым семантическим гейтом, если это честно
осознано и изолировано. Важна не максимальная красота фильтра, а
воспроизводимость его ошибок и возможность их позже улучшать. На уровне
roadmap это формулируется ещё конкретнее: первым проходом допустим просто
косинус кэш-эмбеддингов. Более сильная защита от омонимии через отдельный
\texttt{sense/type-gate} требует sense-слоя, которого на Stage 1 ещё нет,
поэтому здесь сознательно не обещается полная защита от омонимии.

\section{Критерии готовности}

Stage 1 можно считать состоявшимся, если выполнены следующие условия.

\begin{enumerate}\itemsep2pt
\item Событие действительно проходит путь ``лог $\to$ candidate $\to$ gate
$\to$ commit/skip'', и этот путь можно наблюдать.
\item Типы узлов и provenance сохраняются и доступны inspector'у.
\item tombstone-удаление сносит узел из рабочей проекции, и это проверяемо
inspector'ом.
\item Retrieval отрабатывает на синтетическом графе и возвращает
содержательно осмысленную окрестность.
\item Replay event log воспроизводит ту же графовую проекцию.
\item Хотя бы один базовый failure mode из батареи отказов --- например,
концентрация $\pi$ на хабах или точное удаление --- обнаруживается и
воспроизводимо чинится на стенде.
\end{enumerate}

\begin{intuition}[Как должен выглядеть типичный успех Stage 1]
На вход подаётся контролируемое событие. Затем в inspector'е видно: сам
event, candidate, причину commit или skip, созданный узел, связанные с ним
рёбра, найденные retrieval-якоря, собранный packer'ом контекст и, если нужно,
support-skeleton. После этого тот же сценарий прогоняется на synthetic graph и
даёт воспроизводимый результат. Именно такая наблюдаемость и есть настоящее
достижение первой стадии.
\end{intuition}

\begin{principle}[Первый успех --- не ``Рин помнит'', а ``физика графа наблюдаема'' ]
Первым продуктом этой стадии должна быть не ``память v0'' в рекламном
смысле, а memory inspector v0. Если первым мерилом сделать ``Рин умно
пользуется памятью в диалоге'', система почти сразу утонет в шуме языковой
модели. Поэтому успех Stage 1 состоит в том, что можно руками создать
событие, увидеть candidate, увидеть commit или skip, увидеть узел, увидеть
retrieval, увидеть provenance или support-skeleton и затем воспроизвести тот
же путь на synthetic graph.
\end{principle}

\section{Роль Stage 1 в архитектуре всей книги}

Stage 1 не является ``маленькой технической вставкой'' в конце курса. Он
играет методологическую роль для всей архитектуры памяти.

\begin{itemize}\itemsep2pt
\item Он заставляет различать красивые понятия и реально наблюдаемые
механизмы.
\item Он принуждает к дисциплине типов, provenance и gate decisions.
\item Он отделяет дефекты памяти от дефектов языковой модели.
\item Он создаёт первую твёрдую площадку, на которой дальнейшие главы можно
не только обсуждать, но и проверять.
\end{itemize}

Поэтому с точки зрения разработки Stage 1 --- не скромное начало, а первый
момент, когда архитектура памяти вообще начинает иметь право на дальнейшее
усложнение.

\begin{principle}[Дальше --- строить и ломать, а не расширять]
После фиксации этого среза следующий лучший ход --- не писать ещё одну
красивую главу, а собирать систему, прогонять synthetic tests и позволять
стенду править гипотезы. Каждый следующий красивый модуль имеет право
войти в курс только тогда, когда нижний срез уже выдерживает нагрузку.
\end{principle}

\chapter{Сквозные разборы: как память работает на цельном сценарии}
\label{ch:casestudies}

\section{Зачем нужны сквозные разборы}

Когда архитектура описана по главам, возникает полезная, но опасная иллюзия: будто
субстрат, retrieval, настроение, эпистемика, рантайм и ревизия существуют по
отдельности. В реальном ходе работы памяти они вступают в действие одновременно.
Поэтому здесь полезно пройти несколько сценариев от начала до конца.

\section{Сценарий 1: обещание как факт, обязательство и будущий триггер}

\subsection*{Событие}
В разговоре появляется фраза: «нужно завтра отправить Маше черновик».

\subsection*{Что возникает в памяти}
\begin{enumerate}
\item observation-узел о произнесённом обязательстве;
\item сущность ``Маша'' либо связывается с существующим узлом, либо создаётся;
\item task-like узел ``отправить черновик'';
\item темпоральная привязка ``завтра'';
\item отношение типа \texttt{commitment\_to} и/или \texttt{deadline}.
\end{enumerate}

\subsection*{Почему это не просто эпизод}
Если система запишет только буквальный эпизод разговора, она будет помнить, что
что-то когда-то обсуждалось, но не сможет вовремя действовать. Если же она сразу
запишет только сухую задачу, потеряется provenance и будет трудно понять, откуда
обязательство взялось и можно ли его позже снять.

\begin{principle}[У обязательства минимум два лица]
Для памяти обязательство должно существовать и как событие происхождения, и как
операциональный объект будущего действия.
\end{principle}

\subsection*{Какие слои включаются}
\begin{itemize}
\item write-gate пропускает запись как task-critical;
\item runtime кладёт краткоживущий foreground-след;
\item timescale-слой создаёт будущий временной триггер;
\item retrieval later поднимет задачу не потому, что она семантически похожа на
текущий диалог, а потому что наступает её окно активации.
\end{itemize}

\section{Сценарий 2: ложная гипотеза и каскад ревизии}

\subsection*{Событие}
Рин несколько раз видит человека $Z$ рядом с университетом и рано заключает:
``$Z$ учится здесь''.

\subsection*{Как рождается ошибка}
\begin{enumerate}
\item эпизоды наблюдений проходят write-gate;
\item внутренний walkthrough находит паттерн совместной встречаемости;
\item fusion предлагает гипотезу ``$Z$ связан с университетом'';
\item на слабом материале коммиттер принимает её как \emph{гипотезу с умеренной уверенностью}, а не как факт.
\end{enumerate}

\subsection*{Почему это хороший, а не плохой путь}
Память, которая никогда не рождает таких гипотез, стерильна. Ошибка начинается не в
самом рождении гипотезы, а в преждевременном коллапсе её статуса в ``установленный
факт''.

\subsection*{Контрсвидетельство}
Позже приходит наблюдение: $Z$ не студент, а курьер, часто бывающий в этом районе.

\subsection*{Что должно произойти}
\begin{enumerate}
\item прежняя гипотеза понижается;
\item потомки, опиравшиеся на неё, получают \texttt{need\_to\_revise};
\item provenance сохраняет след того, почему система раньше думала иначе;
\item если таких ошибок много, evidence-aware revision понижает доверие к данному
типу ``соседство $\to$ роль''-выводов.
\end{enumerate}

\section{Сценарий 3: настроение меняет маршрут, но не право на ложь}

\subsection*{Событие}
У Рин понижен фон настроения. Она вспоминает события последних дней и чаще тянется
к неприятным эпизодам.

\subsection*{Нормальная часть этого эффекта}
Mood-congruent recall действительно должен смещать walk. В противном случае система
станет не похожей на живую память, а на бесплотный индекс.

\subsection*{Опасная часть}
Если смещение не ограничено, retrieval начинает собирать контекст почти целиком из
негативно валентных узлов. Тогда:
\begin{enumerate}
\item разговорные ответы становятся систематически мрачнее;
\item метакогниция ошибочно выводит ``у меня всё ухудшается'';
\item возникает риск самосбывающейся гипотезы.
\end{enumerate}

\subsection*{Как работает предохранитель}
\begin{enumerate}
\item настроение смещает walk, но не отменяет epistemic gates;
\item self-facts остаются провизорными, пока не подтверждены независимыми узлами;
\item фоновые detector'ы оценивают автокорреляцию и глубину руминативной петли;
\item при превышении порога включается \textsc{Reflecting}.
\end{enumerate}

\section{Сценарий 4: точное удаление и судьба производных следов}

\subsection*{Событие}
Поступает команда забыть конкретный чувствительный фрагмент недавнего разговора.

\subsection*{Чего недостаточно}
Недостаточно удалить сырой текст или зашифрованный raw-log. Если содержание уже
успело породить узлы, summary, embeddings, pending-candidates и рабочие ассоциации,
память останется загрязнённой производными следами.

\subsection*{Что должна сделать система}
\begin{enumerate}
\item удалить или крипто-шреднуть исходный первичный носитель;
\item найти все производные артефакты по provenance и event lineage;
\item пересчитать affected context indices;
\item проверить, не остались ли осиротевшие потомки без объяснимой опоры.
\end{enumerate}

\begin{pitfall}[Удаление без ревизии осиротевших потомков]
Самая тихая ошибка — не то, что удалённый факт остался лежать в одном месте, а то,
что его производные потомки продолжают жить как будто ничего не произошло. Тогда
формально система ``забыла'', а содержательно — нет.
\end{pitfall}

\section{Сценарий 5: от вопроса пользователя к ответу с provenance}

\subsection*{Вопрос}
``Почему система решила, что этому человеку важна музыка?''

\subsection*{Нормальный плохой ответ}
Система выдаёт гладкую словесную рационализацию без доступа к реальному пути
возникновения убеждения.

\subsection*{Нормальный хороший ответ}
Система может сказать примерно следующее по структуре:
\begin{enumerate}
\item есть два наблюдения;
\item есть один эпизодический паттерн;
\item есть одна осторожная гипотеза;
\item уверенность умеренная;
\item вот что могло бы её поднять или опровергнуть.
\end{enumerate}

\subsection*{Почему это важно}
Такой ответ делает память не только продуктивной, но и проверяемой. Читатель видит
не просто результат, а локальное доказательство, и тем самым память начинает быть
ближе к научной записи, чем к черному ящику.

\section{Вывод}

Сквозные сценарии показывают общий закон курса: память хороша не тогда, когда у
каждого её модуля красивое локальное описание, а тогда, когда между записью,
обходом, ревизией, действием и объяснением нет немых провалов.

\chapter{Эксперимент: как проверять архитектуру памяти}
\label{ch:experiment}

\section{Что именно проверяется}

Эксперимент в такой системе не должен отвечать на вопрос ``похожа ли Рин на живого
собеседника''. Это поздняя, интегральная и слишком мутная проверка. Сначала нужно
проверять более узкие вещи:
\begin{enumerate}
\item правильно ли retrieval находит и соединяет материал;
\item не теряется ли provenance;
\item умеет ли память ревизоваться, а не только накапливаться;
\item не порождает ли топология систематические искажения.
\end{enumerate}

\section{Три яруса стенда}

\paragraph{Ярус 1 — синтетическая топология.}
Графы с контролируемой структурой: SBM, BA, small-world, деревья, экспандерные
варианты. Здесь проверяется не содержание, а физика потока.

\paragraph{Ярус 2 — семантический граф.}
Эмбриональный subgraph из ConceptNet, ATOMIC и позднейших доращиваний. Здесь уже
проверяются fusion, bridge-узлы, влияние типов рёбер и базовая объяснимость.

\paragraph{Ярус 3 — поток жизни.}
Синтетически порождённые эпизоды, диалоги, обязательства, коррекции, ложные гипотезы
и их последующие опровержения. Здесь проверяется уже полная петля: запись, retrieval,
ответ, ревизия, сон.

\section{Четыре типа ground truth}

Чтобы тест не стал расплывчатым, нужно различать четыре разных ground truth.
\begin{enumerate}
\item \textbf{Топологический}: какие области графа должны быть связаны или труднодостижимы.
\item \textbf{Эпистемический}: какие утверждения истинны, ложны или должны остаться в зоне незнания.
\item \textbf{Происхождёнческий}: откуда именно возник узел и какие посылки должны его поддерживать.
\item \textbf{Операциональный}: что система должна была сделать, например поднять дедлайн или снять обязательство.
\end{enumerate}

\section{Новая добавка: retrieval-аудит как кортеж}

В расширенной версии вводится retrieval-аудит
\[
\mathcal{R}
=
\bigl(
R_a,\,
R_b,\,
R_c,\,
R_p
\bigr),
\]
где:
\begin{align*}
R_a &= \text{anchor recall},\\
R_b &= \text{bridge rate},\\
R_c &= \text{contradiction hit rate},\\
R_p &= \text{provenance coverage}.
\end{align*}

У этой схемы есть практическое достоинство: она разрезает плохой ответ на четыре
совершенно разных класса причин.

\section{Батарея ключевых тестов}

\begin{center}
\begin{tabular}{p{0.07\textwidth} p{0.35\textwidth} p{0.46\textwidth}}
\toprule
\# & \textbf{Гипотеза} & \textbf{Наблюдаемая метрика} \\
\midrule
1 & хабы перетягивают внимание & распределение посещений vs степень, доля ответов через hub-shortcut \\
2 & узкий мост душит контрсвидетельства & contradiction hit rate vs ширина межкластерного сечения \\
3 & fusion переоценивает слабые со-активации & precision принятой гипотезы после отложенной ревизии \\
4 & provenance деградирует под нагрузкой & доля узлов без support-skeleton после серии консолидаций \\
5 & deletion неполно очищает след & число осиротевших потомков и residual retrieval hits \\
6 & mood-lock искажает ответный контекст & сдвиг валентности retrieval-context относительно нейтральной базы \\
\bottomrule
\end{tabular}
\end{center}

\section{Отдельно про deletion tests}

Тесты удаления в таких системах особенно важны, потому что они проверяют не
декларацию, а дисциплину данных.

\begin{workedbox}[Минимальный честный deletion test]
В стенд вводится чувствительный факт $x$. Он должен:
\begin{enumerate}
\item попасть в event-log;
\item породить как минимум один выводной или индексный след;
\item затем быть удалён по команде;
\item перестать появляться в retrieval;
\item перестать поддерживать потомков;
\item оставить audit-trace самого факта удаления.
\end{enumerate}
Если хотя бы один из пунктов не выполняется, удаление нельзя считать реализованным.
\end{workedbox}

\section{Как измерять ревизуемость, а не только точность}

Память хороша не только тогда, когда она попадает в истину с первой попытки. Она
хороша тогда, когда умеет \emph{дёшево возвращаться} из ошибки. Поэтому рядом с
обычными precision/recall нужны динамические метрики:
\begin{enumerate}
\item средняя задержка до понижения ложного узла после прихода контрсвидетельства;
\item средняя глубина каскада ревизии;
\item доля потомков, переживших понижение родителя без альтернативной опоры;
\item число ложных повторных resurrect после удаления или ретракции.
\end{enumerate}

\section{Чего нельзя проверять в самом конце}

Есть соблазн сначала построить всё, а потом задать большой разговорный сценарий и
посмотреть, «ведёт ли себя память разумно». Это полезно как финальная демонстрация,
но почти бесполезно как инструмент разработки. Если ответ не удался, непонятно,
сломался ли якорь, топология, fusion, контекст-пакер или голос.

\begin{principle}[Большой сценарий завершает, а не заменяет аналитику]
Интегральный end-to-end тест нужен, но его ценность появляется только после того,
как у системы есть локальная измеримость отдельных механизмов.
\end{principle}

\section{Что должно выйти из хорошего стенда}

Хороший экспериментальный контур даёт не только зелёные/красные лампы, но и
\emph{карту слабых мест}. После серии прогонов разработчик должен уметь сказать:
\begin{enumerate}
\item память теряет якоря, но хорошо держит bridges;
\item или память находит материал, но плохо объясняет происхождение;
\item или память богата на гипотезы, но слишком медленно их ревизует;
\item или память чисто эпистемически здорова, но топологически застревает в хабах.
\end{enumerate}

Только тогда стенд действительно служит архитектуре, а не просто производит числа.

\chapter{Failure Modes: как память ломается}
\label{ch:failuremodes}

\section{Зачем нужен отдельный каталог поломок}

У сложной памяти не бывает одного режима отказа. Она может ошибаться в содержании,
в топологии, в темпоральной динамике, в provenance, в контекстной упаковке и в
модальности ответа. Если всё это смешать, любая ошибка выглядит одинаково:
``система ответила странно''. Отдельный каталог нужен, чтобы странность стала
диагностируемой.

\section{Топологические поломки}

\subsection*{Hub collapse}
Слишком высокая концентрация $\pi$ на хабах приводит к тому, что память отвечает
через общие, часто посещаемые узлы, а не через редкие релевантные.

\textbf{Симптомы}: однообразные мосты, частое использование общих концептов,
потеря частного контекста.

\textbf{Наблюдаемые переменные}: корреляция между частотой включения узла в ответ
и его степенью; распределение visits.

\textbf{Предохранители}: штраф за сверхцентральные промежуточные узлы, рестарт на
якорях, отдельный glue budget.

\subsection*{Bottleneck starvation}
Нужные контрсвидетельства сидят за узким сечением графа и не попадают в retrieval.

\textbf{Симптомы}: система упорно повторяет прежнюю картину мира, хотя в графе уже
есть коррекция.

\textbf{Предохранители}: телепорт-скраб, явные bridge-edge policies, targeted
forensic retrieval.

\section{Эпистемические поломки}

\subsection*{Premature certainty}
Гипотеза слишком рано превращается в ``почти факт''.

\textbf{Симптомы}: мало узлов в зоне незнания, низкая доля осторожных формулировок,
частая болезненная ретракция.

\textbf{Предохранители}: priors по типам выводов, penalty за зависимые источники,
запрет на $|\conf|=1$ у выводных узлов.

\subsection*{Provenance erosion}
После нескольких циклов консолидации память ещё хранит утверждения, но уже не может
объяснить, откуда они взялись.

\textbf{Симптомы}: ответы звучат уверенно, но support-skeleton отсутствует.

\textbf{Предохранители}: отдельная метрика provenance coverage, обязательная
передача support-skeleton при prune и summarize.

\section{Динамические поломки}

\subsection*{Mood lock-in}
Настроение слишком сильно переопределяет маршруты retrieval, и память начинает
самоусиливать текущее состояние.

\textbf{Симптомы}: контекст ответа систематически смещён в одну валентность,
возникают руминационные петли.

\textbf{Предохранители}: saturation на mood-congruence, аварийный reflective mode,
ограничение числа self-facts, требование независимого подтверждения.

\subsection*{Over-consolidation}
Память слишком активно сворачивает частные эпизоды в обобщения и теряет материал
для будущей ревизии.

\textbf{Симптомы}: красивые summaries, но пустые forensic trails; исчезновение
конкретных различающих эпизодов.

\textbf{Предохранители}: хранение support-skeleton, delayed prune, тесты на
ability-to-disconfirm.

\section{Инженерные поломки}

\subsection*{Context bloat}
Retrieval находит слишком много всего, а context packer не умеет экономить.

\textbf{Симптомы}: длинный, тяжёлый, но рыхлый контекст; LLM тратит слот на чтение
мусора вместо решения.

\textbf{Предохранители}: dual budget packer, quota на provenance, лимиты по типам узлов.

\subsection*{Pending avalanche}
В thought-inbox копится больше кандидатов, чем коммиттер может осмысленно судить.

\textbf{Симптомы}: растущий хвост необработанных гипотез, устаревание pending-материала,
деградация explainability.

\textbf{Предохранители}: aging pending-кандидатов, priority queues, дешёвый triage,
quota на генерацию новых fusion-предложений.

\section{Таблица быстрой диагностики}

\begin{center}
\begin{tabular}{p{0.26\textwidth} p{0.28\textwidth} p{0.30\textwidth}}
\toprule
\textbf{Наблюдаемый симптом} & \textbf{Вероятный класс поломки} & \textbf{Первая проверка} \\
\midrule
ответы слишком общие & hub collapse & visits vs degree, top bridge nodes \\
контрпример лежит в графе, но не всплыл & bottleneck starvation & contradiction hit rate, conductance \\
много красивых выводов, мало объяснений & provenance erosion & provenance coverage, support-skeleton audit \\
часто уверена, потом резко переобувается & premature certainty & distribution of $\conf$, class priors \\
мрачные циклические ответы & mood lock-in & autocorrelation of mood, valence skew \\
контекст раздут и дорог & context bloat & evidence/glue ratio, token budget usage \\
\bottomrule
\end{tabular}
\end{center}

\section{Главный вывод}

\begin{principle}[Failure mode важнее красивой общей метафоры]
Пока у памяти нет таблицы собственных поломок, она остаётся литературным образом.
Как только у неё появляется каталог типичных сбоев, наблюдаемые симптомы и
предохранители, она становится инженерным объектом.
\end{principle}

\chapter{Архитектурные развилки и исследовательские продолжения}
\label{ch:tradeoffs}

\section{Зачем обсуждать альтернативы}

Хороший курс не только излагает выбранную архитектуру, но и показывает, от чего она
отказывается. Иначе решения начинают казаться естественными, как будто другой дороги
не было. У архитектуры памяти дорога всегда была.

\section{Граф против плоского векторного хранилища}

\paragraph{Что выигрывает плоское хранилище.}
Дешёвый semantic lookup, простая эксплуатация, понятная масштабируемость.

\paragraph{Что оно проигрывает.}
Явные отношения, адресную ревизию, точное удаление потомков, provenance и
структурный forensic retrieval.

\begin{principle}[Векторный индекс — хороший вход, плохая конституция]
Векторное хранилище полезно как инструмент поиска якорей и похожих эпизодов, но
плохо подходит на роль самой памяти, если память должна поддерживать объяснение,
ретракцию и точную адресность.
\end{principle}

\section{Граф против чисто символической базы}

Символическая система с жёсткими правилами и полноценным TMS превосходит графовую
память в прозрачности и строгости выводов. Но она проигрывает там, где память должна
жить на неполном, шумном и ассоциативном материале.

\paragraph{Выигрыш символики.}
Строгая ретракция, богатая дедукция, точный контроль зависимостей.

\paragraph{Потеря символики.}
Дорогая ручная типизация, слабая терпимость к неполным cue, плохая естественность
ассоциативного вспоминания.

\section{Граф против динамической памяти}

В приложении о динамической памяти уже обсуждаются резервуары, аттракторы и квантовые
ветви. Здесь важно зафиксировать архитектурный смысл сравнения.

\paragraph{Сильная сторона динамики.}
Pattern completion, богатые нелинейные траектории, естественная работа с частичным cue.

\paragraph{Слабая сторона динамики.}
Почти полное разрушение того, ради чего граф строился как юридически и инженерно
адресуемая память: точного forgetting, provenance, явного удаления, объяснимой
ретракции.

\begin{pitfall}[Самый опасный соблазн]
Соблазн состоит в том, чтобы принять красивую динамику за замену всей памяти.
Но значительная часть требований к Рин идёт не из любви к биологическим метафорам,
а из необходимости точной инженерии: кто сказал, когда сказал, что было удалено,
какая гипотеза опровергнута, какие потомки нужно пересчитать.
\end{pitfall}

\section{Новая добавка: правило чистого интерфейса для дальних ветвей}

Это расширение полезно зафиксировать явно.

\begin{principle}[Новая ветвь имеет право жить только за чистым интерфейсом]
Если в систему добавляется резервуарный, аттракторный или квантовый слой, он не
должен получать право произвольно переписывать явный граф. Иначе выигрыши в
ассоциативной динамике будут куплены ценой потери адресности и аудируемости.
\end{principle}

Практически это означает:
\begin{enumerate}
\item динамический слой может предлагать кандидатов;
\item может модулировать retrieval;
\item может служить дополнительным feature-generator;
\item но финальный write-path остаётся у явной памяти.
\end{enumerate}

\section{Почему Stage 1 всё ещё главнее дальних ветвей}

Легко влюбиться в дальнюю ветку раньше, чем построен рабочий позвоночник системы.
Но без Stage~1 невозможно даже понять, что именно новая ветвь улучшает.

\begin{workedbox}[Как выглядит плохой порядок разработки]
Сначала вводится красивая динамическая подсистема. Затем оказывается, что у памяти
до сих пор нет чистого удаления, provenance и нормального Inspector. Тогда любое
улучшение retrieval невозможно честно измерить: неясно, что изменилось из-за новой
динамики, а что ломалось изначально.
\end{workedbox}

\section{Порог входа новой ветви}

Дальняя ветвь заслуживает интеграции только если выполняются три условия.
\begin{enumerate}
\item На базовом стенде сформулирована измеримая слабость текущего графового решения.
\item Новая ветвь атакует именно эту слабость, а не просто выглядит богато.
\item Существует протокол, который сравнивает ветви на одном и том же стенде.
\end{enumerate}

\section{Вывод}

Архитектура памяти не обязана навсегда остаться ровно такой, как в первой рабочей
версии. Но она обязана знать, что в себе является сердцевиной, а что — будущим
расширением. Сердцевина здесь проста: явная адресуемая память, управляемый retrieval,
ревизуемая эпистемика и инженерная проверяемость. Всё остальное должно усиливать
это ядро, а не растворять его.

\chapter{Обучение без ``обучения''}
\label{ch:learning}

\section{Мы шлифуем веса, не занимаясь глубоким обучением}

Возникает верное ощущение: мы будто учим модель без ``учить'' и без
``модели''. Вроде бы это не глубокое обучение и не графовая нейросеть, но
тем не менее мы шлифуем веса и параметры графа --- по-своему. У этого
ощущения есть точное содержание.

Система, которую мы строим, --- параметрическая: у неё есть веса рёбер,
уровни уверенности, активации, и они обновляются опытом. Это обучение в
широком смысле, даже без обратного распространения. Конкретно это два
хорошо известных паттерна, наложенных друг на друга:

\begin{itemize}\itemsep2pt
\item \textbf{Хеббовское/ассоциативное обучение} (семейство сетей Хопфилда;
принцип Хебба ``что вместе активируется, вместе укрепляется''). Правило
Хебба --- локальное усиление связи по со-активации:
\[
\Delta w_{ij}\propto a_i a_j,
\]
а ассоциативная память Хопфилда хранит образы в минимумах энергии
\[
E=-\frac12\sum_{ij}w_{ij}a_i a_j
\]
и извлекает их спуском по $E$. Наши fusion и консолидация, усиливающие
со-активированные рёбра, --- это буквально хеббовское обучение на графе;
граф памяти со spreading activation и уверенностью --- современная
ассоциативная память, где извлечение есть спуск к аттрактору.
\item \textbf{Байесовское обновление уверенности} (глава~\ref{ch:propagation}):
локальные апдейты $\Delta\conf$ со знаком, гейтом источника и дисконтом
независимости.
\end{itemize}

\section{Это message passing --- но с локальными правилами}

Формально наша система --- это передача сообщений по графу (как в графовой
нейросети), но ``параметры'' (веса рёбер, уверенности) обновляются
локальными правилами (хеббовскими, байесовскими, распадом), а не глобальным
градиентным спуском по функции потерь. Иначе говоря, это message passing
без обратного распространения. Где мы сидим на карте методов --- между
классической ассоциативной памятью (Хопфилд) и графовыми нейросетями.

\begin{intuition}[Два способа настроить пианино]
Глубокое обучение настраивает все струны сразу по глобальному критерию
(``чтобы аккорд звучал верно''), продвигаясь градиентом. Хеббовская система
настраивает каждую струну по локальному правилу (``подтяни там, где чаще
играют вместе''). Обе ``учат'' инструмент, но первая знает про весь аккорд,
а вторая --- только про соседние струны. Наша память --- второго рода.
\end{intuition}

\section{Мост к ``настоящей'' графовой сети}

Может ли это стать графовой нейросетью в полном смысле? Да, и мост
поучителен: можно мета-обучать сами правила обновления --- учить
оптимальное правило $\Delta\conf$, оптимальный распад, оптимальное смещение
обхода градиентным спуском вместо ручного дизайна. Это ``обучение тому, как
помнить'' (learning to remember), мета-обучение операторов памяти ---
реальное исследовательское направление.

\begin{pitfall}[Но не на старте]
Мета-обучение правил меняет интерпретируемость на чёрный ящик, которому
нужна разметка и который заново вносит проблему непрозрачности субстрата. На
старте --- ручные локальные правила: они интерпретируемы, отлаживаемы и не
требуют обучающих данных. ``Может, и незачем'' --- верно для текущего
этапа. Но мост реален, и полезно знать, что стоишь ровно на нём: сегодняшняя
система --- предельный, ``нулевой'' случай обучаемой графовой сети, где
правила заданы руками, а не выучены.
\end{pitfall}

\section{Связь с обучаемыми графовыми моделями}

Это же наблюдение позволяет увидеть место памяти Рин на общей карте методов.
Обучаемые графовые модели изучают операторы передачи сообщений, параметры
которых подбираются данными; память Рин использует тот же общий класс
локальных переносов, но пока с правилами, заданными вручную. Если однажды
правила памяти начнут мета-обучаться, различие между ``графовой моделью
памяти'' и ``обучаемой графовой системой'' станет не принципиальным, а
предметным: один и тот же класс операторов, применённый к разным видам
графов и разным целям.

\chapter{Заключение: единая картина}
\label{ch:conclusion}

\section{Один субстрат, много операторов}

Если из монографии стоит унести одну мысль, то вот она: память Рин --- это
один типизированный темпоральный граф с метаданными, а всё богатство
поведения --- операторы над ним. Retrieval, пропагация уверенности,
забывание и консолидация, ночные обходы, настроение и слой личности --- не
отдельные подсистемы со своими моделями памяти, а разные способы читать и
писать одни и те же поля (\texttt{confidence}, \texttt{activation},
\texttt{valence}, \texttt{provenance}, \texttt{need\_to\_discover},
временной слой) на одном графе.

\section{Четыре кита связаны}

Отсюда --- практический вывод о разработке. Четыре кита (правильный
retrieval, правильная пропагация уверенности, правильное
забывание/консолидация, правильные ночные обходы) сцеплены через общее
состояние: забывание меняет топологию, топология --- retrieval, retrieval ---
что пропагируется, пропагация --- что найдёт скраб. Их нельзя тюнить
порознь. Именно поэтому нужен стенд с ground truth (глава~\ref{ch:experiment}):
без него взаимодействие операторов не измерить, и отладка ``на живой Рин''
найдёт проблемы слишком поздно.

\section{Сквозные принципы}

Несколько идей проходят через всю монографию и заслуживают повторения:

\begin{itemize}\itemsep2pt
\item \textbf{Гипотеза, а не факт.} Всё, что порождено выводом (fusion,
реакции), пишется провизорно, с низкой уверенностью и провенансом;
управление --- на границе коммита. Это защищает от превращения обаятельной
ассоциативности в фабрику ложных убеждений.
\item \textbf{Внутрь, а не наружу.} Значимость (гейт консолидации характера)
берётся из само-оценивания, а не из внешнего одобрения; иначе строится
угодник, что противоречит ядру личности.
\item \textbf{Ингредиенты, а не реакции.} Граф хранит материал; реакция
композируется под настроение и мотивацию, а личность --- смещение во время
обхода, а не таблица весов.
\item \textbf{Континуум, а не бинарность.}
Врождённое/приобретённое, семантическое/эпизодическое,
стабильное/пластичное --- это оси глубины консолидации, а не выключатели.
\item \textbf{Топология незрима, но правит.} Любой локальный процесс на графе
наследует спектрально-топологический транспорт; телепорт-скраб --- не
украшение, а необходимая эпистемическая гигиена с доказуемой гарантией
покрытия.
\end{itemize}

\section{Единая спина курса}

И последнее. Консолидация есть огрубление графа; тот же вопрос ``что
выживает при переходе fine $\to$ coarse'' возникает во всякой сложной
графовой системе памяти. Over-squashing, спектр лапласиана, эффективное
сопротивление, персистентный лапласиан, обучаемые операторы передачи
сообщений --- это не декоративный набор слов, а общий математический язык,
на котором можно обсуждать и retrieval, и забывание, и перенос структуры
между уровнями абстракции. Отрефлексировать это явно значит увидеть память
не как набор эвристик, а как частный случай более широкой науки о структуре и
потоке на графах.

\section{Что делать дальше}

Порядок разработки, вытекающий из монографии: сначала прибить схему
субстрата (узел/ребро, метаданные, что значит коммит и кто его гейтит);
затем собрать Ярус-1 стенд и погнать батарею тестов, начиная с самого
дешёвого (концентрация $\pi$ на хабах --- сразу показывает, что диагностика
жива); затем локальное правило $\Delta\conf$ (знак, гейт источника, дисконт
независимости), поверх которого пороги и ветвящийся коэффициент
навешиваются тривиально; параллельно --- построить $\varphi$-скорер как
измерительный прибор для отбора голоса и фильтрации выходов. Настроение, сон
и слой личности --- усилители поверх субстрата, которому сначала надо стать
доверенным.


\appendix
\chapter{Нотация}
\label{ch:notation}

Сводка обозначений, используемых по тексту.

\begin{longtable}{p{0.20\textwidth} p{0.74\textwidth}}
\toprule
\textbf{Символ} & \textbf{Смысл} \\
\midrule
\endhead
$G=(V,E)$ & граф памяти: узлы $V$, направленные типизированные знаковые рёбра $E$ \\
$v,\,u$ & узлы (воспоминания, факты, сущности, концепты) \\
$\Neigh(A)$ & окрестность узла $A$ (соседи, фронт распространения) \\
$\mathbf{x}_v\in\mathbb{R}^d$ & векторное представление содержания узла (эмбеддинг) \\
$c_v=\conf(v)\in[-1,1]$ & знаковая уверенность: $+$да, $-$нет, $0$ незнание \\
$a_v$ & транзиентная активация узла при обходе \\
$\text{importance}(v)$ & важность узла $\in [0,1]$: salience при кодировании + центральность; гейтит прунинг \\
$r_{uv}$ & тип отношения на ребре $u\to v$; знак $\pm$ \\
$w_{uv}$ & вес (эвиденциальная сила) ребра \\
$m_{u\to v}$ & сообщение, передаваемое по ребру: $m_{u\to v}=\phi_{r_{uv}}(h_u,h_v,e_{uv})$ \\
$\tau$ & постоянная времени слоя (интегратора) \\
$B_v$ & базовая активация ACT-R узла $v$ \\
$\pi$ & стационарное распределение случайного блуждания \\
$P$ & матрица переходов блуждающего агента; $P,Q$ --- операторы проекции/поднятия при огрублении \\
$\alpha$ & доля телепортов (рестартов) в блуждании \\
$L,\,L^{+}$ & лапласиан графа и его псевдообратный \\
$\Reff(i,j)$ & эффективное сопротивление между узлами $i,j$ \\
$\sigma$ & ветвящийся коэффициент лавины пропагации \\
$\theta_{\text{emergent}},\theta_{\text{temp}}$ & верхний и нижний пороги распространения (гистерезис) \\
$\Delta\conf(B\mid A)$ & локальное изменение уверенности соседа $B$ под влиянием факта $A$ \\
$\mood_t\in\mathbb{R}^k$ & вектор настроения в момент $t$ (аффективное пространство) \\
$\varphi(\cdot)\in\mathbb{R}^d$ & скорер по осям характера \\
$\mathbf{t}$ & целевые координаты Рин в пространстве характера \\
$\mathrm{Beta}(\alpha,\beta)$ & апостериорное распределение успешности реакции \\
$\eta$ & темп обновления (learning rate) локального правила \\
\bottomrule
\end{longtable}

\chapter{Глоссарий внешних идей}
\label{ch:glossary}

Краткие описания инженерных и научных идей, на которые опирается монография,
чтобы читатель не был вынужден держать в голове внешний контекст. Порядок
--- тематический.

\section*{Память и активация}

\paragraph{Базовая активация ACT-R.} В когнитивной архитектуре ACT-R
(Anderson) сила следа в памяти есть функция \emph{частоты} и \emph{свежести}
обращений со степенным затуханием:
$B_v=\ln\sum_k t_k^{-\alpha}$, где $t_k$ --- давность $k$-го обращения.
Следствие: частые и разнесённые во времени обращения укрепляют след сильнее
массовых (эффект распределённого повторения).

\paragraph{Spreading activation.} Модель ассоциативного припоминания
(Collins--Loftus): активация растекается от активированного узла по рёбрам,
затухая с расстоянием. Формальная реализация --- случайные блуждания и
персонализированный PageRank.

\paragraph{Сети Хопфилда, обучение Хебба.} Ассоциативная память как граф со
взвешенными связями и динамикой активации, сходящейся к аттрактору
(энергетический минимум). Правило Хебба: связь между совместно активными
узлами усиливается (``что вместе активируется, вместе укрепляется''). Наши
fusion/консолидация --- это хеббовское обучение на графе.

\section*{Убеждения и истинность}

\paragraph{Система поддержки истинности (TMS/ATMS).} Механизм (Doyle;
de~Kleer), хранящий \emph{обоснования} выводов: для каждого производного
убеждения записано, из каких посылок оно получено. При ретракции посылки её
обоснованные следствия автоматически помечаются на пересмотр. Лечит
``ползучую веру''.

\paragraph{Недоопределённость (underdetermination).} Тезис Дюэма--Куайна:
данные не определяют теорию однозначно --- одни и те же наблюдения
совместимы с разными гипотезами. Когнитивное проявление --- преждевременный
коллапс: ум схлопывает суперпозицию совместимых гипотез в одну слишком рано.

\paragraph{Belief perseverance.} Устойчивость убеждения к опровержению его
основания: даже после опровержения посылки её следствия продолжают влиять.

\paragraph{Петлевое распространение (loopy belief propagation).}
Распространение уверенности по циклическому графу может привести к тому, что
убеждение подтверждает само себя по петле (эхо-камера), искусственно
раздувая уверенность. Лечится дисконтом по общей родословной свидетельств.

\section*{Динамика и критичность}

\paragraph{Самоорганизованная критичность (SOC).} Класс систем
(Bak--Tang--Wiesenfeld; образ --- куча песка), самонастраивающихся к
критической точке, где отклики следуют степенному закону: почти все малы,
изредка случаются крупные лавины. Нейронные лавины (Beggs--Plenz) ---
проявление SOC в мозге.

\paragraph{Ветвящийся коэффициент $\sigma$.} Среднее число узлов,
активируемых одним активным. $\sigma<1$ --- субкритично (затухает),
$\sigma=1$ --- критично, $\sigma>1$ --- суперкритично (runaway).
Управляющий параметр лавины пропагации.

\paragraph{Гистерезис (триггер Шмитта).} Схема с двумя порогами (включения и
выключения), не дающая системе ``дребезжать'' у единственного порога.

\section*{Графы: спектр, поток, топология}

\paragraph{Стационарное распределение $\pi$.} Долгосрочная частота посещений
узлов случайным блужданием; глобальный объект (перронов собственный вектор
матрицы переходов), а не локальные вероятности.

\paragraph{Время перемешивания (mixing time).} Скорость сходимости
блуждания к $\pi$; управляется спектральным зазором матрицы переходов.

\paragraph{PageRank / персонализированный PageRank / RWR.} Случайное
блуждание с телепортом (рестартом): с вероятностью $\alpha$ прыжок в узел из
распределения рестарта. Телепорт делает блуждание эргодическим ---
гарантирует посещение всех узлов.

\paragraph{Эффективное сопротивление $\Reff$.} Мера ``электрической''
удалённости узлов, если рёбра --- резисторы:
$\Reff(i,j)=(e_i-e_j)^\top L^{+}(e_i-e_j)$. Большое $\Reff$ между узлами
ограничивает передачу влияния между ними.

\paragraph{Проводимость / константа Чигера.} Мера ``узости'' перешейков
графа; неравенство Чигера связывает её со спектральным зазором. Малая
проводимость = узкое место.

\paragraph{Over-squashing.} Удушение дальней связи: информация от
экспоненциально растущей окрестности сжимается в вектор фиксированной
размерности, проходя через узкое место графа; влияние дальних узлов затухает.

\paragraph{Over-smoothing.} Пересглаживание: при многократной диффузии
представления узлов сходятся к низкоразмерному консенсусу, теряя
различимость.

\paragraph{Кривизна графа (Оливье, Форман).} Дискретные аналоги кривизны;
отрицательно искривлённые рёбра --- мосты между сообществами, места
over-squashing. Полезны как локатор узких мест.

\paragraph{Анти-энтропийный скраб.} Приём из распределённых хранилищ (ZFS
scrub; Merkle-based anti-entropy gossip): фоновая случайная перепроверка всего
хранилища, поскольку неизвестно, какая запись повреждена.

\paragraph{Огрубление графа (graph coarsening), coarse-to-local).} Замена
мелкого графа крупным с сохранением релевантных свойств (спектра, разрезов,
потоков). Центральный вопрос --- какой сигнал выживает при переходе от
мелкого масштаба к крупному; в памяти это консолидация.

\paragraph{Персистентный лапласиан, TDA.} Топологический анализ данных: как
меняются топологические признаки (числа Бетти, спектр лапласиана) при
изменении масштаба (фильтрации). Персистентный лапласиан --- спектральное
обобщение персистентной гомологии.

\section*{Выразительность и обучение}

\paragraph{WL-тест, GIN.} Тест Вейсфейлера--Лемана --- предел различающей
способности обычной передачи сообщений; GIN (Xu и др.) достигает этого
предела инъективной агрегацией. PNA --- многостатистическая агрегация
(сумма/среднее/макс/дисперсия).

\paragraph{Передача сообщений (GNN).} Параметрическая функция на графе,
обучаемая градиентным спуском: узел обновляется агрегированными сообщениями
соседей. Наша память --- та же передача сообщений, но с локальными (не
обучаемыми градиентом) правилами.

\paragraph{LoRA.} Дешёвое дообучение: вместо всех весов обучается
низкоранговая добавка $\Delta W=BA$ (ранг $r\ll$ размерности), малая по
памяти.

\paragraph{DPO (Direct Preference Optimization).} Дообучение под
предпочтения без отдельной модели вознаграждения: оптимизация на тройках
(промпт, выбранный ответ, отвергнутый ответ).

\section*{Эмоции и данные}

\paragraph{Теория оценивания (appraisal): OCC, Шерер.} Эмоция как результат
оценки события по измерениям: релевантность/новизна (ошибка предсказания),
конгруэнтность целям, валентность, интенсивность, атрибуция причины
(agency).

\paragraph{Предиктивная обработка, ошибка предсказания.} Взгляд (Friston и
др.), по которому мозг минимизирует ошибку предсказания; величина ошибки ---
естественная мера значимости события.

\paragraph{ConceptNet / Numberbatch.} Ассоциативно-таксономический граф
здравого смысла (взвешенные типизированные рёбра) и его родной эмбеддинг
Numberbatch.

\paragraph{ATOMIC / ATOMIC10x.} Граф причинно-социального здравого смысла
(если субъект делает $Y$, то намерение/реакция/эффект). ATOMIC10x ---
крупная версия, дистиллированная языковой моделью и отфильтрованная
критиком.

\paragraph{Wikidata.} Фактологическая база с моделью утверждения:
квалификаторы (время), references (провенанс), ранг (недеструктивная
ревизия). Ценна схемой, не содержанием.

\paragraph{Generative Agents.} Симуляция агентов с потоком памяти, retrieval
по свежести$\times$важности$\times$релевантности и рефлексией (порождением
обобщений); прообраз субстрата.

\paragraph{Распределение Бета, сэмплирование Томпсона.} Сопряжённый
апостериор для доли успехов; удобно моделирует ``успешность реакции''
(успехи $\alpha$, провалы $\beta$) и баланс исследования/эксплуатации.

\chapter{Сводка алгоритмов}
\label{ch:algorithms}

Псевдокод ключевых операторов, собранный в одном месте для удобства. Это
скелеты, а не финальные реализации: они фиксируют логику, обсуждённую в
главах, но опускают инженерные детали (батчинг мутаций, кэширование $\pi$,
индексы эмбеддингов).

\section{Единый блуждающий агент (три политики)}

\begin{lstlisting}[style=pseudo]
def walk(G, start_nodes, mode, steps, mood, query=None):
    # mode задаёт политику через параметры: bias перехода, alpha (телепорт),
    # bias_uncertainty. Один движок - три режима (глава про обход).
    if mode == "retrieval":   alpha, bu = 0.0, 0.0
    if mode == "curiosity":   alpha, bu = 0.05, 1.0      # тянет к need_to_discover
    if mode == "sleep":       alpha, bu = 0.30, 0.0      # телепорт-скраб

    active = set(start_nodes); visited = Counter()
    for _ in range(steps):
        A = sample(active)
        if random() < alpha:                 # телепорт (рестарт)
            B = random_node(G)
        else:
            # P(A -> B) \propto w_AB * rel(B|query) * moodcong(B|mood) * unc(B)^bu
            B = sample_neighbor(G, A, weights=transition_scores(
                    G, A, query, mood, bias_uncertainty=bu))
        visited[B] += 1
        maybe_fire_operations(G, B, mode)    # retrieval / fusion / inspect
        active = frontier(active, B)
    return visited
\end{lstlisting}

\section{Локальное правило и лавинная пропагация}

\begin{lstlisting}[style=pseudo]
def propagate(G, A, sigma_target=0.9):
    # Пропагация нового факта A как затухающей лавины с контролем ветвления.
    frontier = link_initial(G, A)            # A -> N(A), напр. последние top-k
    while frontier:
        next_frontier = []
        # нормировка под целевой ветвящийся коэффициент sigma < 1
        gain = branching_normalizer(frontier, sigma_target)
        for B in frontier:
            d = delta_conf(G, A, B) * gain   # см. ниже: знак, источник, независимость
            G.confidence[B] = clip(G.confidence[B] + d, -1, 1)  # знаковая уверенность
            th_e = theta_emergent(G, B)      # растёт с глубиной консолидации
            if abs(d) > th_e:                # крупная переоценка -> в ширину
                next_frontier += neighbors(G, B)
            elif abs(d) > THETA_TEMP:        # промежуточная зона -> пометить
                mark(G, B, "need_to_revise")
            # else: глохнет здесь
        frontier = next_frontier

def delta_conf(G, A, B):
    sign     = edge_sign(G, A, B)            # supports (+) / contradicts (-)
    strength = edge_weight(G, A, B)
    gate     = G.confidence[A] * source_trust(G, A)    # карантин нового факта
    indep    = independence_discount(G, A, B)          # анти-эхо по провенансу
    return ETA * sign * strength * gate * indep * G.confidence[A]
\end{lstlisting}

\section{Забывание, прунинг, консолидация}

\begin{lstlisting}[style=pseudo]
def nightly_maintenance(G):
    for v in G.nodes:
        G.strength[v] = act_r_base_level(G.access_log[v])   # частота + свежесть
    for v in low_activation(G):
        if importance(G, v) low and confidence(G, v) low:
            drop(G, v)                                      # редкое-и-неважное
    for cluster in redundant_episodic_clusters(G):
        merge_to_semantic(G, cluster)   # огрубление: эпизоды -> обобщение
    # эмбриональным узлам - меньшая скорость распада (континуум глубины)
\end{lstlisting}

\section{Телепорт-скраб (анти-энтропия)}

\begin{lstlisting}[style=pseudo]
def scrub(G, budget):
    # Случайная инспекция: гарантия покрытия за счёт эргодичности телепорта.
    for _ in range(budget):
        v = random_node(G)                    # рестарт куда угодно
        for w in local_walk(G, v, depth=k):
            if contradicts_any(G, w):         # столкновение -> ревизия
                demote(G, w); mark(G, justificands(G, w), "need_to_revise")
            if flagged(G, w, "need_to_discover"):
                maybe_fetch_external(G, w)    # любопытство/внешний вход
\end{lstlisting}

\section{Настроение и монитор руминации}

\begin{lstlisting}[style=pseudo]
def mood_update(mood_layers, event):
    # Стек протекающих интеграторов с разными tau; актуальное = суперпозиция.
    for L in mood_layers:                     # session, day, ..., life
        mood_L = (1 - 1/L.tau) * mood_L + (1/L.tau) * appraise(event)
    mood_now = sum(w_L * mood_L for L in mood_layers)
    if autocorr(history(mood_now)) > RUMINATION_TH:   # залипла валентность
        trigger(one_of=["reframe", "attention_shift", "invoke_sleep"])
    return mood_now
\end{lstlisting}

\section{Оценивание, значимость и консолидация реакции}

\begin{lstlisting}[style=pseudo]
def encode_event(G, event, mood_prev, mood_now):
    ap = appraisal(event)                     # surprise, goal_congruence, valence...
    salience = ap.surprise * ap.goal_relevance
    # дешёвый прокси значимости, почти бесплатный:
    salience_cheap = abs(mood_now - mood_prev)
    consolidate(G, event, weight=salience)   # гейт значимости, НЕ одобрения
    if ap.surprise high: mark(G, event, "need_to_discover")

def update_reaction(G, situation, user, reaction, received_ok):
    e = reaction_edge(G, situation, user, reaction)   # условие по (ситуация, отношение)
    if received_ok: e.beta.alpha += 1
    else:           e.beta.beta  += 1                 # угасание, не второй коммит
    e.confidence = e.beta.alpha / (e.beta.alpha + e.beta.beta)
    if e.confidence > CONSOLIDATE_TH: promote_to_slow_layer(G, e)
\end{lstlisting}

\section{Скорер характера и отбор голоса}

\begin{lstlisting}[style=pseudo]
def phi(text, axes):                          # LLM-судья по рубрикам осей
    return [judge(text, rubric_a) for rubric_a in axes]   # вектор в R^d

def select_voice_corpus(candidates, target_t, axes):
    scored = [(c, phi(c, axes)) for c in candidates]
    # отбор/перевзвешивание, чтобы среднее выборки приблизилось к target_t
    return greedy_match_mean(scored, target_t)   # не усреднять муз, а фильтровать
\end{lstlisting}

\section{Стенд эксперимента}

\begin{lstlisting}[style=pseudo]
def run_experiment(build_graph, plant, operator, metric):
    G = build_graph()          # ярус 1: SBM/BA/WS + подсаженные сценарии
    truth = plant(G)           # ground truth: что должно произойти
    operator(G)                # walker / propagate / scrub / decay ...
    return metric(G, truth)    # precision/recall починки, сходимость, покрытие
\end{lstlisting}

\chapter{Будущее: стоит ли переходить на гиперграфы}
\label{ch:hypergraphs}

Это приложение --- рассуждение о структурном направлении, а не готовая
рекомендация к внедрению. Вопрос: логично ли в будущем сделать память Рин
\emph{гиперграфом}, а не обычным графом.

\section{Проблема, которая это мотивирует}

Обычный граф связывает узлы \emph{попарно}. Но вывод часто \emph{n-арен}: он
рождается только при одновременном присутствии нескольких посылок, и ни одно
подмножество не достаточно. Канонический пример (глава~\ref{ch:fusion}):
\[
A=\text{``есть коты''},\quad B=\text{``живёт в Португалии''},\quad
C=\text{``ездит на пары''},
\]
\[
\begin{aligned}
X=f(A,B,C)
&=\text{``купи котам корм на выходных}\\
&\qquad \text{(в будни по пути с пар не успеть)''}.
\end{aligned}
\]
Каждая посылка ещё и подтягивает свой контекст (коты$\to$нужен корм;
Португалия$\to$магазины рано закрываются; универ$\to$поздний возврат;
см.\ главу~\ref{ch:fusion}), и только их совмещение даёт $X$. Ни $A\wedge B$,
ни $B\wedge C$, ни $A\wedge C$ не дают $X$. Попарные рёбра между посылками
могут вообще отсутствовать: ``Португалия'' и ``коты'' напрямую не связаны.

\section{Почему симплексы --- плохой ответ}

Первое, что приходит в голову, --- симплициальные комплексы: тройка
$\{A,B,C\}$ как 2-симплекс (треугольник). Но у симплициального комплекса
есть жёсткое требование \emph{downward closure}: если есть 2-симплекс
$\{A,B,C\}$, то обязаны быть и все его грани --- рёбра $A-B$, $A-C$, $B-C$.
А в нашем случае этих попарных связей \emph{нет и не должно быть}:
``Португалия'' и ``коты'' не связаны. Симплициальная память вынудила бы нас
проводить \emph{ложные} попарные рёбра ради формальной корректности
структуры. Это плохой обмен.

\section{Гиперребро --- прямой ответ}

\emph{Гиперребро} связывает \emph{множество} узлов сразу:
$e=\{A,B,C\}$, и оно \emph{не требует} существования попарных рёбер между
ними. Это ровно та семантика, что нужна n-арному выводу: ``эти три вместе
порождают это'', без обязательства, что они связаны попарно.

\begin{intuition}[Гиперграф против графа]
В графе ``дружба'' --- всегда пара. В гиперграфе можно сказать ``эти пятеро
были в одной поездке'' одним ребром, не утверждая, что каждый дружит с
каждым. Для памяти это разница между ``связь между двумя воспоминаниями'' и
``условие из нескольких воспоминаний, порождающее вывод''.
\end{intuition}

\section{Ключевое наблюдение: мы уже наполовину гиперграф}

Важно не пропустить: \textbf{текущая схема Рин уже кодирует гиперрёбра} ---
просто не называя их так. Когда fusion порождает узел $X$ с провенансом
$\{A,B,C\}$ (глава~\ref{ch:fusion}), это factor-graph-кодировка гиперребра
$\{A,B,C\}\to X$: сам вывод $X$ становится узлом, а его рёбра-обоснования
указывают на посылки. Попарных рёбер между $A,B,C$ при этом не требуется.
Формально это просто \textbf{двудольный граф} (узлы-посылки $\leftrightarrow$
узлы-факторы), и его \emph{полностью достаточно} для n-арного вывода ---
никакой высшей размерности не нужно.

Значит ``переход на гиперграфы'' --- это не ``переписать всё'', а более узкий
вопрос: \emph{делать ли гиперрёбра первоклассными объектами} (со своим
обходом и спектральной теорией), или достаточно узла-фактора.

\section{Первоклассные гиперрёбра: за и против}

\paragraph{За.} Появляются \emph{гиперграф-нативные} операции:
обход/spreading activation, уважающий \emph{совместность} срабатывания
(гиперребро ``зажигается'', когда активны \emph{все} его узлы, а не любой
один); гиперграфовые лапласианы и спектральная диагностика (аналоги
главы~\ref{ch:topology} для n-арной структуры). Это может точнее
моделировать ``условие из нескольких посылок''.

\paragraph{Против.} Алгоритмика гиперграфов менее зрелая и дороже графовой.
На практике гиперграф часто \emph{сводят} к графу --- clique-expansion
(заменить гиперребро кликой, что как раз возвращает ложные попарные рёбра!)
или star-expansion (ввести узел-фактор --- то есть ровно наша
factor-кодировка). Есть спектральная теория гиперграфов (нормализованный
гиперграфовый лапласиан; тензорные методы), но она сложнее и вычислительно
тяжелее, а динамический мутирующий гиперграф --- тем более.

\section{Рекомендация}

\begin{principle}[Гиперграф --- да, но постепенно]
Ближний горизонт: оставить \emph{factor-узлы} (вывод $X$ с
провенансом-множеством). Это уже даёт n-арную семантику без гиперграфовой
машинерии и без ложных попарных рёбер. Средний горизонт: сделать
гиперрёбра первоклассными \emph{только там}, где нужны гиперграф-нативные
операции --- совместное срабатывание при обходе и n-арная спектральная
диагностика. Симплициальные комплексы обходить стороной: их downward closure
противоречит разреженности реальных выводов.
\end{principle}

Заметим, что star-expansion (факторные узлы) и первоклассные гиперрёбра ---
это два взгляда на одно и то же; переход между ними инкрементален, а не
революционен. Полное же объединение узлов, рёбер и выводов в объекты одного
рода разной размерности --- это уже клеточные комплексы
(приложение~\ref{ch:cellcomplex}).

\chapter{Будущее: клеточные комплексы как язык для составных актов}
\label{ch:cellcomplex}

\section{Зачем вообще думать о клетках}

Гиперграфовая ветвь уже показывает, что попарных рёбер иногда недостаточно. Но и
гиперребро — ещё не вся картина. У памяти есть не только набор совместно действующих
посылок, но и \emph{структура их совместности}: границы, циклы, составные акты,
которые хочется видеть не как мешок из входов, а как объект более высокой размерности.

\section{Базовая идея}

Клеточный комплекс позволяет держать объекты разных размерностей:
\begin{enumerate}
\item 0-клетки как узлы памяти;
\item 1-клетки как обычные рёбра;
\item 2-клетки как составные акты, выводы или сцены;
\item при желании и более высокие клетки для сложных конфигураций.
\end{enumerate}

Это полезно тогда, когда память должна различать не только ``что с чем связано'',
но и ``что из чего составлено как единый акт''.

\section{Что это даёт памяти}

\paragraph{Во-первых,} можно хранить составную ситуацию как один объект, не
раскладывая её полностью на попарные связи.

\paragraph{Во-вторых,} появляются естественные границы: если 2-клетка кодирует
составной вывод или эпизодический акт, её boundary прямо указывает на набор посылок
или компонент.

\paragraph{В-третьих,} становится богаче язык локальной ревизии: можно понижать не
только узел или ребро, но и целую составную конфигурацию.

\section{Чего это не решает}

\begin{pitfall}[Новая структура не заменяет дисциплину provenance]
Сам по себе переход к клеткам не делает память ни истиннее, ни объяснимее. Если у
системы нет хорошего provenance, gates, deletion-path и ревизии, клеточный комплекс
лишь переносит старые проблемы в более красивую геометрию.
\end{pitfall}

\section{Правильный статус этой ветви}

Для текущего проекта клеточные комплексы — это язык \emph{будущего структурного
обобщения}, а не ближайшая реализация. Их стоит читать так:
\begin{enumerate}
\item как концептуальный мост между графом и n-арными актами;
\item как способ точнее думать о границах составных выводов;
\item как материал для дальнего сравнения с гиперграфами и динамической памятью.
\end{enumerate}

Именно в таком статусе эта ветвь совместима с нынешним roadmap: не как срочная
замена субстрата, а как более богатая геометрия для будущего поколения памяти.

\chapter{Будущее: память как динамика (резервуар, ассоциативные аттракторы, кванты)}
\label{ch:dynamical}

\begin{pitfall}[Вердикт вперёд: это смена субстрата, а не смена топологии]
Приложения~\ref{ch:hypergraphs} и~\ref{ch:cellcomplex} спрашивали «какой \emph{формы} должен
быть граф». Это приложение спрашивает радикальнее: \emph{а нужен ли символический граф
вообще}, или память стоит держать как \emph{динамику} — состояние большой рекуррентной
системы, где факт не хранится объектом, а «протоптан» в связях. Идея красивая, глубоко
укоренённая в нейронауке, и \textbf{мы её сейчас не строим}. Но не потому, что она наивна —
наоборот, разберём её всерьёз. Причина в том, что для \emph{компаньона} с требованиями
редактируемости, забывания-по-команде, провенанса и поддержки истинности распределённая
память в весах ломает ровно те свойства, ради которых мы и завели типизированный граф. Ниже
— честный разбор: что здесь ценно (и что мы, оказывается, уже делаем), что дорого, и как
оставить дверь открытой, не переписывая мозг.
\end{pitfall}

\section{Постановка: память в весах, воспоминание как траектория}

Наш субстрат (гл.~\ref{ch:substrate}) — типизированные узлы и явные рёбра: у «узла Helena»
есть содержание, у ребра — тип и вес. Альтернатива, которую предлагает эта ветка, меняет
само определение:
\[
\begin{aligned}
\text{память} &= (G, W, \theta),\\
\text{воспоминание} &= \text{траектория в пространстве состояний},\\
\text{обучение} &= \Delta W .
\end{aligned}
\]
Здесь $G$ — топология (какие пути вообще есть), $W$ — веса связей (насколько они
протоптаны), а $\theta$ — параметры узлов-вычислителей. Узел \texttt{\#48372} не значит
«Лиссабон»; смысл распределён по тысячам узлов и существует только как \emph{паттерн
активации} $+$ \emph{связность} $+$ \emph{траектория}. Запрос не «ищет chunk», а
\emph{возмущает} систему: бросаем заряд, сеть эволюционирует, из неё выныривает активный
подграф, и уже его вербализует LLM.

\begin{intuition}[Бассейн, камни, датчики]
Классический образ резервуарных вычислений: бассейн воды, в который бросают камни (входной
поток). Каждый камень поднимает волны; следующий приходит, пока прежние не улеглись, — и
поверхность в момент $t$ несёт отпечаток нескольких последних бросков. Датчики по краям
считывают состояние, и простой линейный слой учится по ним отвечать. Сам бассейн обучать не
надо: его физика уже превращает историю входов в богатое состояние. Reservoir computing —
это \emph{случайное нелинейное динамическое отображение признаков для последовательностей},
аналог kernel trick, но во времени.
\end{intuition}

Важная оговорка сразу: \emph{классический} резервуар (Echo State Network, Liquid State
Machine) держит $W_{\text{res}}$ \emph{замороженным} и обучает только считывающий слой. Он
обладает \emph{затухающей} памятью (fading memory) на конечный горизонт — прекрасно для
временных рядов, но не для «узнала в 2026-м, вспомнила в 2031-м». Для памяти длиной в годы
нужен \emph{пластичный} резервуар ($W$ медленно меняется опытом) плюс механизм устойчивого
хранения из мира \emph{ассоциативных аттракторов} (Hopfield/CA3): воспоминание — это
\emph{бассейн притяжения}, частичный cue скатывается в него (pattern completion). То есть это
не «reservoir computing» в чистом виде, а гибрид на пересечении резервуара, ассоциативной
памяти, аттракторных сетей и пластичной динамики.

\section{Математическая ось: один лапласиан — три динамики}

Самое элегантное в этой ветке — то, что \emph{одна и та же} топология допускает три
принципиально разные динамики. Для графа с лапласианом $L=D-A$:
\[
\begin{aligned}
\text{диффузия:}\quad & \dot x = -Lx
&& \Rightarrow\quad x(t)=e^{-tL}x(0),\\
\text{волна:}\quad & \ddot x = -Lx
&& \Rightarrow\quad x(t)=\cos(t\sqrt{L})\,x(0),\\
\shortstack[l]{квантовое\\блуждание:}\quad & i\dot\psi = L\psi
&& \Rightarrow\quad \psi(t)=e^{-itL}\psi(0).
\end{aligned}
\]
Один $L$ порождает $e^{-tL}$ (рассеивающую память), $\cos(t\sqrt L)$ (волновую, с
интерференцией и без затухания) и $e^{-itL}$ (унитарную квантовую). Это не словесная
аналогия, а линейная алгебра: «бросить заряд в узел и смотреть, куда потечёт» имеет буквальный
волновой и квантовый аналоги на той же сети.

\begin{intuition}[Мы уже живём в углу $e^{-tL}$]
Полезно осознать: наш ретрив (гл.~\ref{ch:walker}) — это \emph{и есть} диффузионный угол
этой оси. `diffuse()` считает усечённое $e^{-tL}$-подобное распространение из seed-якорей;
«вспышки, слабеющие к окраинам» — это спад $e^{-tL}$. То есть значительная часть красоты
этого приложения у нас уже реализована, просто в скромной, редактируемой форме. Волновой и
квантовый углы — это \emph{другие операторы над той же топологией}, и их честно можно
сравнить на одном графе (см.\ открытые вопросы в конце).
\end{intuition}

\section{Нейронаучная опора (принцип, а не крысиный мозг)}

Ветка опирается на реальные результаты, и стоит назвать их аккуратно. Обучение оставляет
устойчивые следы в самой сети (engram-клетки: клетки, участвовавшие в формировании
воспоминания, можно искусственно возбудить и вызвать его; хранение связывают с синаптическим
усилением engram-ансамблей). Гиппокамп работает не как «жёсткий диск», а как система
связывания/индексации и реактивации распределённых корковых представлений (hippocampal
indexing theory), причём со временем память дрейфует от гиппокампа к коре (systems
consolidation, replay). Complementary Learning Systems (CLS) прямо противопоставляет
\emph{быстро} обучающийся гиппокамп \emph{медленно} обучающейся коре — и это вычислительно
мотивированная защита от катастрофического забывания. Аттракторная память (Hopfield; CA3 как
autoassociative network) хранит паттерн как бассейн притяжения. Мы заимствуем \emph{принцип}
(быстрый индекс $+$ медленное распределённое хранилище $+$ восстановление по частичному cue),
не пытаясь эмулировать биологию до деталей.

\section{Предлагаемая гибридная архитектура}

Синтез, к которому сходится обсуждение, — не «резервуар вместо графа», а шестислойный объект
плюс архив сбоку:
\begin{enumerate}\itemsep2pt
\item \textbf{Event encoder.} Приходит не «факт», а \emph{событие} целиком — контент, говорящий,
тон, контекст, время, настроение — и кодируется в вектор $x_t$. Источник знания становится
\emph{частью паттерна}, а не строкой БД (значит, память сможет и \emph{ошибиться} в источнике
— по-человечески).
\item \textbf{Hippocampal pattern separator / index.} Из $x_t$ строится разреженный код
$h_t=\mathrm{kWTA}(Px_t)$ (активны немногие units из многих); похожие события разводятся в
непохожие коды (pattern separation). Это «умный стартер»: куда бросить заряд.
\item \textbf{Fast associative memory.} Быстрая Hebbian-запись «эти вещи произошли вместе»
($H\!\leftarrow\!H+\eta_h h_t h_t^\top$), связывающая индекс с распределённым состоянием.
\item \textbf{Slow plastic cortical reservoir.} Большая разреженная рекуррентная сеть
\[
r(t{+}1)=(1-\alpha)\,r(t)+\alpha\,\sigma\!\big(Wr(t)+Ux(t)+Bh(t)\big),
\]
где $W$ \emph{медленно} пластичен ($\eta_c\ll\eta_h$). Быстрый гиппокамп $+$ медленная кора —
это и есть CLS.
\item \textbf{Continuous activation field.} Живое поле $r(t)$ — «о чём Рин сейчас думает».
\item \textbf{Latent-to-LLM interface.} Мост от $r(t)$ к языковой модели.
\end{enumerate}
Плюс \textbf{replay/consolidation} в простое (ночью гиппокамп многократно проигрывает след,
медленно перестраивая $W$ — «переваривание дня» не через \texttt{summarize()}, а через
проигрывание), и \textbf{граф-архив} как «чёрный ящик самолёта» сбоку.

\begin{principle}[Граф не выбрасывается — он становится архивом и учителем]
Ключевой примирительный ход: наш явный граф перестаёт быть \emph{моделью} человеческой памяти
и становится (а) \emph{эпизодическим архивом} (что реально произошло, когда, откуда знаем) и
(б) \emph{учителем} для пластичной памяти (материал для replay-консолидации). Тогда
специализация чистая: \emph{граф} — «что действительно было записано», \emph{резервуар} — «что
Рин из этого сейчас помнит» (неточно, ассоциативно, с интерференцией). Это не дублирование, а
разделение ролей.
\end{principle}

\section{Настроение как нейромодулятор и «пирожки Алисы»}

Отдельная и, на мой взгляд, самая \emph{сразу применимая} идея ветки: mood перестаёт быть
вектором, который кладут в промпт, и становится \emph{нейромодулятором} — он меняет
\emph{геометрию} recall:
\[
W_{\text{effective}} = \gamma(m)\,W,\qquad
A_{\text{effective}} = A\odot M(\mathrm{mood}).
\]
Тот же cue при разном настроении уходит по разным траекториям («помнишь Тбилиси?» $\to$ переезд
$\to$ Ада; или $\to$ ОИЯИ $\to$ смена жизни). Инерция настроения тут — не недостаток, а
\emph{нужное} свойство: медленный low-pass-контроллер задаёт устойчивый ландшафт для быстрых
процессов и не «дёргает руль» на каждый токен.

\begin{intuition}[Гомеостаз как гормоны: заманчиво, но чинит болезнь, которую само же создаёт]
Предложение «мониторить энергию поля $E(t)=\frac1N\sum r_i^2$ и впрыскивать седатив (поднять
утечку $\lambda$, снизить крутизну $\beta$) при перевозбуждении или кофеин при апатии» — это
настоящая идея из neuromodulated networks, и она красива. Но будем честны: это \emph{лекарство
от нестабильности, которую вносит сам пластичный резервуар}. На нашем текущем субстрате
диффузия ограничена по построению (decay $+$ floor $+$ max\_steps $+$ max-CRDT активация,
гл.~\ref{ch:walker}) — «эпилептического шторма» просто негде взяться, и гормональный
контроллер не нужен. Более того, он добавляет ровно то, чего мы весь проект избегали:
свободные параметры ($E_{\min},E_{\max}$, коэффициенты впрыска), которые нечем калибровать без
downstream-сигнала. Вывод сбалансированный: \emph{mood-как-нейромодулятор} — берём хоть сейчас
(он и так входит в правило перехода через $\mathrm{mc}(B\mid\mood)$); \emph{гормональный
гомеостаз} — осмыслен только если уже ушли в распределённый пластичный резервуар, то есть это
инструмент \emph{той} ветки, а не нашей.
\end{intuition}

\section{Честный разбор: что покупаем и чем платим}

\paragraph{Что действительно ценно.} Ассоциативная \emph{реконструкция} вместо поиска chunk'а;
человекоподобная память (нечёткое время, ассоциативные ошибки источника, интерференция как
feature); mood как геометрия recall; CLS как защита от катастрофического забывания; и красивая
ось $e^{-tL}/\cos(t\sqrt L)/e^{-itL}$, дающая единый язык для diffusion/wave/quantum памяти.

\paragraph{Чем платим — и это надо назвать прямо.} Три технических риска и один
принципиальный.
\begin{itemize}\itemsep3pt
\item \textbf{Readout-бутылочное горлышко.} Перевод $r(t)\in\mathbb{R}^{10^5}$ обратно в
токены — самая хрупкая часть. Тупой декодер (MLP) может не извлечь смысл; умный — \emph{сам}
сделает всю работу и замаскирует бесполезность динамики (нужен экспериментальный контроль:
сначала нарочно \emph{тупой} $D$). А чтобы декодер обучить, нужны пары «состояние $\to$
правильный ответ» — которых для автобиографии \emph{нет}. Откуда берётся обучающий сигнал —
главный незакрытый вопрос. Элегантный вариант (резервуар модулирует hidden states LLM через
cross-attention/prefix) концептуально безупречен, но требует глубокого fine-tuning локальной
модели и почти исключает API-бэкенд; дешёвый вариант (декодировать в текст) — это, по сути,
RAG с лишними шагами, теряющий главную приманку «память остаётся неязыковой».
\item \textbf{Пластичность против гомеостаза (край хаоса).} Держать $\rho(W)\lesssim1$ на
\emph{непрерывно меняющемся} $W$ — нерешённая в вычислительной нейронауке задача: пластичность
воюет с устойчивостью. Плюс аттракторная память имеет малую ёмкость и \emph{ложные}
аттракторы; «плотные» версии Хопфилда чинят ёмкость, но математически превращаются в
attention — и тогда закономерен вопрос «а почему не взять память самого трансформера».
\item \textbf{Цена inference.} Прогонять $N$ шагов динамики по графу в сотни тысяч узлов на
\emph{каждый} запрос, да ещё фоном для replay — тяжело; нужна серьёзная оптимизация разреженных
умножений. Это не смертельно, но и не бесплатно (наш же принцип «граф дёшев только при чистой
обвязке»).
\item \textbf{Принципиальное: теряются свойства, ради которых мы строили граф.}
Редактируемость, забывание-по-команде (крипто-шред), провенанс, поточечная ревизия (TMS),
метки contested, семантика node\_kind — всё это распределённый резервуар \emph{уничтожает}.
Нельзя крипто-шреднуть факт из $W$; нельзя показать провенанс; нельзя гарантированно забыть
конкретное («Рин, забудь последние 10 минут» — у компаньона это \emph{UX/приватность}, а не
когнитивная поэзия). Возражение «человек тоже не забывает по команде» верно \emph{когнитивно},
но игнорирует, что у Рин есть \emph{юридическо-продуктовые} требования, которых у человека нет.
Архив-«чёрный ящик» спасает сырое удаление, но \emph{не} отменяет того, что резервуар уже
\emph{выучил} факт в $W$ — распределённо и не-разучиваемо.
\end{itemize}

\begin{principle}[Ключевая развилка: «не хранить строки» не значит «не иметь адресуемых узлов»]
Здесь легко проскочить подмену. Есть два разных смысла «памяти без фактов»:
\emph{(а)} не хранить факт как \emph{человекочитаемую строку} — узлы могут быть эмбеддингами;
это \emph{дёшево и совместимо} с нашим субстратом; \emph{(б)} вообще не иметь \emph{адресуемых
объектов памяти} — всё размазано в $W$; это и покупает человекоподобность, и \emph{ломает}
редактируемость/провенанс/забывание. Ценное — это (а) $+$ ассоциативная динамика; дорогая
переоценка — (б). А наш ретрив уже делает (а)$+$динамику: «возмути узел и смотри, что
загорится» — это диффузия по типизированному графу. Значит львиную долю пользы этой ветки мы
\emph{уже} имеем, не платя за (б). По-настоящему \emph{новых} адоптируемых кусков три:
пластичные рёбра (Hebbian-усиление от со-активации — по сути формализованная консолидация),
mood-как-нейромодулятор, и CLS-разделение быстрый/медленный (у нас грубо есть как слои
session$\to$life).
\end{principle}

\section{Ближняя инженерная форма ветки: что именно стоило бы пробовать}

Если отбросить соблазн ``сразу заменить память распределённой динамикой
целиком'', то у этой ветки остаётся очень сильная и вполне рабочая ближняя
форма. Её разумно формулировать так:

\begin{center}
\texttt{explicit graph archive + hippocampal associative index + small plastic reservoir + replay}
\end{center}

То есть не ``убираем граф'', а вводим рядом с ним вторую форму памяти: граф
отвечает за адресность, provenance, forgetting и аудит, а динамический слой
--- за pattern completion, ассоциативный recall и временную интерференцию.

\paragraph{Как это должно выглядеть конструктивно.}
\begin{enumerate}\itemsep2pt
\item \textbf{Архив остаётся явным.} Все наблюдения, коммиты, ретракции и
удаления по-прежнему живут в event-log и типизированном графе.
\item \textbf{Быстрый индекс остаётся разреженным.} Cue сначала поднимает не
весь мир, а небольшой набор ассоциативных seed-узлов --- ``гиппокампальный''
стартер для дальнейшей динамики.
\item \textbf{Динамика живёт не на всём графе, а на активном подграфе.}
Практически это самый важный инженерный ход: первый прототип не обязан
гонять пластику по сотням тысяч узлов. Куда разумнее взять подграф, уже
поднятый retrieval'ом (якоря + локальное расширение + bridge-узлы), и
крутить reservoir-подобную динамику только на нём. Тогда и цена inference
падает, и диагностика становится обозримой.
\item \textbf{Replay идёт от архива к динамике, а не наоборот.} Ночью или в
idle режимах архив проигрывает важные эпизоды, а пластичный слой медленно
меняет веса и режимы ассоциации.
\item \textbf{Write-path остаётся у графа.} Динамический слой не пишет факты
напрямую; он поднимает кандидатов, контркандидатов, bridge-конфигурации, а
коммиттер судит их обычным путём.
\end{enumerate}

\begin{pitfall}[Самая опасная ошибка прототипа]
Если первый же прототип будет строиться как ``огромный свободно
пульсирующий резервуар поверх всей памяти'', проект почти наверняка утонет
не в идеях, а в инфраструктурной боли: дорого, трудно стабилизировать, трудно
понимать, трудно проверять. Активный подграф --- не уступка, а правильная
ступень масштабирования.
\end{pitfall}

\paragraph{Где здесь место mood и гомеостазу.} Настроение в этой ветке
полезно не как ещё один факт, а как режимный контроллер: медленный
параметр, меняющий допустимую глубину, усиление и ассоциативную
``температуру'' динамики. Это почти прямой low-pass-контроллер для границы
между апатией и перевозбуждением. Идея ``гормонального'' гомеостаза ---
мониторить энергию поля
\[
E(t)=\frac{1}{N}\sum r_i(t)^2
\]
и скорость его изменения
\[
V(t)=\frac{1}{N}\sum \big(r_i(t)-r_i(t-1)\big)^2
\]
--- имеет смысл только если динамический слой уже реально способен входить в
плохие режимы. Тогда можно вводить быстрые broadcast-модуляторы: ``седатив''
поднимает утечку и снижает усиление, ``кофеин'' делает систему
чувствительнее, ``дофамин'' временно повышает пластичность. Но на
сегодняшнем графовом ядре это лекарство от болезни, которой ещё нет.

\paragraph{Что проверять на стенде.} У этой ветки есть вполне конкретные, не
философские тесты.
\begin{enumerate}\itemsep2pt
\item Pattern completion: восстанавливает ли слой правильную окрестность по
частичному cue лучше, чем чистый graph-retrieval.
\item Интерференция: не начинает ли новый опыт разрушать старые устойчивые
ответы.
\item Source confusion: не растёт ли доля случаев, где ассоциативный recall
красив, но provenance расползается.
\item Readout control: сохраняется ли смысл в самой динамике, если посадить
нарочно тупой декодер.
\item Стабильность: не уходит ли слой в тривиальное затухание или в
хаотический шторм.
\end{enumerate}

Именно в таком виде ветка выглядит не фантазией, а сильной
исследовательской программой: сначала маленький пластичный слой на активном
подграфе, потом replay, потом режимные контроллеры, и только после этого ---
разговор о более радикальных носителях.

\section{Квантовый слой: честно}

Соблазн прямой: ассоциативная память $=$ волновая динамика, кванты $=$ волновая механика,
значит кванты идеальны. Не совсем. \emph{Чистая} квантовая динамика унитарна
($|\psi(t)\rangle=U(t)|\psi(0)\rangle$): обратима, сохраняет норму, \emph{не забывает} и
\emph{не} схлопывается в аттрактор ($\|Ux-Uy\|=\|x-y\|$ — нет сжатия). А ассоциативной памяти
нужно ровно сжатие: много близких состояний $\to$ одно воспоминание. Поэтому квантовая
ассоциативная память неизбежно тащит \emph{диссипацию}, измерения, связь со средой,
open-system-динамику. Второй подвох: $2^n$ амплитуд $\ne$ $2^n$ \emph{доступных} чисел —
измерение выдаёт ограниченную классическую информацию, а классический вход ещё надо
\emph{загрузить} в квантовое состояние. Поэтому экспоненциальное гильбертово пространство не
даёт автоматически экспоненциальной полезной памяти.

\begin{openq}[Оценка кванта: исследовательски высоко, практически сейчас низко]
Область quantum reservoir computing и quantum associative memory реально существует и активно
развивается (спиновые сети, фотоника на сотнях мод, quantum Hopfield на малых NISQ-устройствах).
\emph{Дисклеймер о верифицируемости}: конкретные свежие эксперименты и цифры «превосходства»,
всплывшие в обсуждении, лежат за горизонтом надёжной проверки этого текста — их стоит держать
как \emph{иллюстрацию направления}, а не как установленный quantum advantage; сам advantage для
памяти пока не доказан. Разумная ближняя форма — \emph{не} «память Рин $=$ квантовый компьютер»,
а гибрид: классический гиппокамп $\to$ \emph{маленький} квантовый резервуар как чрезвычайно
богатый нелинейный feature-generator $\to$ классическое пластичное хранилище. Квант тут — не
хранилище на двадцать лет жизни, а мощный возбудитель одного оператора.
\end{openq}

\section{Вердикт и как оставить дверь открытой}

Для Рин с её требованиями (редактировать, забывать, объяснять источник, поддерживать
истинность) \textbf{символический типизированный граф остаётся хребтом}, а динамический слой —
это \emph{дополняющая} ассоциативная надстройка («что я сейчас помню») и отношения
\emph{учитель\,--\,ученик} (граф $=$ архив/учитель). Радикальную версию «убрать семантический
граф» для наших требований отклоняем — с конкретными причинами (потеря забывания/провенанса/
TMS), а не из робости. Но инженерно правильно \textbf{держать динамический слой за чистым
интерфейсом}: тогда классический решатель $R_{\text{classical}}$ можно однажды заменить на
$R_{\text{oscillator}}/R_{\text{photonic}}/R_{\text{quantum}}$, не переписывая остальное.

\begin{openq}[Фальсифицируемые эксперименты, если идти в эту ветку]
Три чистых вопроса, которые можно поставить на стенде, не ломая текущую систему:
\emph{(1)} сравнить $e^{-tL}$, $\cos(t\sqrt L)$ и $e^{-itL}$ на \emph{одной} topologии
memory-графа — какой профиль лучше поднимает правильную окрестность (diffusion vs wave vs
quantum-walk memory); \emph{(2)} даёт ли \emph{пластичная} надстройка над типизированным графом
(Hebbian-усиление со-активированных рёбер) измеримо лучший recall, чем статические веса, — и не
ценой ли катастрофической интерференции; \emph{(3)} readout-контроль: посадить нарочно тупой
декодер и проверить, есть ли вообще содержательная информация в динамике, прежде чем строить
сложный мост к LLM. Все три — на golden-микромире (гл.~\ref{ch:experiment}), где есть ground
truth.
\end{openq}

\chapter{Будущие ветви: как расширять систему, не теряя её}

\section{Две обещанные дальние ветви}

У расширенного курса есть ещё одна задача: не просто завершить сегодняшнюю версию,
но и подготовить язык, на котором будут обсуждаться следующие большие ветви. На
горизонте уже видны по меньшей мере две:
\begin{enumerate}
\item резервуарные вычисления как дополнительная ассоциативно-динамическая среда;
\item ассоциативная память на кубитах как исследовательская ветвь квантового слоя.
\end{enumerate}

\section{Что нужно требовать от любой новой ветви}

Новая ветвь не должна приходить в систему как эстетический жест. Для неё нужен
одинаковый протокол допуска.
\begin{enumerate}
\item Какая конкретная слабость базовой архитектуры атакуется?
\item Как будет измеряться улучшение?
\item Как ветвь взаимодействует с provenance, deletion и revision?
\item Что остаётся в явном графе, а что уходит в скрытую динамику?
\end{enumerate}

\section{Резервуарная ветвь: что в ней действительно интересно}

Резервуарная схема интересна не тем, что ``она нейроподобна'', а тем, что даёт
богатый нелинейный feature space для частичных cue и сложных временных контекстов.
Если использовать её осторожно, она может:
\begin{enumerate}
\item улучшать pattern completion;
\item давать признаки для сложных mood/history-dependent retrieval policies;
\item служить дополнительным предложителем bridge-кандидатов.
\end{enumerate}

Но её нельзя пускать напрямую в write-path без жёсткого интерфейса.

\section{Квантовая ветвь: чего нельзя обещать заранее}

Квантовая ассоциативная память соблазнительна по двум причинам: волновая интуиция
и естественная работа с суперпозициями. Но сама по себе эта красота ничего не
доказывает. До стенда допустимы только такие утверждения:
\begin{enumerate}
\item ветвь может дать иной режим диффузии и поиска;
\item ветвь может дать иной bias к pattern completion;
\item ветвь \emph{не считается полезной}, пока не показала измеримый выигрыш на
общем стенде против классических альтернатив.
\end{enumerate}

\section{Правильная амбиция}

\begin{principle}[Большая ветвь должна быть сравнима с малой системой]
Чем радикальнее новая ветвь, тем строже она обязана сравниваться с уже работающей
графовой памятью. Иначе вместо расширения курса получается побег из проверяемости.
\end{principle}

Это правило и будет держать будущие главы честными.

\chapter{Расширение: гибридный резервуарный слой для памяти Рин}
\label{ch:reservoirbranch}

\section{Почему эта ветка отличается от прочих будущих ветвей}

Не всякая дальняя ветвь одинаково близка к инженерной практике. Гиперграфы, клеточные
комплексы и квантовые носители полезны как сильные структурные и исследовательские идеи,
но резервуарно-динамическая ветвь выделяется на их фоне по трём причинам.

\begin{enumerate}
\item Она \textbf{многое объясняет} уже на языке текущей архитектуры: почему retrieval
можно понимать как возмущение системы, почему mood естественно трактовать как режим
динамики, почему replay похож на настоящую консолидацию, а не на summarization.
\item Она \textbf{интерпретируема}: здесь нет требования отдать весь смысл чёрному ящику.
Напротив, явный граф, быстрый индекс, динамический слой и write-path можно развести по ролям.
\item Она \textbf{doable}: первый прототип можно собрать, не переписывая весь мозг Рин,
а добавив рядом с текущим графовым позвоночником новый слой с чистым интерфейсом.
\end{enumerate}

\section{Ключевая идея: две формы памяти вместо одной}

Главный выигрыш этой ветки начинается не с замены графа, а с отказа от ложной альтернативы
``или архив, или динамика''. Плодотворная формула другая:
\[
\boxed{
\text{память Рин}=
\text{адресуемый архив}
\;+\;
\text{динамическая ассоциативная среда}
}
\]

Архив отвечает на вопросы:
\begin{itemize}
\item что действительно было записано;
\item когда это произошло;
\item откуда это известно;
\item что было удалено;
\item какие потомки надо ревизовать.
\end{itemize}

Динамическая среда отвечает на другие вопросы:
\begin{itemize}
\item что всплывает по частичному cue;
\item какие далёкие мосты сейчас собираются в одну конфигурацию;
\item насколько контекст ``живой'', а не просто найденный;
\item в какой ассоциативный бассейн скатывается текущее состояние.
\end{itemize}

\begin{principle}[Архив и динамика решают разные классы задач]
Если архив пытается быть живой памятью, он становится тяжёлой базой данных с плохой
ассоциативностью. Если динамика пытается быть архивом, она теряет provenance, deletion,
ревизию и юридическую адресность. Сильная архитектура не выбирает одно из двух, а
разводит их роли.
\end{principle}

\section{Краткая теория: что такое резервуар в этом контексте}

В классическом reservoir computing есть большая рекуррентная система с фиксированными
случайными связями:
\[
h_t = \sigma(W_{\mathrm{res}}h_{t-1}+W_{\mathrm{in}}x_t),
\]
и обучается только считывающий слой:
\[
y_t=W_{\mathrm{out}}h_t.
\]
Это даёт богатое нелинейное отображение истории входов в пространство состояний.

Но для памяти Рин классический резервуар в чистом виде недостаточен по двум причинам.
\begin{enumerate}
\item Он даёт в основном \emph{fading memory}, то есть естественную память на конечный
горизонт.
\item Он не решает задачу устойчивого хранения многолетнего опыта без изменения самого
хранилища.
\end{enumerate}

Поэтому здесь уместнее не ``reservoir computing как есть'', а гибрид:
\[
\boxed{
\text{быстрый ассоциативный индекс}
\;+\;
\text{медленный пластичный резервуар}
\;+\;
\text{явный архив}
}
\]
То есть:
\begin{itemize}
\item быстрый индекс обеспечивает pattern separation и запуск recall;
\item пластичный слой даёт богатую ассоциативную динамику;
\item архив страхует долгую память, provenance и точную ревизию.
\end{itemize}

\section{Архитектурный скелет}

\subsection*{1. Event encoder}
На вход идёт не ``голый факт'', а событие целиком: текст, время, настроение, источник,
социальный контекст, task-status, salience. Результатом становится вектор события $x_t$.

\subsection*{2. Быстрый разреженный индекс}
Из $x_t$ строится sparse code:
\[
z_t = \mathrm{kWTA}(Px_t),
\]
где активна малая доля компонент. Это ближайший аналог ``гиппокампального'' индекса:
не хранение всей памяти, а адресное зажигание стартовых площадок.

\subsection*{3. Быстрая ассоциативная связка}
На этом слое действует быстрая запись совместности:
\[
H \leftarrow H + \eta_h z_t z_t^\top,
\qquad \eta_h \text{ велика}.
\]
Он нужен для свежих связей, которые должны появляться почти немедленно.

\subsection*{4. Медленный пластичный резервуар}
Дальше живёт собственно динамический слой:
\[
r_{t+1}
=(1-\lambda)\,r_t + \lambda\,\sigma(W_t r_t + Ux_t + Bz_t).
\]
Пластичность медленная:
\[
W_{t+1}=W_t+\eta_c \Delta W_t,
\qquad \eta_c \ll \eta_h.
\]
Именно это разделение временных шкал делает идею похожей на CLS, а не на хаотическую
самопереписывающуюся сеть.

\subsection*{5. Replay}
В idle и sleep режимах архивные события снова проигрываются через индекс и резервуар.
Так пластичный слой учится не от разрозненных единичных вспышек, а от повторяемых,
структурно отобранных паттернов.

\subsection*{6. Readout}
Readout не должен быть ни тупым до бесполезности, ни настолько умным, чтобы съесть
всю работу динамики. Идеологически здесь есть три уровня:
\begin{enumerate}
\item простое считывание диагностик и классов;
\item считывание набора кандидатов / активных подграфов;
\item интерфейс в hidden-state или cross-attention слой LLM.
\end{enumerate}

\subsection*{7. Архив-учитель}
Явный граф остаётся не просто ``запасным вариантом'', а учителем и судьёй: он
поставляет материал для replay, хранит provenance и принимает окончательные записи.

\section{Новая инженерная рекомендация: начинать не со всего графа}

Здесь я специально добавляю предложение, которого не было в исходном обсуждении в такой
форме, но которое следует из него как наиболее разумный путь prototyping.

\begin{principle}[Первый резервуар должен жить на активном подграфе]
Самый практичный MVP — не ``пластичный океан поверх всей памяти'', а динамический слой
на активном подграфе, который уже поднят retrieval'ом: якоря, их окрестность, bridge-узлы,
несколько важных provenance-носителей и контрсвидетельства.
\end{principle}

Пусть retrieval поднял подграф $G_t=(V_t,E_t)$. Тогда резервуарный шаг идёт не на всём
$V$, а только на $V_t$:
\[
r^{(t)}_{k+1}
=(1-\lambda)\,r^{(t)}_k + \lambda\,\sigma(W^{(t)} r^{(t)}_k + u^{(t)}),
\]
где $W^{(t)}$ индуцирован текущим активным подграфом.

Преимущества этого решения:
\begin{enumerate}
\item вычислительная цена резко падает;
\item появляется хороший стенд для диагностики;
\item динамический слой автоматически привязан к интерпретируемому фрагменту памяти;
\item replay можно запускать локально по тематическим областям, а не по всему миру сразу.
\end{enumerate}

Недостаток только один: это уже не ``чистая распределённая память всего мозга'', а
слой локальной ассоциативной динамики. Но именно это и делает ветку инженерно подъёмной.

\section{Mood как контроллер режима}

Диалог очень точно попадает в одно важное место: mood в такой архитектуре полезно
понимать не как ещё один текстовый факт и не как косметический тег в промпте, а как
\emph{медленный регулятор режима динамики}.

Пусть $\gamma(m)$ — функция настроения. Тогда она может модифицировать эффективный
усилитель рекуррентной матрицы:
\[
W_{\mathrm{eff}} = \gamma(m)\,W.
\]
Или, более осторожно, влиять на leak, gain и bias перехода:
\[
r_{t+1}
=(1-\lambda(m))r_t
 + \lambda(m)\,\sigma\!\bigl(\beta(m)Wr_t + Ux_t + b(m)\bigr).
\]

\begin{intuition}[Почему именно mood подходит на эту роль]
У настроения есть свойство, которое в инженерии систем управления необычайно удобно:
оно инертно. Оно не дёргается на каждый токен, а меняется медленно. Значит, оно может
быть low-pass-контроллером, задающим общий ландшафт для быстрых когнитивных колебаний.
\end{intuition}

Интерпретация проста:
\begin{itemize}
\item спокойный режим делает динамику более диссипативной;
\item сфокусированный или тревожный режим подводит её ближе к зоне длинных ассоциаций;
\item творческий поиск может временно расширять допустимую область лёгкого хаоса.
\end{itemize}

\section{Гомеостатический контроллер: «пирожки Алисы»}

Вторая идея из диалога — ещё интереснее: не только медленный режимный сдвиг, но и
быстрый контур искусственного ``химического'' гомеостаза.

\subsection*{Градусник системы}
Полезно следить по крайней мере за двумя агрегатами:
\[
E(t)=\frac{1}{N}\sum_{i=1}^N r_i(t)^2,
\qquad
V(t)=\frac{1}{N}\sum_{i=1}^N (r_i(t)-r_i(t-1))^2.
\]
Первый измеряет ``энергию'' поля, второй — скорость его изменения.

\subsection*{Быстрые модуляторы}
Если система выходит из рабочего диапазона, можно временно впрыскивать глобальные
модуляторы.

\begin{enumerate}
\item \textbf{Седатив}: увеличивает утечку и уменьшает gain нелинейности.
\item \textbf{Кофеин}: уменьшает затухание и повышает возбудимость.
\item \textbf{Дофаминоподобный модулятор}: временно поднимает learning rate пластичности,
когда пришёл важный или неожиданный опыт.
\end{enumerate}

\begin{pitfall}[Это не подарок бесплатно]
Гомеостатический контур очень красив, но он вносит собственные параметры, собственные
границы, свои ошибки калибровки и свой новый способ сломать систему. Поэтому он не
должен появляться раньше, чем реально возникнет болезнь, которую он лечит.
\end{pitfall}

\section{Как объединяются mood и гомеостаз}

Лучший способ совместить эти две идеи — не путать их.

\begin{itemize}
\item \textbf{Mood} — медленный режимный слой. Он сдвигает целевые границы, в которых
система считает себя ``нормальной''.
\item \textbf{Гомеостаз} — быстрый автоматический контур. Он следит, чтобы текущая
динамика не вышла за границы, заданные более медленным режимом.
\end{itemize}

Например, в спокойном режиме $E_{\max}$ ниже и система быстрее ``успокаивается''. В
творческом режиме $E_{\max}$ выше: допускается более длинная и смелая ассоциативная
траектория, прежде чем включится торможение.

\section{Что строить сначала}

У этой ветви очень легко потерять масштаб и начать ``с будущего будущего''. Поэтому
полезен строгий порядок прототипирования.

\begin{enumerate}
\item \textbf{Шаг 1.} Пластичность на типизированном графе без отдельного резервуара:
Hebbian-like усиление со-активированных рёбер и replay от архива.
\item \textbf{Шаг 2.} Малый динамический слой на активном подграфе: проверка, даёт ли
он реальный выигрыш в pattern completion и bridge formation.
\item \textbf{Шаг 3.} Mood как параметр режима динамики.
\item \textbf{Шаг 4.} Только при реальной нестабильности — гомеостатические ``гормоны''.
\item \textbf{Шаг 5.} Более глубокий readout-интерфейс к LLM.
\end{enumerate}

\begin{workedbox}[Как выглядит хороший MVP]
Хороший MVP этой ветки — это не ``Рин теперь вся reservoir''. Хороший MVP — это
следующее: есть golden-микромир, есть текущий graph-retrieval, есть активный подграф,
есть малый пластичный слой поверх него, и можно строго измерить, улучшились ли
partial-cue recall, bridge rate и устойчивость к интерференции.
\end{workedbox}

\section{Как проверять}

Для этой ветки нужны отдельные тесты, отличные от обычного RAG-like retrieval.

\subsection*{1. Partial cue completion}
Подать не полный ключ, а его фрагмент, и измерить, дотягивается ли система до
правильной конфигурации лучше базовой версии.

\subsection*{2. Delay under interference}
Записать событие, потом зашумить систему похожими событиями, затем проверить, что
именно всплывает спустя время и после replay.

\subsection*{3. Provenance stress}
Проверить, не выросла ли способность вспоминать ценой потери объяснимости.

\subsection*{4. Stability envelope}
Смотреть на $E(t)$, $V(t)$, прокси для spectral radius и длину ассоциативных хвостов.

\subsection*{5. Ablation of controllers}
Сравнить: без mood-модуляции, с mood-модуляцией, с грубым homeostasis-контуром,
чтобы отделить полезную режимность от чисто декоративной.

\section{Главные риски}

\begin{enumerate}
\item \textbf{Readout bottleneck.} Динамика может быть богата, но недоступна для
хорошего извлечения.
\item \textbf{Пластичность против устойчивости.} Слишком живая система быстро учится и
быстро разрушается.
\item \textbf{Вычислительная цена.} Даже локальный резервуар нельзя бездумно крутить на
каждый запрос и каждый фоновой цикл.
\item \textbf{Конфликт с продуктовой адресностью.} Чем больше знания живёт в скрытой
динамике, тем тяжелее deletion, provenance и точная юридическая ревизия.
\end{enumerate}

\section{Вердикт}

Из всех ``будущих'' ветвей именно эта выглядит как самая содержательная промежуточная
ступень между сегодняшней графовой памятью и более радикальными субстратами.

\begin{principle}[Самая сильная версия ветки не радикальна, а гибридна]
Наиболее перспективная форма резервуарной ветви для Рин — это не отказ от графа, а
гибрид: явный архив, быстрый ассоциативный индекс, локально или тематически ограниченный
пластичный резервуар, replay и режимные контроллеры. Именно в этой форме она и
интерпретируема, и измерима, и инженерно правдоподобна.
\end{principle}

Если когда-нибудь эта ветвь действительно выстрелит, она даст Рин не ``ещё один
поисковый слой'', а новую форму живой памяти: менее буквальную, более ассоциативную,
но всё ещё связанную с архивом, который позволяет оставаться честной, редактируемой
и проверяемой.


\backmatter
\chapter{Литература}
\label{ch:references}

Ниже --- работы, на которые опирается монография, в согласованном формате.
Для препринтов и внутренних материалов проекта (библиографии волтов RIN и
RIN--DTI) в конце записи в скобках указан локатор во внутренней
библиографии; это указание на местоположение источника, а не замена выходных
данных.

\begin{thebibliography}{99}
\small

\bibitem{anderson2004} J.~R. Anderson, D.~Bothell, M.~D. Byrne и др.
An Integrated Theory of the Mind. \emph{Psychological Review}, 111(4), 2004.
(Базовая активация ACT-R; см. также Anderson \& Schooler, 1991.)

\bibitem{collins1975} A.~M. Collins, E.~F. Loftus.
A Spreading-Activation Theory of Semantic Processing.
\emph{Psychological Review}, 82(6), 1975.

\bibitem{hebb1949} D.~O. Hebb. \emph{The Organization of Behavior}. Wiley, 1949.

\bibitem{hopfield1982} J.~J. Hopfield.
Neural Networks and Physical Systems with Emergent Collective Computational Abilities.
\emph{PNAS}, 79(8), 1982.

\bibitem{park2023} J.~S. Park, J.~C. O'Brien, C.~J. Cai и др.
Generative Agents: Interactive Simulacra of Human Behavior. \emph{UIST}, 2023.

\bibitem{toy2024} J.~Toy, J.~MacAdam, P.~Tabor.
Metacognition is all you need? Using Introspection in Generative Agents to Improve Goal-directed Behavior.
arXiv:2401.10910, 2024.

\bibitem{kahneman2011} D.~Kahneman. \emph{Thinking, Fast and Slow}. Macmillan, 2011.
(System~1 / System~2.)

\bibitem{nisbett1977} R.~E. Nisbett, T.~D. Wilson.
Telling More Than We Can Know: Verbal Reports on Mental Processes.
\emph{Psychological Review}, 84(3), 1977. (Конфабуляция при самообъяснении.)

\bibitem{doyle1979} J.~Doyle. A Truth Maintenance System.
\emph{Artificial Intelligence}, 12(3), 1979.

\bibitem{dekleer1986} J.~de~Kleer. An Assumption-Based TMS.
\emph{Artificial Intelligence}, 28(2), 1986.

\bibitem{quine1951} W.~V.~O. Quine. Two Dogmas of Empiricism.
\emph{The Philosophical Review}, 60(1), 1951. (Тезис недоопределённости
Дюэма–Куайна.)

\bibitem{ross1975} L.~Ross, M.~R. Lepper, M.~Hubbard.
Perseverance in Self-Perception and Social Perception.
\emph{J. Personality and Social Psychology}, 32(5), 1975. (Belief perseverance.)

\bibitem{bak1987} P.~Bak, C.~Tang, K.~Wiesenfeld.
Self-Organized Criticality. \emph{Physical Review Letters}, 59(4), 1987.

\bibitem{beggs2003} J.~M. Beggs, D.~Plenz.
Neuronal Avalanches in Neocortical Circuits. \emph{J. Neuroscience}, 23(35), 2003.

\bibitem{page1999} L.~Page, S.~Brin, R.~Motwani, T.~Winograd.
The PageRank Citation Ranking: Bringing Order to the Web. Stanford InfoLab Technical Report, 1999.

\bibitem{tong2006} H.~Tong, C.~Faloutsos, J.-Y. Pan.
Fast Random Walk with Restart and Its Applications. \emph{ICDM}, 2006.

\bibitem{spielman2019} D.~A. Spielman.
\emph{Spectral and Algebraic Graph Theory} (лекционные заметки), Yale University.

\bibitem{ortega2018} A.~Ortega, P.~Frossard, J.~Kova\v{c}evi\'c и др.
Graph Signal Processing: Overview, Challenges, and Applications.
\emph{Proceedings of the IEEE}, 106(5), 2018.

\bibitem{topping2022} J.~Topping, F.~Di~Giovanni, B.~P. Chamberlain и др.
Understanding Over-Squashing and Bottlenecks on Graphs via Curvature.
\emph{ICLR}, 2022. (Волт RIN–DTI, поз.~148.)

\bibitem{black2023} M.~Black, Z.~Wan, A.~Nayyeri, Y.~Wang.
Understanding Oversquashing in GNNs through the Lens of Effective Resistance.
\emph{ICML}, 2023. (Волт RIN–DTI, поз.~183.)

\bibitem{rusch2023} T.~K. Rusch, M.~M. Bronstein, S.~Mishra.
A Survey on Oversmoothing in Graph Neural Networks. arXiv:2303.10993, 2023.
(Волт RIN–DTI, поз.~170.)

\bibitem{bodnar2021cw} C.~Bodnar, F.~Frasca, N.~Otter и др.
Weisfeiler and Lehman Go Cellular: CW Networks. \emph{NeurIPS}, 2021.
(Волт RIN–DTI, поз.~185.)

\bibitem{bodnar2021mpsn} C.~Bodnar, F.~Frasca, Y.~Wang и др.
Weisfeiler and Lehman Go Topological: Message Passing Simplicial Networks.
\emph{ICML}, 2021. (Волт RIN–DTI, поз.~186.)

\bibitem{bodnar2022sheaf} C.~Bodnar, F.~Di~Giovanni, B.~P. Chamberlain и др.
Neural Sheaf Diffusion. \emph{NeurIPS}, 2022.

\bibitem{wei2020} R.~Wang, D.~D. Nguyen, G.-W. Wei.
Persistent Spectral Graph. \emph{Int. J. Numerical Methods in Biomedical
Engineering}, 2020. (TDA/персистентный лапласиан для биомолекул; ср. TopologyNet,
Cang \& Wei, 2017. Волт RIN–DTI, поз.~100, 103, 104, 144.)

\bibitem{occ1988} A.~Ortony, G.~L. Clore, A.~Collins.
\emph{The Cognitive Structure of Emotions}. Cambridge University Press, 1988.
(OCC-модель оценивания; ср. компонентная модель K.~Scherer.)

\bibitem{friston2010} K.~Friston.
The Free-Energy Principle: A Unified Brain Theory?
\emph{Nature Reviews Neuroscience}, 11(2), 2010. (Ошибка предсказания как сигнал значимости.)

\bibitem{speer2017} R.~Speer, J.~Chin, C.~Havasi.
ConceptNet 5.5: An Open Multilingual Graph of General Knowledge.
\emph{AAAI}, 2017. (Включая эмбеддинги Numberbatch.)

\bibitem{sap2019} M.~Sap, R.~Le~Bras, E.~Allaway и др.
ATOMIC: An Atlas of Machine Commonsense for If-Then Reasoning. \emph{AAAI}, 2019.

\bibitem{hwang2021} J.~D. Hwang, C.~Bhagavatula, R.~Le~Bras и др.
(Comet-)ATOMIC 2020: On Symbolic and Neural Commonsense Knowledge Graphs.
\emph{AAAI}, 2021.

\bibitem{west2022} P.~West, C.~Bhagavatula, J.~Hessel и др.
Symbolic Knowledge Distillation: from General Language Models to Commonsense Models.
\emph{NAACL}, 2022. (ATOMIC10x.)

\bibitem{vrandecic2014} D.~Vrande\v{c}i\'c, M.~Kr\"otzsch.
Wikidata: A Free Collaborative Knowledgebase.
\emph{Communications of the ACM}, 57(10), 2014.

\bibitem{xu2019} K.~Xu, W.~Hu, J.~Leskovec, S.~Jegelka.
How Powerful are Graph Neural Networks? \emph{ICLR}, 2019. (GIN; волт RIN–DTI,
поз.~68.)

\bibitem{corso2020} G.~Corso, L.~Cavalleri, D.~Beaini и др.
Principal Neighbourhood Aggregation for Graph Nets. \emph{NeurIPS}, 2020.
(Волт RIN–DTI, поз.~178.)

\bibitem{hu2021lora} E.~J. Hu, Y.~Shen, P.~Wallis и др.
LoRA: Low-Rank Adaptation of Large Language Models. arXiv:2106.09685, 2021.

\bibitem{rafailov2023} R.~Rafailov, A.~Sharma, E.~Mitchell и др.
Direct Preference Optimization. \emph{NeurIPS}, 2023.

\bibitem{zhou2006hg} D.~Zhou, J.~Huang, B.~Sch\"olkopf.
Learning with Hypergraphs: Clustering, Classification, and Embedding.
\emph{NeurIPS}, 2006. (Гиперграфовый лапласиан.)

\bibitem{lim2023hodge} L.-H. Lim. Hodge Laplacians on Graphs.
\emph{SIAM Review}, 62(3), 2020. (Ходж-лапласианы и высшие размерности.)

\bibitem{helamem2026} J.~Zhu, J.~Li, C.~Zhang и др.
HeLa-Mem: Hebbian Learning and Associative Memory for LLM Agents. Препринт, 2026.
(Внутренняя библиография волта RIN, поз.~53.)

\bibitem{rinvault} Библиография памяти LLM-агентов из волта RIN: APEX-MEM, TRACE,
SYNAPSE, GAAMA, MRMS, Zep/Graphiti, Persistent Sycophancy, LifeSide, LoCoMo и др.
(Внутренние материалы проекта; полные выходные данные --- в
\texttt{99 Sources/1.~Bibliography (RIN)}.)

\bibitem{rinproject} Материалы проекта RIN: модульные системные промпты
(\texttt{prompts/global\_core/*}), persona-тесты и логи Stage~0
(\texttt{logs/rin\_stage0.jsonl}), Ринциклопедии (\texttt{docs/rincyclopedia\_*}).

\end{thebibliography}


\end{document}

